% Options for packages loaded elsewhere
\PassOptionsToPackage{unicode}{hyperref}
\PassOptionsToPackage{hyphens}{url}
%
\documentclass[
  letterpaper,
]{scrbook}
\usepackage{amsmath,amssymb}
\usepackage{lmodern}
\usepackage{iftex}
\ifPDFTeX
  \usepackage[T1]{fontenc}
  \usepackage[utf8]{inputenc}
  \usepackage{textcomp} % provide euro and other symbols
\else % if luatex or xetex
  \usepackage{unicode-math}
  \defaultfontfeatures{Scale=MatchLowercase}
  \defaultfontfeatures[\rmfamily]{Ligatures=TeX,Scale=1}
\fi
% Use upquote if available, for straight quotes in verbatim environments
\IfFileExists{upquote.sty}{\usepackage{upquote}}{}
\IfFileExists{microtype.sty}{% use microtype if available
  \usepackage[]{microtype}
  \UseMicrotypeSet[protrusion]{basicmath} % disable protrusion for tt fonts
}{}
\makeatletter
\@ifundefined{KOMAClassName}{% if non-KOMA class
  \IfFileExists{parskip.sty}{%
    \usepackage{parskip}
  }{% else
    \setlength{\parindent}{0pt}
    \setlength{\parskip}{6pt plus 2pt minus 1pt}}
}{% if KOMA class
  \KOMAoptions{parskip=half}}
\makeatother
\usepackage{xcolor}
\usepackage{longtable,booktabs,array}
\usepackage{calc} % for calculating minipage widths
% Correct order of tables after \paragraph or \subparagraph
\usepackage{etoolbox}
\makeatletter
\patchcmd\longtable{\par}{\if@noskipsec\mbox{}\fi\par}{}{}
\makeatother
% Allow footnotes in longtable head/foot
\IfFileExists{footnotehyper.sty}{\usepackage{footnotehyper}}{\usepackage{footnote}}
\makesavenoteenv{longtable}
\usepackage{graphicx}
\makeatletter
\def\maxwidth{\ifdim\Gin@nat@width>\linewidth\linewidth\else\Gin@nat@width\fi}
\def\maxheight{\ifdim\Gin@nat@height>\textheight\textheight\else\Gin@nat@height\fi}
\makeatother
% Scale images if necessary, so that they will not overflow the page
% margins by default, and it is still possible to overwrite the defaults
% using explicit options in \includegraphics[width, height, ...]{}
\setkeys{Gin}{width=\maxwidth,height=\maxheight,keepaspectratio}
% Set default figure placement to htbp
\makeatletter
\def\fps@figure{htbp}
\makeatother
\setlength{\emergencystretch}{3em} % prevent overfull lines
\providecommand{\tightlist}{%
  \setlength{\itemsep}{0pt}\setlength{\parskip}{0pt}}
\setcounter{secnumdepth}{5}
% Make \paragraph and \subparagraph free-standing
\ifx\paragraph\undefined\else
  \let\oldparagraph\paragraph
  \renewcommand{\paragraph}[1]{\oldparagraph{#1}\mbox{}}
\fi
\ifx\subparagraph\undefined\else
  \let\oldsubparagraph\subparagraph
  \renewcommand{\subparagraph}[1]{\oldsubparagraph{#1}\mbox{}}
\fi
\RedeclareSectionCommand[beforeskip=-2.5ex plus -1ex minus -.2ex,afterskip=0.3ex plus -.2ex]{section}
\usepackage[top=0.5in, left=0.5in, right=0.5in, bottom=1in]{geometry}
\bibliographystyle{catalog}
\usepackage{makeidx}
\makeindex
\renewcommand\toprule[2]\relax
\renewcommand\bottomrule[2]\relax
\usepackage{fancyhdr}
\lfoot{\textit{Coe College (2022-2023)}}
\renewcommand{\part}[1]{\addcontentsline{toc}{part}{#1}}
\makeatletter
\@ifpackageloaded{bookmark}{}{\usepackage{bookmark}}
\makeatother
\makeatletter
\@ifpackageloaded{caption}{}{\usepackage{caption}}
\AtBeginDocument{%
\ifdefined\contentsname
  \renewcommand*\contentsname{Table of contents}
\else
  \newcommand\contentsname{Table of contents}
\fi
\ifdefined\listfigurename
  \renewcommand*\listfigurename{List of Figures}
\else
  \newcommand\listfigurename{List of Figures}
\fi
\ifdefined\listtablename
  \renewcommand*\listtablename{List of Tables}
\else
  \newcommand\listtablename{List of Tables}
\fi
\ifdefined\figurename
  \renewcommand*\figurename{Figure}
\else
  \newcommand\figurename{Figure}
\fi
\ifdefined\tablename
  \renewcommand*\tablename{Table}
\else
  \newcommand\tablename{Table}
\fi
}
\@ifpackageloaded{float}{}{\usepackage{float}}
\floatstyle{ruled}
\@ifundefined{c@chapter}{\newfloat{codelisting}{h}{lop}}{\newfloat{codelisting}{h}{lop}[chapter]}
\floatname{codelisting}{Listing}
\newcommand*\listoflistings{\listof{codelisting}{List of Listings}}
\makeatother
\makeatletter
\makeatother
\makeatletter
\@ifpackageloaded{caption}{}{\usepackage{caption}}
\@ifpackageloaded{subcaption}{}{\usepackage{subcaption}}
\makeatother
\ifLuaTeX
  \usepackage{selnolig}  % disable illegal ligatures
\fi
\IfFileExists{bookmark.sty}{\usepackage{bookmark}}{\usepackage{hyperref}}
\IfFileExists{xurl.sty}{\usepackage{xurl}}{} % add URL line breaks if available
\urlstyle{same} % disable monospaced font for URLs
\hypersetup{
  pdftitle={Course Catalog 2022-2023},
  hidelinks,
  pdfcreator={LaTeX via pandoc}}

\title{Course Catalog 2022-2023}
\author{}
\date{2022-08-01}

\begin{document}
\frontmatter
\maketitle

\renewcommand*\contentsname{Table of contents}
{
\setcounter{tocdepth}{1}
\tableofcontents
}
\mainmatter
\bookmarksetup{startatroot}

\chapter*{PREFACE}\label{preface}

\markboth{PREFACE}{PREFACE}

\section*{Non-Discrimination}\label{non-discrimination}

\markright{Non-Discrimination}

Coe College does not discriminate on the basis of race, color,
ethnicity, age, religion, national origin, sexual orientation, gender
identity, sex, marital status, disability, or status as a U.S. Veteran.
All students have equal access to the facilities, financial aid, and
programs of the College.

\section*{Higher Education Opportunity Act
(HEOA)}\label{higher-education-opportunity-act-heoa}

\markright{Higher Education Opportunity Act (HEOA)}

The College complies with Readmission Requirements for Service Members
as outlined in the Higher Education Opportunity Act section 487. This
applies to active duty in the Armed Forces, whether voluntary or
involuntary, including service as a member of the National Guard or
Reserve, for a period of more than 30 days under a call or order to
active duty.

The HEOA provides that a prompt readmission of a previously enrolled or
admitted student may not be denied to a service member of the uniformed
services for reasons relating to that service.~ In addition, a student
who is readmitted under this section must be readmitted with the same
academic status as the student had when they attended the college.

\section*{Equal Opportunity in
Employment}\label{equal-opportunity-in-employment}

\markright{Equal Opportunity in Employment}

Coe College is an equal opportunity employer in the recruitment and
hiring of faculty and staff.

\section*{Family Educational Rights and Privacy Act
(FERPA)}\label{family-educational-rights-and-privacy-act-ferpa}

\markright{Family Educational Rights and Privacy Act (FERPA)}

The provisions of the Family Educational Rights and Privacy Act (FERPA)
prohibit the College from releasing grades or other information about
academic standing to parents unless the student has released such
information in writing. Further information concerning Coe College
procedures in compliance with FERPA is available in the Office of the
Registrar and included on p.~68 of this catalog.

\section*{Solomon Amendment of 1997}\label{solomon-amendment-of-1997}

\markright{Solomon Amendment of 1997}

Pursuant to the regulations of the Solomon Amendment of 1997, Coe
College is required to make student recruiting information available to
military recruiters who request it.

\section*{Reservation of the Right to
Modify}\label{reservation-of-the-right-to-modify}

\markright{Reservation of the Right to Modify}

The provisions of this catalog are to be considered directive in
character and not as an irrevocable contract between the student and the
College. The College reserves the right to make changes that seem
necessary or desirable, including course and program cancellations.
Responsibility for understanding and meeting graduation requirements as
stated in the Coe College Catalog rests entirely with the student.
Faculty advisors and the Registrar will assist in every way possible.

\part{INTRODUCTORY RESOURCES}

\chapter{MISSION STATEMENT OF THE
COLLEGE}\label{sec-mission-statement-of-college}

\begin{center}
\includegraphics{catalog_sections/graphics/coe_bell.png}
\end{center}

\index{Mission Statement}

Coe College is a national, residential liberal arts college offering a
broad array of programs in the arts, sciences and professions. Our
mission is to prepare students for meaningful lives and fulfilling
careers in a diverse, interconnected world. Coe's success will be judged
by the success of our graduates.

\emph{Coe College admits students without regard to sex, race, creed,
color, handicap, sexual orientation, national, or ethnic origin. All
students have equal access to the facilities, financial aid, and
programs of the College.}

\chapter{FAST FACTS}\label{sec-fast-facts}

\textbf{COE COLLEGE} is a private, four-year co-educational liberal arts
college that was founded in 1851 and is historically affiliated with the
Presbyterian Church (U.S.A.), yet is ecumenical in practice and outlook.

\subsection{LOCATION}\label{location}

Coe is located just 225 miles west of Chicago in Iowa's second largest
city, Cedar Rapids (metropolitan population 175,000). The city is easily
accessible by the interstate highway system, bus services and several
airlines.

\subsection{CAMPUS}\label{campus}

Situated on 65 acres in the center of the metropolitan area, the campus
is urban but enclosed, with 30 buildings occupying an attractive
landscape.

\subsection{ENROLLMENT}\label{enrollment}

The student body of approximately 1,400 students represents most states
and around 15 foreign countries. All students are required to live on
campus unless they are residents of Cedar Rapids.

\subsection{LIBRARY}\label{library}

Stewart Memorial Library is in the very center of the main campus. It
contains over 500,000 volumes and 16,000 pieces of media, and subscribes
to 3,500 print and online periodical subscriptions providing access to
over 100 databases. Coe's library offers students a variety of research
assistance, study areas, a small theater, preview room, a media editing
room, and a 3D printer and laser cutter. The Library houses the Learning
Commons and the College's permanent art collection.

\subsection{ACCREDITATION}\label{accreditation}

Coe College is accredited by the Higher Learning Commission
(hlcomission.org), an institutional accreditation agency recognized by
the U.S. Department of Education. Coe College's Bachelor of Music
program is accredited by the National Association of Schools of Music;
its education program is accredited by the Iowa Department of Education,
and its Bachelor of Science in Nursing is accredited by the Commission
on Collegiate Nursing Education
(http://aacn.nche.edu/ccne-accreditation).~ Coe College's chemistry
program is certified by the American Chemical Society.~ Copies of
accrediting and approval statements are available online at
https://www.coe.edu/why-coe/consumer-information.

\subsection{MEMBERSHIPS}\label{memberships}

Coe is a charter member of the Associated Colleges of the Midwest, whose
other members are: Beloit, Carleton, Colorado, Cornell, Grinnell, Knox,
Lake Forest, Lawrence, Luther, Macalester, Monmouth, Ripon, and
St.~Olaf. Other memberships include: Phi Beta Kappa, Phi Kappa Phi, the
National Association of Independent Colleges and Universities, the
American Association of Colleges and Universities, Council on
Undergraduate Research, the Association of Presbyterian Colleges and
Universities, the American Rivers Conference, and the National
Collegiate Athletic Association.

\subsection{ACADEMIC PROGRAM}\label{academic-program}

Academic areas of study are offered in managerial or public accounting,
African American studies, art, art history, Asian studies, biology,
business administration, chemistry, communication studies, computer
science, creative writing, data science, economics, elementary
education, English, film studies, French \& francophone studies,
interdisciplinary French \& francophone studies, general science,
history, international business, international economics, international
studies, kinesiology, literature, mathematics, music (B.A. or B.M.),
nursing (B.S.N.), philosophy, physics, political science, psychology,
religion, social \& criminal justice, sociology, Spanish, Spanish
studies, theatre arts, and writing. Interdisciplinary and/or collateral
majors are also available in biochemistry, environmental science,
environmental studies, molecular biology, neuroscience, organizational
science, and public relations. Coe also offers certificate programs in
primary and secondary education.

\subsection{EXTRA-CURRICULAR
ACTIVITIES}\label{extra-curricular-activities}

Students can participate in 11 men's and 11 women's NCAA Division III
varsity athletic teams, 5 coed varsity athletic teams, 8 club sports, 10
nationally affiliated men's and women's social fraternities, and 65
student organizations including but not limited to student government,
newspaper, intramural sports, departmental clubs, and residence hall
activities. Various vocal and instrumental ensembles are available for
course credit.

\subsection{FINANCIAL ASSISTANCE}\label{financial-assistance}

Coe College is committed to assisting those families in need of
financial assistance. The average aid package for incoming students
enrolled during the 2021--2022 academic year totaled more than \$48,000.
The total cost of full-time tuition, room, board, and activity fee for
the 2021--2022 academic year is \$59,324.

\subsection{HISTORY OF COE COLLEGE}\label{sec-history-of-coe-college}

Coe College claims the shortest name of any American institution of
higher education, but the school has actually carried five titles
through its history. When the Rev.~Williston Jones founded the college
in 1851, he called it The School for the Prophets. Cedar Rapids' first
resident minister opened the parlor of his home to a group of young men
with the goal of educating them for the ministry to serve churches in
the Midwest. Two years later, while Jones was canvassing churches in the
East for money to send three of his students to Eastern seminaries, a
Catskills farmer named Daniel Coe stepped forward with a pledge of
\$1,500 and urged Jones to start his own college in the frontier town of
Cedar Rapids. Legend has it that the \$1,500 raised by Coe was brought
west from New York, sewn into the petticoat of a lady visitor traveling
by stagecoach to Iowa. Coe gave this generous gift with the stipulation
that the proposed institute should be ``made available for the education
of females as well as males.'' Accordingly, Coe was coeducational from
its founding.

With Jones' blessing, the Cedar Rapids Collegiate Institute was
incorporated in 1853 by a group of Cedar Rapids leaders chaired by Judge
George Greene. They used Daniel Coe's money to purchase two downtown
lots for the school and 80 acres of farmland on what was then the edge
of town. The farm evolved into today's campus. In 1868, in a failed
attempt to secure the Lewis Parsons estate, the trustees renamed the
school Parsons Seminary. After a period of severe financial
difficulties, the institution was reestablished in honor of its original
benefactor as the Coe Collegiate Institute in 1875.

T.M Sinclair, founder of the Sinclair Meat Packing Company, played the
key financial role in the final step toward the firm establishment of
Coe College. Sinclair liquidated all the debt from Parsons Seminary and
the Cedar Rapids Collegiate Institute. The Sinclair gift made it
practical for the property of the Coe Collegiate Institute---including
the original land paid for by Daniel Coe---to be transferred to Coe
College with the Iowa Presbyterian Synod to assume major responsibility
for the institution. Coe College has operated continuously since its
incorporation under that name on Feb.~2, 1881. From the first, the
College was committed to intellectual excellence. It has continued in
this tradition ever since.

The compact campus on the east edge of Cedar Rapids grew with many
building projects in its early years, including Old Main (1868),
Williston Hall (1881), Marshall Hall (1900), the first gymnasium (1904),
and the first T.M. Sinclair Memorial Chapel (1911). In 1907, Coe earned
accreditation from the North Central Association of Colleges and
Universities. Over the decades, Coe's reputation as a superior liberal
arts college has continued to grow. One recognition of this came in
1949, when Coe was granted a Phi Beta Kappa chapter, a distinction
reserved for about 10\% of American colleges and universities.

Central to the educational philosophy of Coe College is the belief that
a liberal arts education is the best preparation for life. Students have
the opportunity to experience a variety of subjects outside their
respective programs of study. Coe offers more than 40 areas of study
that cover a range of fields. The College awards the following
undergraduate degrees: Bachelor of Arts (B.A.), Bachelor of Music
(B.M.), and Bachelor of Science in Nursing (B.S.N.).

There are a number of factors that contribute to Coe College's strong
academic quality. The key to Coe's tradition of excellence in academic
quality relates directly to small class sizes and the interest shown by
professors to make learning a personalized experience. At Coe, the
average class size is 16, and the student-faculty ratio is 11:1. Classes
are taught by our involved and committed faculty, 91\% of whom hold the
highest degree in their field. This means classes are taught by
experienced professionals who have in-depth knowledge of their subjects.
To provide students with a well-rounded experience and solid preparation
for the future, Coe offers, along with quality instruction from superb
faculty, an abundance of out-of-class opportunities which include
student-faculty collaborative research, honors projects, service
learning, and internships. Within one year of graduation, according to
the annual survey results from the Center for Creativity and Careers,
98\% of Coe graduates are engaged in post-graduate activity such as
employment, graduate school, military, or travel/adventure.

With the addition of the east campus, Coe has nearly doubled in size
since 1989. New facilities on the east side of College Drive include
Athletic Recreation Center (2017), Clark Racquet Center and athletic
fields (1989), Clark Alumni House (1993), Nassif House (1999), and four
student apartment buildings (Morris House and Schlarbaum House in 2000,
Brandt House and Spivey House in 2002). McCabe Hall (2005), named in
honor of former Coe President Joseph E. McCabe, houses the offices of
the president, provost, advancement and alumni relations, and The Center
for Creativity, Careers and Community (C3) making way for the remodeling
of Coe's oldest building, Stuart Hall, and the first significant
addition of classroom space since Peterson Hall was built in the 1960s.
In 2012 and 2013, Peterson Hall of Science was completely renovated to
support Coe's science programs. To further enhance the campus
environment, Coe completed the largest capital project in its history in
2017. Make Your Move -- the Campaign for Eby and Hickok -- included \$24
million in essential enhancements, including an Athletic and Recreation
Complex project as well as the renovation and expansion of Hickok Hall,
one of the College's main academic buildings. The result is vastly
improved academic, recreational, wellness and competition facilities to
benefit future generations of students.

\part{ACADEMIC RESOURCES}

\chapter{2022--2023 ACADEMIC CALENDAR}\label{sec-academic-calenar}

\subsection{FALL TERM 2022}\label{fall-term-2022}

\begin{longtable}[]{@{}
  >{\raggedright\arraybackslash}p{(\columnwidth - 2\tabcolsep) * \real{0.3333}}
  >{\raggedleft\arraybackslash}p{(\columnwidth - 2\tabcolsep) * \real{0.6667}}@{}}
\toprule\noalign{}
\endhead
\bottomrule\noalign{}
\endlastfoot
Tues, August 23 & Open Registration \\
Wed, August 24 & Classes Begin \\
Tues, August 30 & Last Day to Add or Drop a Full-Term, or 1st Half-Term
Course without a W \\
Fri, September 2 & Census Date, Attendance Due on My.Coe \\
Mon, September 5 & No Classes (Holiday) Labor Day \\
Fri, September 9 & Date of Record \\
Thurs, September 22 & Last Day to Withdraw \&/ or Change Method of
Grading for First-Half Courses \\
Mon--Tues, October 10--11 & Fall Term Break \\
Thurs--Sat, October 13--15 & Homecoming \\
Mon, October 17 & Mid-Term Progress Report Due on My.Coe (11:59 PM) \\
Mon, October 17 & Begin Second-Half Term Courses \\
Fri, October 21 & May Term deposit and registration due \\
Fri, October 21 & Last Day to Add or Drop a Second-Half Term Course
without a W \\
Mon, October 24-Fri, November 11 & Advising Season \\
Fri, October 28 & Last Day to Withdraw \&/or Change Method of Grading
for Full-Term Courses \\
Mon--Fri, November 14--18 & Registration for Spring Term 2023 \\
Tues, November 15 & Last Day to Withdraw \&/or Change Method of Grading
for Second-Half Courses \\
Mon--Fri, November 21--25 & No Classes (Holiday) Thanksgiving Recess \\
Thurs, December 1 & Open Online Registration until Last Day of Finals \\
Fri, December 2 & Last Day of Fall Term Classes \\
Sat, December 3 & Reading Day \\
Mon--Thurs, December 5--8 & Final Exams \\
Wed, December 14 & Final Grades and Attendance are Due from the Faculty
on My.Coe (11:59 PM) \\
\end{longtable}

\subsection{SPRING TERM 2023}\label{spring-term-2023}

\begin{longtable}[]{@{}
  >{\raggedright\arraybackslash}p{(\columnwidth - 2\tabcolsep) * \real{0.3333}}
  >{\raggedleft\arraybackslash}p{(\columnwidth - 2\tabcolsep) * \real{0.6667}}@{}}
\toprule\noalign{}
\endhead
\bottomrule\noalign{}
\endlastfoot
Wed, January 11 & Classes Begin \\
Mon, January 16 & No Day Classes (Holiday) Martin Luther King Jr.~Day \\
Wed, January 18 & Last Day to Add or Drop a Full-Term or 1st Half-Term
Course without a W \\
Fri, January 20 & Census Date, Attendance Due on My.Coe \\
Thurs, January 26 & Date of Record Mon, February 6 \\
Mon, February 6 & Last Day to Withdraw \&/or Change Method of Grading
for First-Half Courses \\
Mon, March 6 & Begin Second-Half Term Courses \\
Tues, March 7 & Mid-Term Progress Report Due Online (11:59 PM) \\
Fri, March 10 & Last Day to Add or Drop a Second-Half Term Course
without a W \\
Mon--Fri, March 13--17 & Spring Term Recess (Starting after Class on
Friday, March 10) \\
Mon--Fri, March 20--April 7 & Advising Season Thurs, March 30 \\
Thurs, March 30 & Last Day to Withdraw \&/or Change Method of Grading
for Full-Term Courses \\
Mon--Fri, April 10--14 & Registration for Fall Term 2023 \\
Mon, April 10 & Summer Registration Opens \\
Tues, April 11 & Last Day to Withdraw \&/or Change Method of Grading for
Second-Half Courses \\
Wed, April 12 & Student Research Symposium ** No Day Classes ** Evening
Classes Will Meet \\
Fri, April 28 & Last Day of Spring Term Classes \\
Sat, April 29 & Reading Day \\
Mon--Thurs, May 1--4 & Final Exams \\
Sat, May 6 & Honors Convocation / Baccalaureate \\
Sun, May 7 & Commencement \\
Tues, May 9 & Final Grades and Attendance are Due for Non-Graduating
Students on My.Coe (11:59 PM) \\
\end{longtable}

\subsection{MAY TERM 2023}\label{may-term-2023}

\begin{longtable}[]{@{}
  >{\raggedright\arraybackslash}p{(\columnwidth - 2\tabcolsep) * \real{0.3333}}
  >{\raggedleft\arraybackslash}p{(\columnwidth - 2\tabcolsep) * \real{0.6667}}@{}}
\toprule\noalign{}
\endhead
\bottomrule\noalign{}
\endlastfoot
Wed, May 10 & Classes Begin \\
Fri, May 12 & Last Day to Add or Drop a Course Without a W \\
Mon, May 15 & Attendance Due Online; Last Day to Change Method of
Grading for May Term \\
Wed, May 17 & Last Day to Withdraw from May Term Courses \\
Fri--Sat, May 19--20 & Meeting of the Board of Trustees \\
Mon, May 29 & No Classes (Holiday) Memorial Day \\
Fri, June 2 & Last Day of May Term Classes \\
Sat, June 3 & Residence Halls Close \\
Sun, June 4 & Final Grades and Attendance for May Term Due on My.Coe
(11:59 PM) \\
\end{longtable}

\subsection{SUMMER TERM 2023*}\label{summer-term-2023}

\begin{longtable}[]{@{}
  >{\raggedright\arraybackslash}p{(\columnwidth - 2\tabcolsep) * \real{0.3333}}
  >{\raggedleft\arraybackslash}p{(\columnwidth - 2\tabcolsep) * \real{0.6667}}@{}}
\toprule\noalign{}
\endhead
\bottomrule\noalign{}
\endlastfoot
Mon, June 5 & Classes Begin \\
Fri, June 9 & Last Day to Add or Drop a Full-Term Course Without a W \\
Mon, June 19 & No Classes (Holiday) Juneteenth \\
Fri, June 23 & Last Day to Change Method of Grading \&/or Withdraw from
Block A Courses \\
Tues, July 4 & No Classes (Holiday) Independence Day \\
Fri, July 7 & Last Day of Block A Courses \\
Mon, July 10 & Classes Begin: Block B Courses \\
Tues, July 11 & Mid-Term Progress Report Due Online (11:59 PM); Block A
Final Grades Due \\
Fri, July 14 & Last Day to Add or Drop a Course Without a W for Block B
Courses \\
Mon, July 17 & Last Day to Change Method of Grading \&/or Withdraw from
Full Term Courses \\
Tues, July 25 & Last Day to Change Method of Grading \&/or Withdraw from
Block B Courses \\
Fri, August 11 & Last Day of Term Classes: Block B and Full-Term
Courses \\
Tues, August 15 & Final Grades and Attendance Due (Block B and
Full-Term) on My.Coe (11:59 PM) \\
\end{longtable}

*Courses taught at the Wilderness Field Station are subject to the broad
dates of the summer, but will provide a specific add/drop and withdraw
calendar to students at the time of application.

\chapter{THE ACADEMIC CALENDAR}\label{the-academic-calendar}

The academic year consists of three terms (see Academic Calendar,
p.~9-11). Students normally take four course credits in the Fall Term
and four course credits in the Spring Term. Thus, eight course credits
are completed in an academic year. During optional May Term, students
may enroll for up to one course credit in one of the limited selection
of courses. Summer term is limited in scope and is not considered a
regular term. (The maximum course load is described in Course Load,
p.~41).

\chapter{MAY TERM (OPTIONAL)}\label{may-term-optional}

Students may enroll for up to one course credit during May Term in one
of the limited selection of courses. All May Term courses require
consent of instructor prior to registration. May Term courses are
designed to meet at least two of the following shared learning outcomes:

\begin{itemize}
\tightlist
\item
  Evaluate and engage with complex interdependent systems and
  demonstrate understanding across diverse contexts.
\item
  Critically describe and break apart issues or problems through
  systematic analysis and illustrate logic for conclusions.
\item
  Engage with experiential learning practices such as learning by doing,
  while utilizing abilities to think critically, problem solve and make
  connections between knowledge gained in the classroom and experience
  beyond.
\end{itemize}

May Term courses are expected to have the same amount of contact time
and academic rigor per course credit as courses which meet over a Fall
or Spring Term.~ For every 1 course credit of May term students are
expected to complete 140 hours of work. Such contact time includes class
meetings, lectures by the instructor, supervised course related
activities and independent out of class activities.~ Off-campus May Term
courses at an off-campus location provide certain educational benefits
through site visits, guest lectures, etc., that also contribute to the
contact time for the course.

\chapter{COE PLAN}\label{sec-coe-plan}

Coe College's requirements for graduation, commonly known as the Coe
Plan were developed with the following outcomes in mind:

\begin{itemize}
\tightlist
\item
  Creation of a bridge from high school to Coe College that helps
  students understand the importance of a liberal arts education, the
  ways to develop the skills needed by any learner, and the
  opportunities they will have by going to Coe College.
\item
  Development of required curriculum that exposes the students to ways
  of learning in various contexts, big ideas in a myriad of disciplines,
  ways of being and understanding of cultures around the world, and
  processes to develop the skills needed by any learner.
\item
  Creation of a bridge from Coe College to life after Coe.
\end{itemize}

These outcomes are met through the College's First-Year Experience,
General Education program, Writing Emphasis courses, and the College's
Practicum experiences and areas of study, described in this section of
the Catalog.

\chapter{GRADUATION REQUIREMENTS}\label{sec-graduation-requirements}

All students who graduate from Coe College must complete at least one
major and earn at least 32 course credits (cc) with grades leading to a
cumulative grade point average (GPA) of 2.0 or higher. (The course is
the unit of academic credit.) Courses are one credit unless otherwise
indicated. Students are expected to complete 180 hours of work to earn
one course credit, although class times vary from course to course.
Other institutions may convert Coe credit to their system by considering
one course credit to be 6 quarter hours, or 4 semester hours.) No more
than a total of two course credits from courses which are less than 0.5
credit can be used to meet the 32-credit graduation requirement. No more
than eight course credits earned of Advanced Placement or International
Baccalaureate credit can be used to satisfy this requirement.

Students must meet one of the following requirements:

\begin{itemize}
\tightlist
\item
  Complete at least the final academic year of required courses
  registered through Coe.
\item
  Earn a total of 16 course credits or the equivalent at Coe. The last
  eight course credits needed for graduation must include at least four
  earned at Coe. Approved off-campus study programs and internships can
  be used to fulfill this requirement.
\end{itemize}

A student may be simultaneously awarded two degrees (B.A., B.M., B.S.N.)
after satisfactorily completing 40 course credits and the requirements
for both degrees. However, a simultaneous Bachelor of Music plus a
Bachelor of Arts with a music major is not permitted.

In addition, students must fulfill the requirements of the First-Year
Experience, General Education, Writing Emphasis, and Practicum.

To participate in Commencement exercises, students must submit a
completed Intent to Graduate form to the Office of the Registrar,
preferably three terms prior to Commencement.

\section{Second Baccalaureate Degree}\label{second-baccalaureate-degree}

A student who holds a baccalaureate degree from another institution may
earn a second baccalaureate degree at Coe, if the following criteria are
met:

\begin{itemize}
\tightlist
\item
  The first degree must be from a regionally accredited institution as
  recognized by the US Department of Education, or another appropriate
  accrediting body.
\item
  The first degree must be completed (not in progress) before beginning
  the second degree at Coe.
\end{itemize}

Students accepted at Coe to pursue a second degree are granted a maximum
of 24 course credits in transfer credit towards the 32 course credits
required for graduation. To graduate, at least eight course credits must
be earned at Coe College and all requirements for the major area of
study must be met with at least 40\% of the major course credits taken
at Coe. Students must earn a cumulative GPA of at least 2.0 as well as
meet any GPA or grade requirements in their area of study.

Second baccalaureate students are exempt from the following
requirements: first-year experience, general education, writing
emphasis, and practicum. They are not eligible to graduate with Latin
Honors or for induction in Phi Beta Kappa or Phi Kappa Phi.

\section{Transfer Student
Information}\label{transfer-student-information}

To honor its mission and to preserve its academic integrity as a liberal
arts institution, the College accepts a course in transfer if that
course meets the spirit of the College's mission and is from a
regionally accredited institution. This section includes information, in
addition to that included in the section, Graduation Requirements (see
p.~19), germane to students who are transferring to Coe College from
another college or university.

Courses transferred to Coe can, as approved by the Registrar, fulfill
some graduation requirements. From institutions on a semester hour
system (at Coe, 1 course credit = 4 semester hours), only courses with
three or more semester hours can be used to fulfill any major or general
education requirements. From institutions on other than a semester hour
system, only courses equivalent to at least 0.75 course credits can be
used to fulfill any major or general education requirement. In some
cases, in consultation with the Registrar, multiple courses within the
same field may be used to fulfill one requirement.

Transfer credits earned after high-school graduation and before Coe
matriculation count towards the eight term, full-time residence
requirement (see p.~211). Full-time enrollment may include participation
in Coe College exchange programs, ACM off-campus study programs, and
other approved off-campus study programs.

All students must complete at Coe at least 40\% of the total course
credits required for each declared major or minor or three course
credits, whichever is greater. In addition to completing at least one
major area of study, transfer students must abide with the following to
complete the requirements for graduation:

\begin{itemize}
\tightlist
\item
  First-Year Experience. Transfer students are not required to fulfill
  the requirements of the First-Year Experience, if they have completed
  at least one full-time college term since graduation from high school.
\item
  Writing Emphasis. (See Writing Emphasis Courses, p.~29).
\item
  General Education. Requirements include Liberal Arts selections in the
  four divisional areas (Natural and Mathematical Sciences, Social
  Sciences, Humanities, Fine Arts) and Diverse Cultural Perspectives
  courses. Any courses accepted in transfer for at least 0.75 course
  credit that fit the criteria of the Liberal Arts and/or Diverse
  Cultural Perspectives core groups can be applied towards the general
  education requirements as determined by the Registrar. Advanced
  Placement and International Baccalaureate courses may not be used to
  meet any part of the General Education requirements.
\item
  Academic Practicum. Transfer students are required to fulfill this
  requirement.
\end{itemize}

\chapter{AREAS OF STUDY}\label{sec-areas-of-study}

The three undergraduate degrees have areas of study associated with
them. The Bachelor of Science in Nursing's area of study is nursing; the
Bachelor of Music's areas of study are performance, composition, and
education. The Bachelor of Arts' areas of study, commonly referred to as
majors, are listed below.

Students should declare an area of study by the end of their sophomore
year. All students must earn at least a 2.00 GPA in courses required to
complete their areas of study, as well as meet specific requirements set
forth for the areaof study.

\begin{itemize}
\tightlist
\item
  Accounting, Managerial
\item
  Accounting, Public
\item
  African American Studies
\item
  Art
\item
  Art History
\item
  Asian Studies
\item
  Biology
\item
  Business Administration
\item
  Chemistry
\item
  Communication Studies
\item
  Computer Science
\item
  Creative Writing
\item
  Data Science
\item
  Economics
\item
  Elementary Education
\item
  English
\item
  Film Studies
\item
  French \& Francophone Studies
\item
  General Science
\item
  History
\item
  Interdisciplinary
\end{itemize}

\begin{itemize}
\tightlist
\item
  French \& Francophone Studies
\item
  Interdisciplinary Studies*
\item
  International Business
\item
  International Economics
\item
  International Studies
\item
  Kinesiology
\item
  Literature
\item
  Mathematics
\item
  Music
\item
  Philosophy
\item
  Physics
\item
  Political Science
\item
  Psychology
\item
  Religion
\item
  Social \& Criminal Justice
\item
  Sociology
\item
  Spanish
\item
  Spanish Studies
\item
  Theatre Arts
\item
  Writing (Rhetoric)
\end{itemize}

* A coherent interdisciplinary sequence of courses devised by the
student, in consultation with faculty, suited to the student's
individual goals and approved by the Academic Policies Committee (see
p.~135). In addition to the areas of study/majors listed above, the
following \textbf{COLLATERAL MAJORS} are offered, which require a
student to satisfy the requirements of a major from the list above in
addition to the selected collateral major.

\begin{itemize}
\tightlist
\item
  Biochemistry
\item
  Molecular Biology
\item
  Public Relations
\item
  Environmental Science
\item
  Neuroscience
\item
  Environmental Studies
\item
  Organizational Science
\end{itemize}

\subsection{AREAS OF STUDY (MINOR) FOR
B.A.}\label{areas-of-study-minor-for-b.a.}

\begin{itemize}
\tightlist
\item
  African American Studies
\item
  Anthropology
\item
  Art
\item
  Art History
\item
  Asian Studies
\item
  Chemistry
\item
  Classical Studies
\item
  Communication Studies
\item
  Computer Science
\item
  Creative Writing
\item
  Data Science
\item
  Economics
\item
  English
\item
  Film Studies
\item
  French \& Francophone Studies
\item
  Gender and Sexuality Studies
\item
  Health \& Society Studies
\item
  History
\item
  Interdisciplinary French \&
\item
  Francophone Studies
\item
  International Economics
\item
  Mathematics
\end{itemize}

\subsection{AREAS OF STUDY (MAJORS) FOR
B.M.}\label{areas-of-study-majors-for-b.m.}

\begin{itemize}
\tightlist
\item
  Keyboard or Instrumental Performance
\item
  Vocal Performance
\item
  Composition
\item
  Instrumental Music Education
\item
  Vocal Music Education
\end{itemize}

\subsection{AREA OF STUDY (MAJOR) FOR
B.S.N.}\label{area-of-study-major-for-b.s.n.}

\begin{itemize}
\tightlist
\item
  Nursing
\end{itemize}

\chapter{ACADEMIC ADVISING}\label{sec-academic-advising}

The role of the academic advisor is to acquaint students with their
academic options at Coe and assist them in selecting courses that
reflect individual interests and abilities. Advisors also help students
create a four-year education plan that allows students to make
connections between disparate areas of study, and between academic,
co-curricular, and non-academic areas. Students may change advisors at
any time upon request to the Registrar.

First-Year Seminar instructors serve as the primary academic advisors
for first-year students in their respective sections. Students thus see
their advisors frequently during their first term at Coe and have the
opportunity to work closely with them in developing overall programs of
study and long-range goals. After the first term, students may decide to
choose departmental faculty for academic advising or they may continue
to be advised by their First-Year Seminar instructors. Students are free
to speak at any time with professors in their major departments to
answer specific questions regarding requirements and courses in those
departments.

\chapter{FIRST-YEAR EXPERIENCE}\label{sec-first-year-experience}

The student's First-Year Experience at Coe is a deliberate strategy to
engage first-year students across multiple dimensions of college life in
the first year. Components of the First-Year Experience include writing
exercises, various campus events, and the First-Year Seminar (FYS).

The First-Year Seminar is required for all students who have not yet
completed a full-time college term after high school graduation. During
the Fall Term, a variety of First-Year Seminars---topics courses
exploring issues from multiple perspectives---are offered. The seminars
emphasize critical thinking, writing, speaking and research skills.

All First-Year Seminars carry the writing emphasis designation.
First-Year Seminar courses cannot fulfill any distributional, cultural
perspective, or major requirements. Students who drop or fail their
First-Year Seminar are required to complete a replacement course
designated by the Registrar the following term. The completion of this
Spring Term course makes it possible for the student to fulfill the FYS
graduation requirement. For students who fail the Fall Term First-Year
Seminar, successful completion of the Spring Term course also allows the
grade in the FYS to be changed from ``F'' to ``NP'' on the transcript.

Students who begin their college enrollment in the Spring Term must also
complete a course designated as a replacement.

\chapter{WRITING EMPHASIS COURSES}\label{sec-writing-emphasis-courses}

\textbf{Writing Across the Curriculum: Statement of Guiding Principles}

As stated in the Coe College Mission Statement, our reason to exist as
an institution is to ready students intellectually, professionally, and
socially to lead productive and satisfying lives in the global society
of the 21st century. In accordance with this mission, our curriculum
requires that students undertake ``a series of intensive writing
experiences, spread across four years of study.'' This requirement is
known as Writing Across the Curriculum, the guiding principles of which
are detailed below.

As a philosophy, Writing Across the Curriculum asserts that writing is
most effectively learned in context: to varied audiences, with varied
purposes. In adopting this philosophy, Coe College has committed to
making writing-intensive experiences available to students in all
disciplines. Known as ``Writing Emphasis'' credits, these courses are
divergent in subject matter but aligned in their commitment to giving
students content-rich and context-specific writing experiences that
foster a critical flexibility in transferring knowledge about effective
writing to multiple, even unknown, contexts.

Therefore, students who complete the requisite number of writing
emphasis courses will graduate from Coe knowing that writing is both a
means and an end: a method for exploring ideas and deepening one's
knowledge as well as a tool for sharing that knowledge and expressing
one's point of view. Similarly, graduating students will know that the
practice of writing is a recursive process rooted in revision, which
refers both to the reconsideration of one's ideas and to the refinement
of prose, and this process unfolds over a lifetime. Understanding that
students will encounter innumerable writing situations in their lives,
many of which may be unfamiliar to them, they will graduate from Coe
having developed the confidence and self-efficacy necessary to adapt or
draw from their existing knowledge in order to navigate new contexts.

Although many courses at Coe may include writing activities, courses
endorsed by Writing Across the Curriculum are those in which:

\begin{itemize}
\item
  Students are given opportunities to practice revision, whether via
  multiple drafts of a single project or multiple iterations of the same
  type of assignment;
\item
  Students receive instructor feedback on written work to facilitate
  revision; and
\item
  Writing assignments are frequent enough that they are integral to the
  learning throughout the course, enabling students to develop in one or
  more of the following learning outcomes, as appropriate to discipline
  and course objectives:

  \begin{itemize}
  \item
    \textbf{How to write for specific purposes and audiences}

    \begin{itemize}
    \tightlist
    \item
      \emph{Including attention to the ways purposes and audiences shape
      form, mode, voice, method, organization, engagement with and
      citation of research, and/or creativity and imagination}
    \end{itemize}
  \item
    \textbf{How to engage in critical thinking}

    \begin{itemize}
    \tightlist
    \item
      \emph{Particularly how to conduct analysis, how to synthesize
      information, how to interpret and/or use evidence and data, and
      how to present one's ideas coherently and stylistically}
    \end{itemize}
  \item
    \textbf{How to practice, assess, and develop effective habits for
    writing}

    \begin{itemize}
    \tightlist
    \item
      \emph{Specifically, how to read deeply, how to begin and later
      practice selection in research, how to be a skilled reader of
      one's own and others' in-process writing, how to accept and
      interpret feedback, and how to revise writing, all of which are
      parts of the process of learning how to have and develop ideas}
    \end{itemize}
  \item
    \textbf{How to engage in writing as a method for learning and
    discovery}

    \begin{itemize}
    \tightlist
    \item
      \emph{For deepening knowledge, thinking through questions and
      problems, and reflecting on connections and growth}
    \end{itemize}
  \item
    \textbf{How to name and describe one's own knowledge for others}

    \begin{itemize}
    \tightlist
    \item
      \emph{For example, in the form of personal statements for
      application to graduate study, cover letters and resumes for
      entering and advancing in the workforce, or other forms of
      self-summary that facilitate transition from the baccalaureate
      environment through the broader public.}
    \end{itemize}
  \end{itemize}
\end{itemize}

Many writing emphasis sections are offered each year, and, in addition,
the College's rhetoric department offers several interdisciplinary
writing courses designed to guide students learning to write effectively
at the college level. (See course descriptions starting on p.~87).
General Education or major courses that are also designated as writing
emphasis courses may be used to satisfy both requirements.

In this catalog, writing emphasis courses are designated by a (WE) after
the course title, e.g.~``RHE-200 Rhetorical Theory \& Practice (WE).''
In other contexts, the course code may end with a ``W''
(e.g.~RHE-200-W). All sections of such a course, regardless of
instructor, will carry writing emphasis credit.

Fulfilling the Writing Emphasis Requirement:

Only designated WE courses in which a student earns a grade of ``C''
(2.0) or better count toward fulfillment of this requirement.

Only designated courses taken at Coe College count toward fulfillment of
this requirement.

Undergraduate degree-seeking* students

\begin{itemize}
\item
  who start at Coe College or who transfer fewer than 8 course credits
  earned after graduation from high school must complete at least 5
  writing emphasis courses. Of these 5, one will usually be the FYS. Of
  the remaining 4, it is recommended (but not required) that at least 1
  be taken outside the student's intended major(s).
\item
  who transfer at least 8 but fewer than 16 course credits to Coe,
  earned after graduation from high school, must complete at least 3
  writing emphasis courses.
\item
  who transfer 16 or more course credits to Coe, earned after graduation
  from high school, must complete at least 2 writing emphasis credits.
\end{itemize}

For all students, regardless of transfer credit, it is recommended (but
not required) that at least one writing emphasis course be taken in the
upper division within the student's intended major(s). ``Upper
division'' courses are defined in this case as those numbered 300 and
above; upper division writing emphasis courses could include both
scheduled classroom courses and arranged writing-based capstone projects
(such as a thesis). See Departmental Writing Plans for more information
on writing in the major(s).

*Students seeking a second undergraduate degree from Coe should see
Second Baccalaureate Degree p.~13.

\subsubsection{COURSE NUMBERING}\label{course-numbering}

\begin{itemize}
\tightlist
\item
  Courses numbered 100--199 are introductory to the subject.
\item
  Courses numbered 200--299 assume a capacity for the independent
  acquisition of material and generally build on the methods and subject
  matter of 100-level courses.
\item
  Courses numbered 300--399 are typically oriented toward a major or
  minor. They require a strong foundation of knowledge specific to the
  discipline.
\item
  Courses numbered 400--499 are courses in the major or minor that are
  designed to challenge students to integrate discipline-specific
  knowledge in advanced ways. These courses typically contain advanced
  disciplinary coursework, capstone projects, and/or independent
  research.
\end{itemize}

\chapter{GENERAL EDUCATION COURSES}\label{sec-general-education-courses}

Completion of the General Education Program, described here, is required
for all students earning B.A. or B.S.N. degree at Coe College, but not
required of students earning a B.M. degree. An appropriate transferred
course, determined by the Office of the Registrar, with a grade of ``C''
(2.0) or better may be accepted to meet an individual requirement.\\
Advanced Placement and International Baccalaureate courses may not be
used to meet any part of this requirement. No more than two courses with
the same prefix may be used to fulfill the General Education Core
Requirements.

\textbf{A) Liberal Arts Core}

\begin{enumerate}
\def\labelenumi{\arabic{enumi}.}
\tightlist
\item
  A total of at least one course credit in the \textbf{Fine Arts Core}
  (courses with a prefix of ARH, ART, FLM, MU, MUA, THE)
\item
  A total of at least two course credits in the \textbf{Humanities Core}
  (courses with a prefix of AAM, AMS, CLA, COM, CRW, ENG, FRE, GER, GRK,
  HIS, HUM, LTN, JPN, PHL, REL, RHE, or SPA)
\item
  A total of at least two course credits in the \textbf{Natural Sciences
  and Mathematics Core} (courses with a prefix of BIO, CHM, CS, DS, MTH,
  PHY, STA, one of which must be a lab science with a prefix of BIO,
  CHM, or PHY)
\item
  A total of at least two course credits in the \textbf{Social Sciences
  Core} (courses with a prefix of ANT, ECO, POL, PSY, or SOC)
\end{enumerate}

\textbf{B) Diverse Cultural Perspectives (DCP) Core}

A liberally educated person should have some knowledge of other cultures
and some tools to aid in seeing one's own culture from other
perspectives.

The Diverse Cultural Perspectives courses (Non-Western Perspectives
(course number ends in 6), United States Pluralism (course number ends
in 7), Diverse Western Perspectives (course number ends in 8) help
students to understand their own cultural identities and to develop
appreciation for the range of different cultures to be found in the
world, in the nation, and on campus. These courses explore other
cultures in their own terms and as they interact with American culture.
As a group, they encourage reflection on different ways in which
cultural identities are formed, expressed and contested. Students are
encouraged to combine these courses with the study of a world language
and with study abroad.

Completion of one of the following options (see p.~34-38 for
descriptions and approved courses):

\begin{itemize}
\tightlist
\item
  Non-World Language Option (all of the following):

  \begin{itemize}
  \tightlist
  \item
    Any DCP course (course number that ends in 6, 7, 8)
  \item
    Non-Western Perspectives (course number that ends in 6) (A
    semester-long study abroad experience can fulfill the Non-Western
    Perspective and DCP requirement.)
  \item
    United States Pluralism (course number that ends in 7) (A
    semester-long U.S. off-campus study experience can fulfill the U.S.
    Pluralism and DCP requirement.)
  \end{itemize}
\item
  Elementary World Language Option (all of the following):

  \begin{itemize}
  \tightlist
  \item
    Any DCP course (course number that ends in 6, 7, 8)
  \item
    Two elementary world language courses in the same language not
    previously studied (may also fulfill one of the Humanities Core
    courses)
  \end{itemize}
\item
  Intermediate World Language Option (all of the following):

  \begin{itemize}
  \tightlist
  \item
    Any DCP course (course number that ends in 6, 7, or 8)
  \item
    One intermediate (-215) or above World Language course (\emph{may
    also fulfill one of the Humanities Core course requirements})
  \end{itemize}
\end{itemize}

\section{Diverse Cultural Perspectives: Non-Western
Perspectives}\label{sec-diverse-cultural-perspectives-non-western}

The Non-Western Perspectives (NWP) group includes courses in which a
preponderance of the content analyzes human experience from the
perspectives of diverse peoples outside of Western culture. Such
cultures often are characterized by values and beliefs different from
those of the United States and Western Europe.

\begin{itemize}
\tightlist
\item
  ANT 116 Cultural Anthropology
\item
  ANT 286 Topics in Anthropology:NWP
\item
  ARH 106 World Art
\item
  ARH 296 Topics in Art History:Global Persp
\item
  ASC 106 May Term in Asia
\item
  ASC 176 China and Japan
\item
  ASC 186 Modern South Asia
\item
  ASC 196 Modern South East Asia
\item
  +++MISSING INFO: c.asc216.long +++
\item
  BUS 446 International Business Management
\item
  BUS 466 Adv Top Mrktg:non-west persp
\item
  BUS 476 Ad Top Mgmt:non-west persp
\item
  COM 236 Intercultural Communication
\item
  ECO 336 Divergent Economic Growth
\item
  ECO 436 Econ Development
\item
  ECO 446 International Econ
\item
  ENG 146 Intro Postcolonial Literature
\item
  ENG 206 Gender and Literature: NWP
\item
  FRE 146 French Literature Translation:NWP
\item
  FRE 446 Colonial \& Multicultural Narratives
\item
  FSA 146 Turkey: History \&Culture
\item
  GS 136 Gender in Non-Western World
\item
  HIS 136 East Asian Civilization
\item
  HIS 216 History of Modern Korea (WE)
\item
  HIS 246 History of Modern China
\item
  HIS 256 History of Modern Japan
\item
  HIS 276 The ``Discovery'' of America: Clash
\item
  HIS 286 Modern Middle East
\item
  HIS 306 Revolution, Social Struggle, Testim
\item
  HIS 316 Topics in History:Non-Western Persp
\item
  HIS 466 Seminar Modern East Asian History
\item
  IS 116 Intro to International Studies
\item
  IS 126 HumanRightsBurmeseMigrant
\item
  IS 136 May Term in Mongolia
\item
  JPN 106 Images Foreign Culture
\item
  MU 166 Topics in Music:Non-Western Perspct
\item
  PHL 206 Buddhist Thought
\item
  POL 266 Latin American Politics
\item
  POL 276 African Politics
\item
  POL 286 Asian Politics
\item
  POL 296 Topics Pol Sci: Non-West Persp
\item
  POL 386 International Development
\item
  REL 106 Eastern Religions
\item
  REL 116 Buddhism
\item
  REL 136 Religions of China
\item
  REL 196 Hinduism
\item
  REL 206 Buddhist Thought
\item
  REL 226 Religions of China:Daoism
\item
  REL 236 Zen Buddhism
\item
  REL 296 Topics in Religion NWP
\item
  REL 306 Comparative Religion
\item
  REL 336 Tibetan Buddhist Cultrue
\item
  +++MISSING INFO: c.rel386.long +++
\item
  RHE 146 Creative Nonfiction:Global Perspect
\item
  SOC 226 Gender and Globalizaton
\item
  SOC 236 Topics in Sociology NWP
\item
  SPA 336 Hispanic Life/Cult-Latin America
\item
  SPA 446 Latin Am. \& Spanish Short Stories
\item
  SPA 486 Topics in Hispanic Lit:Latin Amer
\item
  THE 486 Spc Top Theatre or Film: NWP
\item
  WSH 286 Topics in Washington, D.C:NWP
\end{itemize}

\section{Diverse Cultural Perspectives: United States
Pluralism}\label{sec-diverse-cultural-perspectives-united-states}

The United States Pluralism (USP) group includes courses in which a
preponderance of the content addresses one or more of the groups within
the United States whose values, beliefs, and experiences differ from or
oppose those of the majority culture. These courses increase students'
knowledge of the history of such groups; of the ways members of these
groups have experienced democracy and culture in America differently
because of factors like social class, race, gender, and religion; and of
reform movements like feminism and civil rights, through which such
groups have attempted to achieve social and economic equality.

\begin{itemize}
\tightlist
\item
  AAM 107 Intro to African American Studies
\item
  AAM 137 African American Literature\\
\item
  AAM 217 Sport and Black Culture\\
\item
  AAM 227 Blackness \& Identity in America
\item
  AAM 287 Topics in African American Studies
\item
  AAM 367 Topics in AfricanAmericanLiterature
\item
  AAM 387 Adv Topics in African American Stud
\item
  AAM 447 Drtd Learn in African American Stds
\item
  AAM 457 Drtd Learn in African American Stds
\item
  AAM 467 Seminar in African American Lit
\item
  ARH 107 Gender and Art\\
\item
  ARH 297 Topics in Art History: US Pluralism\\
\item
  ARH 307 Modern and Contemporary Art
\item
  BUS 387 Adv Top:Human Res Mgt
\item
  BUS 437 Strategic Compensation\\
\item
  BUS 457 Employment and Discrimination Law
\item
  BUS 467 Consumer Behavior
\item
  COM 157 Introduction to Media Analysis
\item
  COM 237 Interpersonal Communication
\item
  COM 337 Persuasion\\
\item
  COM 357 Sex, Race, \& Gender in Media
\item
  COM 437 Special Topic Applied Communication
\item
  COM 447 Special Topics Production
\item
  COM 457 Special Topics in Media Studies
\item
  COM 467 Special Topics in Public Discourse
\item
  ECO 237 Labor Economics\\
\item
  ECO 247 Health Economics\\
\item
  ECO 447 Urban Economics\\
\item
  ECO 457 US Econ History\\
\item
  EDU 117 Exceptional Learners\\
\item
  EDU 187 Human Relations\\
\item
  EDU 237 English Language Learners
\item
  EDU 247 Foundations of Reading\\
\item
  ENG 107 Exploring Literature:US Pluralism
\item
  ENG 117 Asian American Literature
\item
  ENG 127 Social Justice and Literature
\item
  ENG 137 African American Literature
\item
  ENG 157 Latinx/Chicanx Literature
\item
  ENG 207 Gender \& Lit:US Pluralism
\item
  ENG 327 Literature of American Renaissance
\item
  ENG 337 American Realism \& Naturalism
\item
  ENG 347 Study in Modern or Contemp Amer Lit
\item
  ENG 357 Studies in Latinx/ChicanxLiterature
\item
  ENG 367 Studies in African Am Literature
\item
  ENG 467 Seminar inLit:USPluralism
\item
  +++MISSING INFO: c.evs137.long +++
\item
  GS 107 Intro Gender \& Sexuality Studies
\item
  GS 127 Dress, Gender, and Identity
\item
  GS 247 Gender \& Sexuality StudiesSymposium
\item
  GS 327 Thry\&Mthds/Gender \&Sexuality Stdy
\item
  GS 387 Topics: Gender \& Sexuality Studies
\item
  HIS 217 American War in Vietnam
\item
  HIS 227 American Civil War\\
\item
  HIS 257 Native American History
\item
  HIS 297 Women in America\\
\item
  HIS 317 Topics in History:US Pluralism\\
\item
  HIS 347 African American History
\item
  HIS 387 American Colonial History
\item
  KIN 347 Adapted Physical Education
\item
  MU 157 Introduction to Jazz History
\item
  NUR 137 Human Sexuality\\
\item
  NUR 297 Parent Child Relationships\\
\item
  NUR 387 Alternative Therapies for Hlth/Heal
\item
  PHL 277 Philosophy of Gender \& Race
\item
  POL 207 Religion \& American Politics
\item
  POL 277 Women \& Poltics in US
\item
  PSY 137 Human Sexuality\\
\item
  REL 217 Religion in America\\
\item
  RHE 137 Creative Nonfiction U.S.-Pluralism
\item
  RHE 257 Environmental Rhetoric\\
\item
  RHE 377 Cultural Studies\\
\item
  SOC 107 Introductory Sociology\\
\item
  SOC 207 Sociology of the Family\\
\item
  SOC 217 Sociology of Religion
\item
  SOC 237 Topics in Sociology:U S Pluralism
\item
  SOC 247 Sociology of Race
\item
  SOC 417 Sociology of Sex \& Sexuality
\item
  SPA 457 US LatinX Literature
\end{itemize}

\section{Diverse Cultural Perspectives: Diverse Western
Perspectives}\label{sec-diverse-cultural-perspectives-western}

The Diverse Western Perspectives (DWP) group includes courses in which a
preponderance of the content addresses one or more subgroups of the
Western world outside of the United States and the ways in which they
experience Western culture. These courses increase students' knowledge
of the history of particular groups and the ways they have interacted
with Western values. They typically address issues of difference and
conflict between and within Western cultures by examining the influence
of factors such as class, race, gender, and religion.

\begin{itemize}
\tightlist
\item
  ANT 288 Topics Anthropology/Archaeology:DWP
\item
  ANT 488 Adv Top Anthro/Archaeo:DWP
\item
  ARH 128 Introduction to Art History
\item
  ARH 218 The World of Renaissance Art
\item
  ARH 248 Baroque, Rococo, and Neoclassicism
\item
  ARH 268 History of Architecture
\item
  ARH 298 Topics in Art History:Div West Pers
\item
  CLA 108 Images of Foreign Culture
\item
  ENG 108 Exp Lit:Diverse Western Perspective
\item
  ENG 208 Gender \& Lit:DWP
\item
  ENG 378 Studies in Transatlantic Literature
\item
  ENG 388 Romantic Literature
\item
  FRE 148 French Literature Translation:DWP
\item
  FRE 158 France \& Francophone World
\item
  HIS 208 The First World War (WE)
\item
  HIS 218 The Second World War (WE)
\item
  HIS 238 Modern France
\item
  HIS 248 The French Revolution
\item
  HIS 268 Latin America
\item
  HIS 288 Renaissance \& Reformation
\item
  HIS 308 Legacies of the Cold War inLatin Am
\item
  HIS 318 Topics in History :Div West Persp
\item
  HIS 328 Modern France
\item
  MU 458 Music History \& Literature III
\item
  NUR 268 Cult Diver \& Health
\item
  PHL 128 Morality \& Moral Controversies
\item
  PHL 138 Freedom, State, and Society
\item
  POL 108 Introduction to Politics
\item
  POL 248 Political Violence and the Violent
\item
  POL 258 World Politics
\item
  POL 298 European Politics
\item
  POL 398 Religion \& World Politics
\item
  PSY 208 Gender Psychology
\item
  REL 108 Western Religions
\item
  REL 128 Judaism
\item
  REL 138 Modern Judaism
\item
  REL 148 Islam
\item
  REL 178 Christianity
\item
  REL 278 Mysticism
\item
  REL 338 Modern Religious Thought
\item
  SOC 238 Topics in Soc Div West Perspectives
\item
  SOC 328 Urban Sociology
\item
  SOC 338 Political Sociology
\item
  SPA 148 Spanish Literature in Translation
\item
  SPA 258 Spanish Lang Learn in Spain
\item
  SPA 338 Hispanic Life/Culture:Europe
\item
  SPA 418 Gender \& Sexuality in Hispanic Wrld
\item
  SPA 428 Indigeneity, Blackness, \& EthnicLit
\item
  SPA 458 Travel Writing \& Transatlantic Lit
\item
  THE 118 Theatre \& Arts in Serbia
\item
  THE 228 History of Theatre and Drama I
\item
  THE 238 History of Theatre and Drama II
\item
  THE 288 History of Dress
\item
  THE 488 Special Topics in THE/ FLM
\end{itemize}

\section{Independent Studies}\label{sec-independent-studies}

The one-credit independent study or directed readings students are
expected to complete a minimum of 140 hours of academic work, including
meeting with faculty members and independent work between meetings.

\section{Internships}\label{sec-internships}

The internship is a work or volunteer experience in the context of an
independent academic investigation of site-related issues and personal
aptitudes, values, and goals. The one-credit internship includes a
minimum of 140 hours of on-site or remote experience and the required
documentation and/or academic journal or paper as determined by the
faculty member. The details of the academic component are determined by
prior arrangement with the faculty internship advisor.

Ordinarily internships are completed during the academic year as one of
the 32 credits for graduation, or over the summer for a credit-bearing
or non-credit bearing-practicum fulfillment activity. In unusual
circumstances where an internship presents an opportunity to extend the
educational component of the experience significantly, an internship may
earn two credits. Application for non-departmental, two-credit
internships requires consultation with the Internship Faculty Advisor
and approval of the Committee on Petitions; departmental two-credit
internships, when permitted, are overseen by the department.

A combined maximum of 2.0 course credits may be counted toward the 32
credits required for graduation through Internships or Community-Based
Project (see p.~87). Many internships are completed in the Cedar Rapids
area; however, it is permissible to complete an internship outside of
the area during the summer, or done remotely if authorized by the
Internship Faculty Advisor.

Students interested in internships should consult with their Career
Specialist in C3: Creativity, Careers, Community, as well as with the
appropriate academic department who will be overseeing the internship
experience.

Each internship must include the consent of an Internship Faculty
Advisor and completion of the internship request form housed on the
College's online platform for internships.

Courses with an INT prefix and those on this list count as internship
credit:

\begin{itemize}
\tightlist
\item
  AAM 494 Internship in African American Stds
\item
  ANT 494 Internship in Anthropolgy
\item
  ARH 494 Internship in Art History
\item
  ARH 494 Internship in Art History
\item
  ART 494 Internship in Art
\item
  AT 494 Internship in Athletic Training
\item
  BIO 494 Internship in Biology
\item
  BUS 494 Internship in Business
\item
  CHM 494 Internship in Chemistry
\item
  COM 494 Internship in Journalism/Communicat
\item
  CRW 494 Internship in Creative Writing
\item
  CS 494 Internship in Computer Science
\item
  DS 494 Internship in Data Science
\item
  EDU 494 Internship in Education
\item
  ENG 494 Internship in English
\item
  FLM 494 Internship in Film
\item
  FRE 494 Internship in French
\item
  HIS 494 Internship in History
\item
  HSS 494 Health Professions Externship
\item
  INT 115 May Term: Topics 2
\item
  INT 494 Internship
\item
  INT 499 Summer Internship
\item
  KIN 494 Internship in Kin, Health \& Rec
\item
  MTH 494 Internship in Mathematics
\item
  NUR 494 Internship in Nursing
\item
  NYT 394 Internship in New York City
\item
  PHL 494 Internship in Philosophy
\item
  PHY 494 Internship in Physics
\item
  POL 494 Internship in Political Science
\item
  PR 494 Internship in Public Relations
\item
  PSY 494 Internship in Psychology
\item
  REL 494 Internship in Religion
\item
  RHE 494 Internship in Writing
\item
  SOC 494 Internship in Sociology
\item
  SMT 494 Internship in Sports Management
\item
  SPA 494 Internship in Spanish
\item
  THE 494 Internship in Theatre Arts
\item
  WSH 494 Washington Experience
\end{itemize}

Students completing internships that are not department specific should
register for one of the INT-494 or INT-499 courses on p.~49 (see the
\textbf{\emph{Coe Student Accounts Handbook}} for fee).

\section{Practicum}\label{sec-practicum}

A practicum experience is required of all students for all undergraduate
degrees, except those earning second degrees.

Typically completed in the student's junior or senior year, all practica
are experiences that integrate academic components with career or other
life goals and are significant educational exercises outside the
classroom. A practicum experience can consist of an internship,
off-campus study, community-based project, honors project, or some other
kind of independent activity.

Depending upon the type selected, some practica are graded A--F, while
others are P/NP. Some practica are credit bearing, while others are not.
In some instances, the practicum must be approved by the student's major
department.

\begin{enumerate}
\def\labelenumi{\arabic{enumi}.}
\tightlist
\item
  Full-Term (16-week) Off-Campus Study
\item
  Wilderness Field Station Summer Courses
\item
  Crimson Fellows Thesis or Crimson Fellows Project, etc. as stated
\item
  Independent Project (in list of courses that follows starred courses *
  require department approval for practicum credit):

  \begin{enumerate}
  \def\labelenumii{\arabic{enumii}.}
  \setcounter{enumii}{4}
  \tightlist
  \item
    †Internship (see a complete listing of internships on p.~39)
  \item
    †Community-Based Project (see course description on p.~87)
  \end{enumerate}
\end{enumerate}

†A maximum of two course credits earned through any combination of
Internships and Community-Based Projects may be included in the 32
course credits required for graduation.

\begin{itemize}
\tightlist
\item
  AAM 444 Ind Study-Afr-Am St
\item
  ANT 205 Archaeological Field Schl
\item
  ANT 444 Independent Study: Anthropology
\item
  ANT 474 Research Participation:Anthropology
\item
  ARH 444 Independent Study: Art History
\item
  ARH 474 Senior Seminar II
\item
  ART 394 Directed Learning in Art
\item
  ART 444 Independent Study in Art
\item
  ART 474 Senior Seminar II \& Senior Exhibit
\item
  +++MISSING INFO: c.at40.long +++ (\emph{successful completion of
  sequence of AT-20/-30 and -40 required to receive full credit})
\item
  BIO 115 Marine Biology
\item
  BIO 444 Independent Study
\item
  BIO 454 Research Participation
\item
  BIO 462 Advanced Research Lab I
\item
  BUS 444 Ind Study-Bus Admin
\item
  BUS 454 Research in Business
\item
  CHM 444 Independent Study Chemistry
\item
  CHM 454 Undergraduate Summer Research
\item
  COM 394 Directed Learning in Communication
\item
  COM 444 Independent Study in Comm Studies
\item
  CRW 112 Advanced Literary Magazine Editing
\item
  CRW 394 Directed Studies in Creative Writin
\item
  CRW 492 Manuscript Workshop
\item
  CS 444 Ind Study-Comp Sci
\item
  CS 454 Research in Computer Science
\item
  DS 444 Independent Study in Data Science
\item
  DS 454 Research in Data Science
\item
  ECO 444 Ind Study-Economics
\item
  ECO 454 Research in Economics
\item
  EDU 215 Practicum in Education
\item
  EDU 481 Stu Tchg Sec: ART
\item
  EDU 482 Stu Tchg Sec: Phys Education
\item
  EDU 483 Std Teaching Elem: ART
\item
  EDU 485 Std Teaching Elem: Phys Ed
\item
  EDU 489 Student Teaching Sr HS
\item
  EDU 490 Student Teaching Jr HS
\item
  EDU 491 Student Teaching 4-6
\item
  EDU 492 Student Teaching K-3
\item
  EDU 444 Ind Study-Tchr Ed
\item
  ENG 394 Directed Learning in English
\item
  ENG 454 Honors Research
\item
  FLM 442 Independent Study in Film
\item
  FLM 464 Seminar in Film II: (0.5 cc) and FLM 474 Senior Seminar II in
  Film \& Senior (0.5 cc)
\item
  FRE 394 Directed Learning in French
\item
  FRE 444 Ind Study-French
\item
  HIS 444 Ind Study-History
\item
  KIN 444 Ind Study-KIN
\item
  MTH 444 Ind Study-Math
\item
  MTH 454 Research in Mathematics
\item
  MU 421 Student Teaching Elementary Music
\item
  MU 422 Student Teaching Secondary Music
\item
  MU 444 Ind Study-Music
\item
  MUA 490 Senior Recital
\item
  NUR 444 Ind Study-Nursing
\item
  NUR 455 Leadership \& Cont Issues in Nursing
\item
  PHL 394 Directed Learning in Philosophy
\item
  PHL 444 Ind Study-Philos
\item
  PHL 490 Philosophy Colloquium (0.0 cc)
\item
  PHY 255 Advanced Laboratory I
\item
  PHY 355 Advanced Laboratory II
\item
  PHY 444 Ind Study-Physics
\item
  POL 444 Ind Study-Pol Sci
\item
  PSY 354 Research Participation
\item
  PSY 444 Ind Study-Psychology
\item
  PSY 455 Advanced Experimental Psychology
\item
  REL 394 Directed Learning in Religion
\item
  REL 444 Ind Study-Relig
\item
  RHE 394 Drtd Learning in Writing \& Rhetoric
\item
  RHE 444 Independent Study in Writing (WE)
\item
  RHE 490 Publications Practicum
\item
  SCJ 444 Independent Study in Social and Cri
\item
  SOC 365 Research Participation I
\item
  SOC 444 Ind Study-Soc
\item
  SPA 394 Directed Learning: Spanish
\item
  SPA 444 Ind Study-Spanish
\item
  THE 442 Adv Proj-Design/Tech Production
\item
  THE 444 Ind Study-Theatre
\item
  THE 452 Advanced Projects in Acting
\item
  THE 462 Advanced Projects in Directing
\end{itemize}

\chapter{SPECIAL PROGRAMS AND
OPPORTUNITIES}\label{special-programs-and-opportunities}

\section{Clinical Laboratory Sciences/Medical
Technology}\label{clinical-laboratory-sciencesmedical-technology}

In cooperation with the St.~Luke's Methodist Medical Laboratories in
Cedar Rapids, or upon arrangement with other accredited laboratories and
the approval of the College, Coe offers a four-year course leading to a
Bachelor of Arts degree and registration as a Clinical Laboratory
Scientist/Medical Technologist.

The first three years are spent in residence at Coe, where candidates
must complete all of the requirements for the B.A. degree, including
general education requirements and an approved major. The minimum
requirements of the Clinical Laboratory Sciences/Medical Technology
program in biology and chemistry are five course credits in each field
and at least one course in mathematics. The fourth year is a full
calendar year spent at St.~Luke's or another accredited medical
laboratory approved by Coe.

The St.~Luke's Hospital Medical Laboratory is approved as a school of
clinical laboratory sciences/medical technology by the Committee on
Allied Health Education and Accreditation of the American Medical
Association. Candidates completing the course are eligible to take the
certification examinations of the American Society of Clinical
Pathologists and the National Certification Agency and, if approved, may
practice anywhere in the United States.

\section{Crimson Fellows Program}\label{crimson-fellows-program}

Steffens, Westberg (Program Directors).

To graduate as a Crimson Fellow, a student must earn at least a 3.3
cumulative grade point average for all courses taken at Coe College and
complete all of the following:

\begin{enumerate}
\def\labelenumi{\arabic{enumi}.}
\tightlist
\item
  CFP 104 Topics in Crimson Fellows Program (0.2 cc)
\item
  +++MISSING INFO: c.CFP205.long +++
\item
  +++MISSING INFO: c.CFP301.long +++ (0.5 cc)
\item
  +++MISSING INFO: c.CFP302.long +++ (0.5 cc)
\item
  +++MISSING INFO: c.CFP401.long +++ (0.5 cc)
\end{enumerate}

Applications to the Crimson Fellows Program are accepted and reviewed on
a rolling basis.

Completed applications are assessed on ACT/SAT score, the high school
transcript, teacher recommendations, and the student essay. Although
there are no minimum thresholds, the historical average ACT score has
been over 27 with an average high school GPA above 3.70. Emphasis will
be placed on the student essay.

\begin{itemize}
\tightlist
\item
  CFP 104 Topics in Crimson Fellows Program

  \begin{itemize}
  \tightlist
  \item
    Reading and discussion of one or more classic texts from across
    intellectual disciplines. Students are expected to demonstrate
    mastery of the material and actively engage in class discussions.
    S/U basis only. May be taken more than once for credit, provided the
    topics are substantially different. Prerequisite: consent of
    instructor. (0.2 course credit)
  \end{itemize}
\item
  CFP 145 CFP: Culture and Revolution

  \begin{itemize}
  \tightlist
  \item
    Designed to examine cultures that subsequently undergo revolutionary
    change. The first part is devoted to giving a sense of the culture
    that is destroyed by the revolution, and the second part to the
    forces that lead to the revolution. This course focuses on
    masterworks in history, political science, philosophy, and
    literature.
  \end{itemize}
\item
  CFP 155 Style and Transformation in the Art

  \begin{itemize}
  \tightlist
  \item
    Focuses on periods during which the arts undergo a major
    transformation. The class studies both the artistic ideas that are
    being changed and the characteristic styles that result from these
    changes. The scope of this course may include literature, painting,
    music, and architecture.
  \end{itemize}
\item
  CFP 175 Continuity \&Transition Non-West Soc

  \begin{itemize}
  \tightlist
  \item
    Focuses on the great traditions in non-western cultures. By
    selectively dealing with the major traditions present in such
    cultures, the class deals with the dynamics of continuity and
    transition, which are crucial to understanding non-western
    societies.
  \end{itemize}
\item
  CFP 184 Topics in Scientific Inquiry

  \begin{itemize}
  \tightlist
  \item
    Designed to lead students to an intimate understanding of how the
    scientific process works and how scientific thought develops.
    Historical readings and discussions develop students' understanding
    of the course topic sufficiently for them to focus on particular
    scientific questions. Experimental approaches to these questions are
    discussed and developed into research projects. Results are shared
    and integrated, providing group members with greater knowledge of
    the course topic and an acute awareness of the process and
    limitations of science.
  \end{itemize}
\item
  CFP 205 Introduction to Engaged Scholarship

  \begin{itemize}
  \tightlist
  \item
    Establishes the foundation of the Fellows experience. Fellows write
    essays on a series of connected subjects designed to encourage
    critical thinking and reflection. The course includes discussions on
    education reform ideas, action research, student-driven education,
    the passion of life-long learning, global citizenship, community
    advocacy and civic engagement. Grounded in critical pedagogy, the
    course introduces students to the ideas of engaged scholarship.
    Students take command of their own education, developing their sense
    of agency, selfdiscipline, initiative, and self-direction.
    Prerequisite: admission to the Crimson Fellows Program or consent of
    instructor. (Offered Spring Term)
  \end{itemize}
\item
  CFP 301 Communicating Across Disciplines I

  \begin{itemize}
  \tightlist
  \item
    Focuses on developing competencies in communication across multiple
    audiences and perspectives. Students engage content/material related
    to their understanding and development of a year-long research
    project. The central objective of the course is to enhance student
    collaboration and communication skills. Prerequisite: Introduction
    to Engaged Scholarship (CFP-100) or consent of instructor. (0.5
    course credit)
  \end{itemize}
\item
  CFP 302 Communicating Across Disciplines II

  \begin{itemize}
  \tightlist
  \item
    Students work in multi-disciplinary teams to address a topic of
    local and/or global social interest. Prerequisite: Communicating
    Across Disciplines I (CFP-301) or consent of instructor. (0.5 course
    credit)
  \end{itemize}
\item
  CFP 401 Action Research I

  \begin{itemize}
  \tightlist
  \item
    Students work in multi-disciplinary teams to conduct a year-long
    project that addresses a specific challenge/issue for a
    local/regional community partner (e.g.~sustainability, water
    quality, soil degradation, health care, partner violence, refugees,
    immigration, civil rights, etc.). Faculty provide specific content
    related to working with community groups/agencies at the onset of
    the course. Prerequisite: Communicating Across Disciplines II
    (CFP-302) or consent of instructor. (0.5 course credit)
  \end{itemize}
\end{itemize}

\subsection{Crimson Fellows Projects \& Graduating with
Distinction}\label{crimson-fellows-projects-graduating-with-distinction}

To graduate with Distinction, at the time of graduation a student must
have:

\begin{enumerate}
\def\labelenumi{\arabic{enumi}.}
\tightlist
\item
  earned at least a 3.2 cumulative grade point average for all courses
  taken at Coe College,
\item
  earned at least a 3.5 GPA (or higher if set higher in the program in
  which you are pursuing distinction) in the courses taken toward the
  major or minor in which distinction is sought,
\item
  earned at least 14 course credits of graded courses at Coe College,
\item
  satisfactorily completed a Crimson Fellows project in a major or
  minor,
\item
  completed the ``Graduating with Distinction Form'' found on My.Coe and
  submitted it to the Office of the Registrar no later than March 15 of
  the Spring Term before graduation,
\item
  submitted a thesis or project artifact, approved by the majority of
  the student's Crimson Fellows Program examining committee, to the
  Director of Library Services no later than Reading Day of Spring Term.
\end{enumerate}

\subsection{Latin Honors}\label{latin-honors}

Cum laude is awarded to all graduating seniors with a cumulative GPA of
3.60 or higher. Magna cum laude is awarded to graduating seniors with a
cumulative GPA of 3.80 or higher who have completed an honors project.
Summa cum laude is awarded to graduating seniors with a cumulative GPA
of 3.98 or higher who have completed an honors project.

\subsection{Dean's List}\label{deans-list}

Special recognition is given to students who show exceptional academic
performance during Fall and Spring Terms. The designation ``Dean's
List'' is awarded a student if, at the end of a given grade reporting
period, the student:

\begin{enumerate}
\def\labelenumi{\arabic{enumi}.}
\tightlist
\item
  was enrolled as a full-time, degree-seeking student;
\item
  earned at least a 3.5 GPA for the grading period, having no incomplete
  marks, no repeat courses, and at least three letter graded courses;
  and
\item
  ranked in the top ten percent of the student body for that grading
  period.
\end{enumerate}

\section{Cross-Registration with Mount Mercy
University}\label{cross-registration-with-mount-mercy-university}

This agreement:

\begin{enumerate}
\def\labelenumi{\arabic{enumi}.}
\tightlist
\item
  Covers only courses that are not offered at Coe College in the same
  term unless a time conflict exists that cannot be resolved.
\item
  Is permitted on a space-available basis two weeks after the regular
  registration at Mount Mercy.
\item
  Holds students subject to administrative rules of the host institution
  for the courses taken.
\item
  Requires the student to register at both institutions.
\item
  Requires that a student be full-time and degree seeking in the term of
  the request and for at least one previous term at Coe College.
\end{enumerate}

Coe students wishing to enroll at Mount Mercy University may not be on
academic probation (see p.~75) and may not have been dismissed from Coe
College. Prior to registering, students must submit a cross-registration
request form to the Registrar, who grants approval to students wishing
to register at Mount Mercy University. If the course is to be counted
toward a major or minor, the approval of the appropriate Coe department
chair is also required. Declarations of Pass/Not Pass options are made
at Coe according to Coe policies. Both course credit and the letter
grade given at Mount Mercy are recorded in the student's permanent
record, as well as the fact that the course was taught at Mount Mercy.
Under the agreement no additional fees are charged for cross
registration, although the sum of the credits registered at both
institutions are used to determine full-time status and/or the need to
petition to take 5.0 course credits or more.

A cross-registered student missing a class at a cooperating college
because of calendar differences shall not be penalized for missing the
class. The student, however, is responsible for making up any work
missed in the class. Before registering for a course at Mount Mercy
under this agreement, a student must complete the Mount Mercy--Coe
College Cross Registration Form to be processed to ensure that all
stipulations of the agreement are met.

\section{English As A Second Language
Program}\label{english-as-a-second-language-program}

Drexler (Director), Welsh.

English as a Second Language (ESL) programs at Coe College are designed
to help study abroad and undergraduate international students assimilate
into the College and local community through English language
instruction.

Coe College offers the following programs:

\begin{itemize}
\tightlist
\item
  Intensive English Language program (IELP). IELP accepts students whose
  TOEFL ITP score falls lower than 500. IELP students enroll in 18 hours
  of English language courses per week for one or two terms. Students
  who successfully complete IELP may matriculate to the College.
\item
  English Academic Bridge (ELAB) program. ELAB accepts students whose
  TOEFL ITP score falls between 500 and 520, the Admission requirement.
  ELAB students are admitted as undergraduates with the following course
  requirements for their first term at the College: one ESL course, one
  Coe-credit course designed for English language learners, First-Year
  Seminar, and their choice of another Coe-credit course.~
\item
  Short Intensive English Language program (SIELP). SIELP accepts study
  abroad students for seven weeks during the Spring Term. SIELP students
  enroll in 18 hours of English language course per week and live on
  campus.
\item
  Summer Academic Orientation program (SAOP). SAOP accepts students who
  intend to study abroad or become undergraduates in colleges and
  universities in the United States or Canada. SAOP students enroll in
  20 hours of ESL courses per week for the first week of August.
\end{itemize}

In addition to these programs, Coe international students who are not
enrolled in an ESL program may register for ESL courses and work with
ESL faculty on a one-on-one basis to support~their language needs.

\begin{itemize}
\tightlist
\item
  ESL 195 Reading Writing Workshop

  \begin{itemize}
  \tightlist
  \item
    Strengthens critical reading, academic writing, and reasoning skills
    by engaging with a variety of texts. Students identify, challenge,
    and write arguments by practicing summary, analysis, paraphrase, and
    response to published work as well as work of their peers. May be
    taken more than once with consent of the program director provided
    the topics are substantially different. Prerequisite: enrolled in
    ELAB or undergraduate program with appropriate TOEFL score or
    consent of instructor.
  \end{itemize}
\item
  ESL 180 Topics in Listening

  \begin{itemize}
  \tightlist
  \item
    Develops academic listening and note-taking skills with a focus on
    micro listening such as listening for numbers, word stress, thought
    groups, and accent differences as well as macro listening like
    listening for headings, transition words and phrases, and
    distinguishing main ideas from details. May be taken more than once
    with consent of the program director provided the topics are
    substantially different. Prerequisite: enrolled in ELAB or
    undergraduate program with appropriate TOEFL score or consent of
    instructor.
  \end{itemize}
\item
  ESL 185 Topics in Speaking

  \begin{itemize}
  \tightlist
  \item
    Develops speaking fluency through vocabulary building by using
    language in authentic contexts and practicing pronunciation by
    studying language prosody (intonation and rhythm) and morphemes
    (minimal parts of language). May be taken more than once with
    consent of the program director provided the topics are
    substantially different. Prerequisite: enrolled in ELAB or
    undergraduate program with appropriate TOEFL score or consent of
    instructor.
  \end{itemize}
\item
  ESL 190 Topics in Structure

  \begin{itemize}
  \tightlist
  \item
    Develops grammar with the goal of using accurate grammar in
    presentations and written work. Students study various grammar
    points, practice grammar interactively, and integrate learned
    grammar into presentations and written work on a wide spectrum of
    topics. May be taken more than once with consent of the program
    director provided the topics are substantially different.
  \end{itemize}
\item
  ESL 210 Seminar in Culture

  \begin{itemize}
  \tightlist
  \item
    Develops language skills and explores American culture through
    texts, art, film, and community engagement. May be taken more than
    once with consent of the program director provided the topics are
    substantially different. Prerequisite: enrolled in ELAB or
    undergraduate program with appropriate TOEFL score or consent of
    instructor.
  \end{itemize}
\end{itemize}

\section{Pre-Professional Programs}\label{pre-professional-programs}

\subsection{Pre-Law}\label{pre-law}

J. Christensen (Program Director).

The cooperative 3+3 program allows qualified undergraduates from Coe
College to earn both a bachelor's degree and a law degree in six years
through partnership with the University of Iowa College of Law.

If eligible, students admitted under the 3+3 program will receive three
years of financial aid and pay Coe College tuition for the first three
years. Qualified undergraduates from Coe must have satisfied all
graduation requirements with the exception of the 32-course credit
requirement by the conclusion of their junior year for admission into
the College of Law. Students are not eligible for campus housing or
extracurricular activities at Coe during their fourth year. During the
fourth year of the program, which is the first year of law school,
students will pay tuition only to Iowa Law and apply for financial aid
through the University of Iowa. Credits earned during the first year of
law school at Iowa, which would have been their senior year at Coe, will
also apply to their undergraduate degree to complete the final credit
requirements at Coe. At the end of their fourth year of study, students
in the program will receive their bachelor's degree from Coe College,
while also having a year of law school completed at the University of
Iowa College of Law.

There is no prescribed curriculum for students intending to enter law
school after graduation. Law schools report that their most successful
students are those who have acquired a broad academic background in the
liberal arts, developed a capacity for logical analysis, and mastered
the ability to write clearly.

Study in one or several of a variety of disciplines will prepare
students to undertake legal training. Students interested in law should
consult their department advisor and the pre-law program director.

\subsection{Health Professions}\label{health-professions}

Storer (Program Director).

Coe's program for those interested in the health professions, such as
medicine and dentistry, is a flexible one based on the requirements of
the health professions' schools. Students interested in a health
profession usually major in one of the sciences, but all majors offered
by the College are acceptable.

\subsection{Professional School Degree Completion
Plan}\label{professional-school-degree-completion-plan}

Students who complete three years of coursework at Coe (24 course
credits), including general education and area of study requirements,
and who enter a college of architecture, engineering, or a physical
therapy program, can receive a baccalaureate degree from Coe. Required
for satisfactory completion of this program are (1) approval of the
program by the Provost and Dean of the Faculty before transferring to
the professional institution and (2) one year of full-time acceptable
study there.

\subsection{Cooperative Degree Program With The University Of Iowa's
College Of Public
Health}\label{cooperative-degree-program-with-the-university-of-iowas-college-of-public-health}

This combined undergraduate and graduate 5-year program allows students
to earn a Bachelor of Arts degree from Coe College and a Master of
Public Health (MPH) degree from the University of Iowa. This program is
available to students electing any undergraduate major offered by Coe.
The first four years in the program are spent in residence at Coe
College. Interested students take one University of Iowa undergraduate
course in the spring of their second year (\emph{Fundamentals of Public
Health}). In their third year, students take the GRE and apply to the
Master's program in Public Health. If accepted, the student completes up
to four graduate-level MPH courses during the fourth year at Coe College
(\emph{Introduction to Biostatistics, Introduction to Health Promotion
and Disease Prevention, Global Environmental Health, and/or
Epidemiology}). UI Public Health courses are accepted in transfer to Coe
College as elective credit.

\chapter{OFF-CAMPUS STUDY}\label{off-campus-study}

The College endorses a wide variety of off-campus experiences for
students. Coe strongly believes that students can benefit from study in
Washington, New York, and the Wilderness Field Station, as well as in
programs around the world.

Numerous domestic and international study programs are available to Coe
students. The four programs sponsored by Coe College are Asia Term, New
York Term, Washington Term and Wilderness Field Station courses. In
addition, programs are offered by the Associated Colleges of the Midwest
(ACM) and others, including several exchange programs, by colleges,
universities, and educational agencies in America and abroad. Students
who wish to study off-campus on Coe's sponsored programs must apply to
the individual program's director. Student proposals to study on
exchange programs must be submitted to the Director of Off-Campus
Studies at least six weeks before the end of the term just prior to the
off-campus experience. Student proposals to study on all other
off-campus programs must be submitted to the Director of Off-Campus
Studies by the last day of classes of Fall Term of the academic year
prior to the program.

Credits earned on off-campus programs are applied toward graduation on
the same basis as credits earned on campus. Any academic credit earned
from programs not sponsored by Coe may be transferred back to the
College in accordance with the College's general policy on transfer
credit. Application is open to sophomores, juniors, and seniors who have
a minimum 2.25 GPA at both the time of application and of enrollment for
the off-campus study.

NOTE: \emph{International students may not receive Coe College financial
aid for off-campus study outside the U.S.}

\section{Domestic Programs}\label{domestic-programs}

\subsection{New York Term}\label{new-york-term}

Carson (Program Director)

The New York Term is open to all students who meet the basic
requirements for off-campus study. Offered in odd-numbered years in the
Spring Term, this program provides abundant opportunity for experiences
in the performing and visual arts, as well as internship opportunities
for students from any major.

The central course, Fine Arts in New York City (NYT-250), which includes
attendance at concerts, theatre, and dance productions as well as tours
to art exhibits and film screenings, consists of five 0.4 credit
courses: art, music, theatre, dance, and film. A faculty member for each
area grades the respective course. An internship (or other project
approved in advance by the College) completes the program.

There is an extra fee for New York Term. All Coe financial aid applies,
and students are eligible to apply for additional financial aid based on
the additional costs of the term. (See p.~72 for descriptions and course
offerings.)

\subsection{Washington Term}\label{washington-term}

B. Nesmith (Program Director).

The Washington Term is open to all students who meet the basic
requirements for off-campus study. Students accepted for Washington Term
spend Fall or Spring Term in the nation's capital. Washington provides
an unusual opportunity to study national politics and government and to
enjoy a variety of cultural activities in the fine arts. (See p.~75 for
descriptions and course offerings.)

\subsection{Wilderness Field Station}\label{wilderness-field-station}

Ellis (Program Director).

The Coe College Wilderness Field Station, located on remote Low Lake in
Minnesota's Superior National Forest, offers students a unique and
unparalleled opportunity for off-campus study. Courses take advantage of
the serene surroundings for field observation, wilderness study, and
outdoor learning. Biology courses are at the heart of the field station
and often include aquatic biology, animal behavior, ornithology, and
behavioral ecology of vertebrates. Students use the base camp's
laboratories, herbarium, and library to supplement their field work.
Non-science electives, such as nature writing and wilderness and the
law, are also offered. There is an independent study option as well.
Participants take one course during a four-week session.

The program runs from mid-June to mid-July, mid-July to mid-August, or
both. Each course is limited to eight students. Each course offered at
the field station is one course credit. Particular courses satisfy lab
science and other general education requirements and can be used as
major elective credits. Any course taken at the field station satisfies
the College's practicum requirement. The regular application deadline is
March 1st; the final deadline is April 15th, with rolling applications
after that date. For more information and application materials, visit
The Field Station Webpage:
www.coe.edu/academics/coe-difference-centers-and-programs/off-campus-study/wilderness-field-station.

\section{ACM \& International
Programs}\label{acm-international-programs}

Coe offers opportunities to study in locations around the world. Some
programs are for students wishing to broaden their liberal arts
perspectives, while others allow intensive research and study in a
specific academic area. Although some programs provide grades in the
courses, all grades transfer back to Coe as P/NP. For detailed
information and applications, students should contact directors for each
program or visit http://www.acm.edu/off\_campus\_study/index.html.

\subsection{OCC-350 --- Japan Study}\label{occ-350-japan-study}

Nordmann (Program Director).

Students study at Waseda University's School of International Liberal
Studies in Tokyo after a brief orientation providing intensive language
practice and cultural discussions. In addition to required language
study, electives may be chosen from a wide range of Asian studies
courses taught in English. A family living experience in Tokyo provides
an informal education in Japanese culture and is in many ways the
dominant feature of the program, offering total immersion in the
Japanese way of life. The program is recommended for a full year of
study, although students usually cannot receive Coe gift aid for more
than one term. The full-year program includes a month-long cultural
practicum or internship in another region of Japan, usually in February
or March. Administered by Earlham College, Japan Study is recognized by
both ACM and GLCA. Learn more at
http://www.acm.edu/programs/8/japan/index.html.

\subsection{OCC-360 --- Newberry Seminar: Research in the
Humanities}\label{occ-360-newberry-seminar-research-in-the-humanities}

Swenson Arnold (Program Director).

Students in the Newberry Seminar do advanced independent research in one
of the world's great research libraries. They join ACM and GLCA faculty
members in close reading and discussion centered on a common theme, and
then write a major paper on a topic of their choice, using the Newberry
Library's rich collections of primary documents. The Fall Seminar runs
for a full term; the Spring Seminars are month-long. Students live in
Chicago apartments and take advantage of the city's rich resources. The
Newberry Seminar is for students looking for an academic challenge, a
chance to do independent work, and possibly planning to attend graduate
school. Administered by ACM, the Newberry Seminar is also recognized by
GLCA. Learn more at http://www.acm.edu/programs/14/newberry/index.html.

\subsection{OCC-325 --- Field Museum Semester: Research in Natural
History}\label{occ-325-field-museum-semester-research-in-natural-history}

Hughes (Program Director).

An intensive research-and class-based experience for upper-level
students interested in natural history research with a background in
evolutionary biology, anthropology, or related discipline. The program
provides the opportunity for students to explore scientific research and
the Field Museum collections through a substantive internship, a course
taught by the visiting Faculty Director, and a seminar led by Field
Museum professional staff. See ACM website for more information
(https://acm.edu/students/off-campus-study-programs/field-museum-semester-research-in-natural-history/).

\subsection{OCC-365 --- Oak Ridge Science
Semester}\label{occ-365-oak-ridge-science-semester}

St.~Clair (Program Director).

The Oak Ridge Science Semester is designed to enable qualified
undergraduates to study and conduct research in a prestigious and
challenging scientific environment. As members of a research team
working at the frontiers of knowledge, participants engage in long-range
investigations using the facilities of the Oak Ridge National Laboratory
(ORNL) near Knoxville, Tennessee. The majority of a student's time is
spent in research with an advisor specializing in biology, engineering,
mathematics, or the physical or social sciences. Students also
participate in an interdisciplinary seminar designed to broaden their
exposure to developments in their major field and related disciplines.
In addition, each student chooses an elective from a variety of advanced
courses. The academic program is enriched in informal ways by guest
speakers, departmental colloquia, and the special interests and
expertise of the ORNL staff. Administered by Denison University, Oak
Ridge Science Semester is recognized by both ACM and GLCA. Learn more at
http://www.acm.edu/programs/15/oakridge/index.html.

\section{Asia Term}\label{asia-term}

Nordmann (Program Director).

The Asia Term is open to all students who meet the basic requirements
for off-campus study. Usually offered in the Spring Term, this program
provides students an opportunity to experience a variety of Asian
cultures in such countries as Thailand, Vietnam, and Cambodia. At each
site, students study language, read works in English about the culture,
engage in service learning, and work with students at the host
universities to gain a functional understanding of how each culture
works. Students are accompanied to Asia by Coe faculty members.

Students take four credits of coursework, typically one credit of Asian
Tonal Languages, one credit of Asian studies, and two credits of
independent study. In some iterations of the program, students take an
elective course in art, English, history, sociology, education, or
another discipline, depending on the field of the instructor leading the
program, and one credit of independent study. (See p.~70 for
descriptions and course offerings.)

\section{Exchange programs}\label{exchange-programs}

Coe College sponsors a number of programs with cooperating foreign
universities, offering Coe sophomores, juniors and seniors each year the
opportunity to study in a foreign setting. Coe College accepts in return
junior-level students from the foreign institution. Applications of the
recommended students are sent to the host institution, with the host
reserving the right to admit or reject each student nominated.

Any student who applies for one of these programs must have completed at
least one year of continuous study at their home institution. Students
may apply to any appropriate academic program offered at the host
institution as full-time, non-degree seeking, or unclassified students.
Any academic credit earned at the host institution is transferred back
to the home institution in accordance with the rules of that
institution. The length of stay may not exceed one academic year. Upon
completion of the time period specified at the host institution, the
participating students must return to their home institution. Any
extension of stay must be approved by both cooperating institutions. The
exchange student must abide by all rules and regulations of the host
institution.

An exchange student must register and pay tuition and required fees at
his or her home institution. In return, the student receives a tuition
waiver at the host institution. The host institution helps arrange the
necessary visa documents and also provides appropriate advising and
other assistance to the incoming students from Coe College. Please see
individual program descriptions for information regarding housing costs.
The host institution assists in finding housing on the foreign
university campus; Coe College assists in finding residence housing for
students from the foreign university. At the end of the school year, the
host institution submits to the home institution and official
transcripts of grades and credits earned. Grades from exchange programs
transfer to Coe as P/NP.

Beyond tuition and fees, the participating student is responsible for
the following expenses: meal expenses; transportation to and from the
host institution; medical insurance and/or medical expenses; textbooks,
clothing, and personal expenses; passport and visa costs; and all other
debts incurred during the course of the year.

\subsection{OCC-205 --- Coe/Kongju National University (South
Korea)}\label{occ-205-coekongju-national-university-south-korea}

Nordmann (Program Director).

Course offerings in Business, Economics, and Asian studies. Students pay
the cost of living at the destination. Credits earned from the Kongju
exchange program are evaluated on a P/NP basis.

\subsection{OCC-210 --- Coe/Chiang Mai University
(Thailand)}\label{occ-210-coechiang-mai-university-thailand}

Chaimov (Program Director).

Coe students usually pursue an independent research project based on
prior study in Thailand, as CMU offers no courses in English. Students
pay cost of living at destination.

Credits earned from the Chiang Mai exchange program are evaluated on a
P/NP basis.

\subsection{OCC-213 --- Coe/Rangsit University
(Thailand)}\label{occ-213-coerangsit-university-thailand}

Chaimov (Program Director).

Rangsit University. Coe students are responsible for securing their own
accommodations. Rangsit University offers English language bachelor's
degree programs in communications, international business, and
international political economy. Students pay cost of living at
destination.

Credits earned from the Rangsit University exchange program are
evaluated on a P/NP basis.

\subsection{OCC-215 --- Coe/Mid Sweden University
(Sweden)}\label{occ-215-coemid-sweden-university-sweden}

Carstens (Program Director).

Courses in English are available in such areas as business, social
sciences, and environmental studies. Students pay cost of living at the
destination. Credits earned from the Mid Sweden exchange program are
evaluated on a P/NP basis.

\subsection{OCC-220 --- Coe/Nagoya-Gakuin University
(Japan)}\label{occ-220-coenagoya-gakuin-university-japan}

Nordmann (Program Director).

One year of Japanese language study is recommended for students applying
for this program. Exchange students from Coe pay room expenses at Coe
College. In return, students receive a room expense waiver. Other costs
of living are paid at Nagoya-Gakuin University. Credits earned from the
Nagoya-Gakuin exchange program are evaluated on a P/NP basis.

\subsection{OCC-225 --- Coe/Northern Ireland Scholars Program (Northern
Ireland,
UK)}\label{occ-225-coenorthern-ireland-scholars-program-northern-ireland-uk}

Farrell (Program Director).

Students with a high GPA may be selected to study at one of several
universities in Northern Ireland, including Queens University Belfast
and the University of Ulster. Applications for this consortial exchange
are due in December of the year before study. Students pay cost of
living at destination.

Credits earned from the Northern Ireland exchange program are evaluated
on a P/NP basis.

\subsection{OCC-230 --- Coe/University of Jaume I (Castello,
Spain)}\label{occ-230-coeuniversity-of-jaume-i-castello-spain}

Rodríguez Moreno (Program Director).

Courses in Spanish in a wide range of topics. Students pay cost of
living at destination. UJI requires Coe students to have completed two
Spanish courses at Coe. Spanish language courses are available for an
extra cost.

Credits earned from the Jaume I exchange program are evaluated on a P/NP
basis.

\subsection{OCC-235 --- Coe/Sookmyung University (South
Korea)}\label{occ-235-coesookmyung-university-south-korea}

Nordmann (Program Director).

Courses offered in English on areas including the arts, linguistics,
international studies, business, biology, and political science.
Students pay cost of living at destination. Credits earned from the
Sookmyung exchange program are evaluated on a P/NP basis.

\subsection{OCC-240 --- Coe/University of Landau
(Germany)}\label{occ-240-coeuniversity-of-landau-germany}

Chaimov (Program Director).

Courses in English are available in, art, English literature,
linguistics, other topics. Also, a wide range of subject areas taught in
German. Students pay cost of living at the destination.

Credits earned from the Landau exchange program are evaluated on a P/NP
basis.

\subsection{OCC-245 --- Coe/University of Quebec (Saguenay,
Canada)}\label{occ-245-coeuniversity-of-quebec-saguenay-canada}

Janca-Aji (Program Director).

Courses in French are offered in a wide range of topics for students who
pass a proficiency test in French. Students pay cost of living at
destination. Credits earned from the Quebec exchange program are
evaluated on a P/NP basis.

\subsection{OCC-250 --- Coe/Izmir Institute of Technology
(Turkey)}\label{occ-250-coeizmir-institute-of-technology-turkey}

Akgun (Program Director).

Coe may send a student of Chemistry and a student of Physics to study
those subjects in an English-language setting at a science university in
Turkey. Students are responsible for housing, food, and all other costs
of living.

Credits earned on the Izmir exchange program are evaluated on a P/NP
basis.

\subsection{OCC-255 --- Coe/Polytechnic University of Upper France
(France)}\label{occ-255-coepolytechnic-university-of-upper-france-france}

Janca-Aji (Program Director).

Students choose from courses in English on business, communications, and
marketing or a wide range of courses in French. Students pay costs of
living at the destination. Credits earned from the France exchange
program are evaluated on a P/NP basis.

\subsection{OCC-260 --- Coe/National University of Villa Maria
(Argentina)}\label{occ-260-coenational-university-of-villa-maria-argentina}

Students with a good command of Spanish can take courses in Spanish in a
wide range of areas, including literature, rural development, social
sciences, environmental studies, communication, and computer science.
Students pay costs of living to Coe before departure and must transfer
at the same time as an incoming student from UNVM.

Credits earned on the Argentina exchange program are evaluated on a P/NP
basis.

\subsection{OCC-265 --- Coe/Istanbul Altinbas University
(Turkey)}\label{occ-265-coeistanbul-altinbas-university-turkey}

Duru (Program Advisor)

Teaches entire majors in English in psychology, sociology, international
relations, political science, economics, business. Students pay cost of
living at the destination. Credits earned from the Altinbas exchange
program are evaluated on a P/NP basis.

\subsection{OCC-270 --- Coe/Ashesi University
(Ghana)}\label{occ-270-coeashesi-university-ghana}

Eichhorn (Program Director).

Courses in African studies (sociology, anthropology, political science,
history, arts), computer science, business. Coe students pay room and
board expenses at Coe College and receive a waiver of these expenses in
Ghana. They must exchange at the same time as an Ashesi student.

Credits earned on the Ashesi exchange program are evaluated on a P/NP
basis.

\subsection{OCC-275 --- Coe/University of Salford
(England)}\label{occ-275-coeuniversity-of-salford-england}

Kuennen (Program Director).

The University of Salford offers courses in business, contemporary
European history, psychology, and many other areas. Students pay cost of
living at the destination and must arrange their own housing.

Credits earned from the Salford exchange program are evaluated on a P/NP
basis.

\subsection{OCC-285 --- Coe/University of Neuchatel
(Switzerland)}\label{occ-285-coeuniversity-of-neuchatel-switzerland}

Janca-Aji (Program Director).

Coe students take courses at the Institute of French Language and
Civilization. Students pay costs of living at the destination. Credits
earned from the Neuchatel exchange program are evaluated on a P/NP
basis.

\subsection{OCC-290 --- Coe/Jinan University
(China)}\label{occ-290-coejinan-university-china}

Nordmann (Program Director).

Located in southern China, this international university offers courses
in English in international economics and business, journalism, computer
science, and Chinese studies as well as the study of Chinese language.
Students pay costs of living at the destination.

Credits earned from the Jinan exchange program are evaluated on a P/NP
basis.

\subsection{OCC-291 --- Coe/Istanbul Kultur University
(Turkey)}\label{occ-291-coeistanbul-kultur-university-turkey}

Duru (Program Advisor).

Offers English curriculum in psychology, business, economics,
international relations, and other areas. Students pay costs of living
at the destination. Credits earned from the Istanbul Kultur University
exchange program are evaluated on a P/NP basis.

\part{CAMPUS RESOURCES}

\chapter{CAMPUS RESOURCES}\label{campus-resources-1}

\subsection{Libraries}\label{libraries}

The College libraries---Stewart Memorial Library, located at the center
of the campus and Fisher Music Library in Marquis Hall---contain over
500,000 volumes and 16,000 pieces of media. Current subscriptions to
some 3,500 periodicals and serials are maintained in print or electronic
format, and over 200,000 electronic resources with books and journal
volumes added annually.

The collections ably support undergraduate education and are especially
strong in the areas of literature, history, and music. The Fisher Music
Library contains over 5,000 compact discs and records, 5,300 scores and
books, and is equipped with listening facilities. Media services to the
campus are provided through the Media Technologies Department in the
library. These services include a circulating collection of over 8,000
DVDs, two media-equipped auditorium styled classrooms, editing stations,
an innovation studio that houses a 3D printer and laser cutter, and a
variety of cameras, recording equipment available for use.

The main library houses the Learning Commons (see description below) and
the college archives. The library provides an outstanding research
collection consisting of both print and electronic books, journals, and
reference resources. The Reference Department assists students with
their research needs through one on one research assistance, evaluating
resources and websites, citation assistance and multimedia evaluation.
In addition to library orientations, research classes are offered on
specific course related topics. The library offers computer stations,
iPad and laptop checkouts, study areas for individual and group study,
including technology enhanced study rooms. These resources are greatly
augmented by providing access to over 100 scholarly databases and an
extensive webpage: coe.edu/library.

The George T. Henry College Archives includes a research room and a
climate-controlled vault located on the lower level of the Stewart
Memorial Library. It houses and preserves the institutional records of
Coe College and the papers of staff, students, and alumni of the
college. Archive staff is available to aid students, faculty, and
scholars in navigation and use of more than 900 linear feet of primary
source documents. The Archives maintains the papers of journalist,
author, and World War II broadcaster William L. Shirer, Coe class of
1925 author of \emph{The Rise and Fall of the Third Reich}; the literary
works and selected private papers of Iowa poet Paul Engle, Coe class of
1931; and the photographs of longtime Coe College photographer George T.
Henry.

\subsection{Learning Commons}\label{learning-commons}

The Coe Learning Commons in the Stewart Memorial Library integrates all
of the College's academic support resources in a single location at the
heart of campus delivered through peer education and by professional
staff. Services and resources include academic coaching, supplemental
advising, Writing Center, AAP-TRIO program, tutoring, accessibility
support and accommodations, academic technology, Office of Off-Campus
Study, and fellowship and graduate school advising.

\subsection{Art Collections}\label{art-collections}

Selections from the College's Permanent Collection of Art totaling over
800 works by 200 artists are displayed in and near many of the campus
buildings. Most visible are the large outdoor sculptures on the campus,
yet almost every building features selections from the Permanent
Collection.~ For example, the Ella Poe Burling collection of
nineteenth-century American and French art and antiques is exhibited in
the upper lobby of Voorhees Hall.

A large portion of the Permanent Collection can be found in Stewart
Memorial Library.~ Four special galleries contain works by renowned
American painters Grant Wood, Marvin Cone (Coe class of 1914), and
Conger Metcalf (Coe class of 1936). Six large farm murals by Grant Wood
constitute the heart of the Permanent Collection's Regionalist works.
These murals are supplemented by nine smaller yet significant works by
Wood, including \emph{Daughters of Revolution}, a charcoal, pastel, and
pencil on paper drawing of Wood's painting of the same name.

Another signature feature of the Regionalist collection is the work of
Marvin Cone. A 1914 Coe graduate, Cone later became a faculty member who
founded the College's Art Department. Cone personally selected many of
the paintings and drawings in the collection as representative of his
own artistic development, underscoring the historic role of art as a
core element in Coe's teaching mission.

In addition to the works by native Iowans Cone and Wood, Coe College has
acquired a distinguished collection of 70 works by Conger Metcalf, an
American modernist painter, as well as paintings, drawings, and prints
by notable artists such as Milton Avery, Mauricio Lasansky, Henri
Matisse, Pablo Picasso, and Andy Warhol.

More information about the Permanent Collection can be found on either
the Permanent Collection's website,
http://picovado.com/jrogers/\#-h2-introduction-h2-, or the library's
webpage, www.coe.edu/academics/stewart-memorial-library.

\subsection{Information Technology}\label{information-technology}

The Information Technology Office~provides a wide range of technology
support to students,~faculty and staff. This includes management of
computer labs, classroom technology, college-wide software licensing,
wi-fi/internet, printers, cable TV, sound equipment, My Coe (my.coe.edu)
and more. Coe's technology facilities include over 3,000 ethernet ports,
full ethernet and wi-fi internet services within all campus buildings,
250 public/lab computers available~for student use and full access to
G-Suite services. The Office of Information Technology is located in
Voorhees Hall. Assistance from the IT staff can be requested through a
Help Desk/Spiceworks ticket
(https://coecollege.on.spiceworks.com/portal/).~

\subsection{Public Events and
Artists-in-Residence}\label{public-events-and-artists-in-residence}

Guest lecturers and artists provide an essential dynamism to the
educational climate at Coe. Programs are free to Coe students. In
addition to hearing speakers of national and international note,
students may have the opportunity to talk with them during a carry tray
lunch, to attend special issues dinners with the speakers, or to have
them as guests in a class. Performing groups appearing as
artists-in-residence often stay two or three days on campus to work with
students.

\subsubsection{Marquis Lecture \& Performance
Series}\label{marquis-lecture-performance-series}

The Marquis Lecture \& Performance Series hosts performances and
presentations throughout the academic year. The Marquis Series is
endowed by a gift from Sarah Marquis, Coe class of 1918, in honor of her
father, John A. Marquis, who was president of Coe from 1909 to 1919.

\subsubsection{Coe College Contemporary Issues
Forum}\label{coe-college-contemporary-issues-forum}

The Coe College Contemporary Issues Forum brings to audiences of the
College community the presence and

views of distinguished professionals whose work has received national
recognition. The forum is normally presented during the month of
February.

\subsubsection{Phi Beta Kappa Visiting
Scholar}\label{phi-beta-kappa-visiting-scholar}

The Coe chapter of Phi Beta Kappa sponsors a Phi Beta Kappa Visiting
Scholar who, in addition to presenting a public lecture, meets with
individual classes as appropriate to the scholar's area of expertise.

\chapter{STUDENT LIFE}\label{student-life}

Student Life provides personalized support to students, enabling each to
gain the best possible undergraduate education adding substantially to
the educational program. Residence accommodations, along with living and
learning values, an attractive campus social life, a sound health
program, good recreational facilities, and a program of co-curricular
activities are among the opportunities offered Student Life.

\subsection{Campus Civility Statement}\label{campus-civility-statement}

This statement was written by students in order to address standards of
civility and respect within the Coe College community. This statement is
a living document and is intended to evolve over time.

\begin{itemize}
\tightlist
\item
  We, the members of the Coe College community, expect our campus
  climate to be safe, mutually supportive, academically encouraging,
  egalitarian, and tolerant of all its members.
\item
  We expect the academic experience to extend beyond the classroom into
  our living environment.
\item
  We expect a campus free of incidents that create a hostile living
  environment.
\item
  We expect a healthy and responsible attitude to accompany all social
  gatherings.
\item
  We expect that intoxication will not be an excuse for incidents that
  occur while under the influence.
\item
  We expect that diversity of opinion should be cultivated and
  encouraged as well as respected within our community.
\item
  We expect that everyone will have the right to be respected for his or
  her individuality.
\item
  We expect all campus community members to respect the rights of other
  persons regardless of their actual or perceived age, color, creed,
  disability, gender identity, national origin, race, religion, sex, or
  sexual orientation.
\end{itemize}

A community is made up of individuals who model these standards and hold
each other accountable. In order for the community to encompass the
goals outlined above, each individual must be responsible and
accountable for her or his own actions and words.

\subsection{Student Contribution to College
Policy}\label{student-contribution-to-college-policy}

Coe is proud of its traditions and its ability to change. While
cognizant of the past, the College is also sensitive to the changing
nature and needs of students. Coe students play an integral part in the
initiation of change by utilizing available channels and by the creative
development and use of new ones. The Student Senate is a frequent forum
for the resolution of student concerns.

\subsection{Student Senate}\label{student-senate}

Student Senate is the representative government of Coe students and
coordinates many co-curricular activities. Through the student activity
fee, the Senate funds the weekly newspaper, the Cosmos, and other
student publications. The Student Activities Committee (SAC) of the
Student Senate sponsors bands, comedians, multicultural programming, and
other kinds of entertainment and activities.

\subsection{Student Handbook}\label{student-handbook}

The \emph{College Policies and Student Handbook} outlines the College's
expectations for responsible behavior reflecting maturity, mutual
respect, and cooperation among all members of the Coe community. Student
Life develops policies for conduct procedures, residence hall living,
student organizations, and other areas of student affairs for approval
by the Board of Trustees. The College Policies and Student Handbook is
available online at
https://www.coe.edu/student-life/student-development/college-policies-student-handbook.

\subsection{Committee Participation}\label{committee-participation}

Students serve on various committees, which aid in making educational
policy at the college. Most committees (Academic Policies, Assessment,
Athletics, Campus Technology, Diversity, Enrollment, Financial Aid and
Academic Progress, Executive, Finance and Facilities, First-Year
Program, Internationalization, Marquis Series, Petitions, Sustainability
Council, Wellness, and Writing) include students appointed by the
Student Senate as voting members.

\section{Student Services}\label{student-services}

\subsection{Residence Life}\label{residence-life}

Coe is a residential college, in that the residence experience is an
integral part of the educational process. Students are expected to live
on campus for four years and take meals in the College hall (see
On-Campus and Off-Campus Resident Students, Section 27.1.4). The
residence halls and apartments vary in style, size, and personality. All
of the residential facilities have generous visitation policies, and
campus life functions around the concept of the living units.

Residence hall and apartment regulations are published in the
\emph{Student Reference Book}. As room charges do not include Winter
Break or Spring Break, when residential facilities and the dining hall
are closed, an additional fee is assessed to students staying on campus
during those times. Information regarding housing is sent to students
who have accepted admission to the College.

\subsection{Student Health Service}\label{student-health-service}

Coe's Student Health Service offers students medical care provided by a
part-time Advanced Registered Nurse Practitioner (ARNP) and a full-time
nurse. The Health Services staff can diagnose, manage, and treat certain
medical diagnoses, free to full-time Coe students. Referrals are
available to a family physician or specialist in the Cedar Rapids
community as needed. In the event that a student needs hospitalization,
St.~Luke's Hospital or Mercy Medical Center is utilized. The student is
responsible for all health costs beyond those provided by Health
Services. These include hospitalizations, emergency room visits,
physician visits, and prescription medications. Therefore, all students
are expected to carry medical insurance. Provisions for special diets or
other arrangements which deviate from Coe's policies require a
recommendation from the student's healthcare provider and are available
through SODEXO food services. All students are required to have a
current immunization record on file prior to registration for classes at
Coe. Students without immunizations on file will have their registration
held. Students who wish to utilize the Student Health Service should
have a physical on file, which is required of all students in order to
play collegiate sports.

\subsection{Personal Counseling}\label{personal-counseling}

Realizing that students have concerns in areas other than academic
matters and career options, Coe provides appropriate individual and
group counseling. St.~Luke's Family Counseling Center, located next to
the Coe campus, and the College work together to provide for students'
counseling needs. St.~Luke's Family Counseling Center provides
assessment, short-term counseling, and, when appropriate, referral to
community resources. Individual counseling is available from a number of
counselors on an appointment basis for students with personal, social,
and family concerns. In addition, students may seek counseling from a
pastoral and spiritual perspective from the College Chaplain. If
long-term counseling is needed, Coe's counselor and Chaplain work with
students to identify cost-effective solutions on and off campus.

\subsection{Religious Life}\label{religious-life}

Coe College believes that it is important to foster an environment that
accepts and respects the religious faith and beliefs of all its
students, staff, and faculty. The Coe community is diverse in its
religious makeup, and all members of the community are encouraged to
express and practice their particular religious traditions. To this end,
the College Chaplain seeks to provide and create an atmosphere that is
consistent with the religious heritage of the College and conducive to
the development of spiritual and moral values.

Opportunities for worship, interfaith dialogue, Bible study, retreat,
small groups, theological study, outreach, mission, meditation, and
fellowship are abundant. There are also several active religious student
organizations on campus and a specialized leadership program for
students considering vocational ministry. The Chaplain is available for
pastoral care, guided prayer, theological dialogue, pre-marital
counseling, and other spiritual needs.

\subsection{Diversity, Equity \&
Inclusion}\label{diversity-equity-inclusion}

The Office of Diversity, Equity \& Inclusion is comprised of the Dean of
Students, Coordinator of Multicultural Affairs, Director of
International Affairs, and the College Chaplain. This team works closely
with LBGTQAI+ students, multicultural students, international students,
and student allies, with the goal of fostering an environment where all
Kohawks will thrive in an inclusive learning environment. Students
interested in getting involved with Coe's diversity and inclusion
efforts are encouraged to stop by the Student Life Office in Upper Gage
and speak with a team member.

\subsection{Campus Activities}\label{campus-activities}

There is much to do on the Coe campus and in the Cedar Rapids community.
Programming of campus activities is designed to meet the educational and
recreational needs of the Coe community in a creative way. The diversity
of the student body is considered in the scheduling of recitals, plays,
exhibits, lectures, films, and concerts, as well as all school events
and relaxing evenings in Charlie's. The Director of Campus Life
coordinates the events organized by the Student Activities Committee
(SAC).

\subsection{Student Activity Groups}\label{student-activity-groups}

Students earn credit for their participation in Coe's music ensembles
(the Jazz and Concert Bands, the Symphony Orchestra, the Concert Choir,
Chorale, and the Choral Chamber Ensemble), several of which have
completed study/concert tours of Europe, Great Britain, and Asia in the
past. Coe drama and forensics activities have received local and
national recognition for their presentations. Special interest
organizations are also represented on campus, as are national social
fraternities and sororities and honor societies (Phi Beta Kappa, Phi
Kappa Phi, Mortar Board, Alpha Lambda Delta, and Alpha Sigma Lambda).

\subsection{Athletics}\label{athletics}

Coe College sponsors 11 men's and 11 women's athletic teams that compete
in the American Rivers Conference of NCAA Division III. Our goal is to
provide our student athletes with positive educational and athletic
experiences.

Through hard work, intense training, and positive interactions with
coaches, student athletes are provided opportunities to succeed. The
College recognizes that many of its students enjoy participating in
organized athletics or watching athletic contests as forms of
recreational campus life. Basketball, volleyball, softball, table
tennis, flag football, and wrestling are representative events in a
year-round intramural program for both men and women. All students are
eligible to participate.

\subsection{Recreational Facilities}\label{recreational-facilities}

Gage Memorial Union is the center of student interest and activity.
Offices of the Student Activities Committee and other student
organizations are there, plus the College dining hall, and mailroom.
Informal programs and lectures are given there, and it serves as the
College's ``open house'' for students, faculty, and visitors.
``Charlie's,'' located in the adjacent P.U.B., is home to a coffee shop
(serving Starbucks coffee), a convenience store, and grill. Charlie's is
a relaxing place to meet friends or take in one of the many
performances.

The brand-new Coe College Athletics and Recreation Center includes two
pristine courts for basketball and volleyball, a wrestling room that
boasts three oversized mats, strength and conditioning room with
brand-new equipment, and a fitness center that overlooks the campus.
This is where Coe's basketball, volleyball, and wrestling teams host
their home events.

Moray Eby Fieldhouse includes three recently-renovated courts for
basketball and volleyball. Eby also has a natatorium, athletic training
rooms, indoor baseball/softball batting cages, and a rock-climbing wall.
All outside playing fields and tennis courts are also available to Coe
students.

The Clark Racquet Center offers a state-of-the-art facility for runners,
racquet enthusiasts, and everyone wanting to stay in shape. An aerobic
room, dance studio, and classrooms complement the indoor track, tennis
courts, and racquetball and squash courts. Professional staff manage the
center and offer instruction to students.

\part{COLLEGE REGULATIONS}

\chapter{EFFECTIVE CATALOG}\label{effective-catalog}

Students ordinarily are graduated under the provisions of the catalog of
their matriculation date. A student's matriculation date is the day of
first enrollment following admission. If the student is readmitted, the
matriculation date becomes the day of first enrollment following
readmission. However, students will be expected to satisfy, to the
extent practical, the graduation requirements of the catalog in effect
at the anticipated date of graduation. Any necessary modification of
general degree requirements will be worked out by the Provost and Dean
of the Faculty, the Registrar, the appropriate department chair, and the
Academic Policies Committee. Modification of major or minor requirements
will be worked out by the department chair involved, the Registrar, and
the student's advisor. A student has the right of petition to the
Committee on Petitions.

\subsection{Grading}\label{grading}

\textbf{GRADES}

\begin{longtable}[]{@{}
  >{\raggedright\arraybackslash}p{(\columnwidth - 4\tabcolsep) * \real{0.0533}}
  >{\raggedright\arraybackslash}p{(\columnwidth - 4\tabcolsep) * \real{0.4667}}
  >{\raggedright\arraybackslash}p{(\columnwidth - 4\tabcolsep) * \real{0.4800}}@{}}
\toprule\noalign{}
\endhead
\bottomrule\noalign{}
\endlastfoot
A & Excellent & 4.0 grade points per course credit \\
A- & & 3.7 grade points per course credit \\
B+ & & 3.3 grade points per course credit \\
B & Above average & 3.0 grade points per course credit \\
B- & & 2.7 grade points per course credit \\
C+ & & 2.3 grade points per course credit \\
C & Satisfactory, minimum expectation credit & 2.0 grade points per
course credit \\
& & \\
C- & & 1.7 grade points per course credit \\
D+ & & 1.3 grade points per course credit \\
D & Passing, below expectation & 1.0 grade points per course credit \\
D- & & 0.7 grade points per course credit \\
\end{longtable}

\subsection{Status Marks}\label{status-marks}

\begin{itemize}
\tightlist
\item
  \textbf{W} Approved withdrawal from a course.
\item
  \textbf{X} Course extends beyond term. An ``X'' status grade is given
  in courses designated in the Catalog as having coursework that extends
  beyond the end of the term. Under no circumstances can an ``X'' grade
  remain unresolved for more than one calendar year.
\item
  \textbf{O} No mark had been reported by the instructor by the time
  academic reports were processed.
\item
  \textbf{I} Incomplete. An ``I'' status grade is reported only for
  students who are unable to complete the work in the course due to
  extenuating circumstances. The normal length of time for resolution of
  an incomplete is within four weeks of the next Fall or Spring term in
  which the student enrolls. If the instructor believes the resolution
  of the incomplete will take longer, the instructor will note the later
  deadline when reporting the incomplete status grade. An unresolved
  incomplete will automatically become a failing grade after the
  deadline, unless the instructor notes otherwise when reporting the
  incomplete. Under no circumstances can an ``I'' grade remain
  unresolved for more than one calendar year.
\item
  \textbf{R} The prefix ``R'' to a grade (i.e., RA, RB, RC, RD, RF)
  indicates a grade of repeated course. A student may repeat a course
  previously taken, and registration must indicate this repeat. Failure
  to register for a repeat course properly results in no recognition of
  the second attempt. Only the grade earned when the course is retaken
  is used in computing the GPA. Credit may be earned only once for a
  given course. Courses may be repeated only once, although, students
  may petition for a second repeat if extenuating circumstances exist.
\item
  \textbf{EQ} Equivalent credit recognized; no credit given. Does not
  increase courses attempted. For a student who has completed four
  course credits of student teaching and who has high school or life
  experience equivalent to a regular catalog course, the said course,
  upon recommendation of the Education chair and the chair of the
  department in which equivalent credit is being recognized, may be
  listed on the student's transcript with the appropriate credit and a
  status mark of ``EQ.''
\end{itemize}

\subsection{Audited Courses}\label{audited-courses}

Students may audit courses with the consent of the instructor. In doing
so, they attend class but are not required to take tests or submit
papers. Audited courses receive no credit and do not appear on
transcripts. Auditors pay a reduced tuition charge.

\subsection{First Course Grading For Non-Traditional
Students}\label{first-course-grading-for-non-traditional-students}

Individuals who have been away from the collegiate routine for several
years may register on an audit basis in their first course, participate
fully in the class (including tests, papers, etc.), and decide at the
conclusion of the course if they wish to pay the other half of the
tuition and receive credit for the course. This policy applies only to
the first course---not to succeeding courses.

\subsection{Pass / Not-Pass Graded
Courses}\label{pass-not-pass-graded-courses}

Courses are graded A--F unless otherwise indicated in the course
description that only P/NP is an option (e.g., student teaching,
internships). Courses graded P/NP count as credits towards graduation
but do not affect the student's GPA. In addition, students may opt to
take up to four A--F graded courses on a P/NP basis. When students opt
for P/NP grading, the Office of the Registrar converts any grade a
faculty member provides of D- or better to a P, and any grade of F to an
NP. In order to count for Writing Emphasis credit, a submitted grade of
C or higher must be earned.

The following regulations apply:

\begin{enumerate}
\def\labelenumi{\arabic{enumi}.}
\tightlist
\item
  Students are permitted to change the method of grading for a course
  from a letter grade to P/NP. This change is allowed from the start of
  the term of enrollment through the last day to withdraw from courses
  during that term. See Academic Calendar Section 3.0.1 for official
  dates.
\item
  Students may elect to take up to four graded courses during their
  undergraduate career on a P/NP basis. However:

  \begin{enumerate}
  \def\labelenumii{\alph{enumii}.}
  \tightlist
  \item
    Students cannot use a course for which they elected P/NP grading to
    satisfy the requirements for a major or minor.
  \item
    Students cannot elect P/NP grading while on academic probation,
    though provisionally-admitted First Year students are allowed the
    option of P/NP grading.
  \end{enumerate}
\item
  The Registrar will not inform the instructor of the student's request
  for P/NP grading and the instructor must provide the Registrar with an
  appropriate letter grade.
\item
  A student's request for P/NP grading may be nullified at any time
  before the end of the third business day after grades are due in the
  student's final term. A written notice must be sent to the Office of
  the Registrar to communicate the student's intent to replace the P or
  NP grade with the instructor's letter grade.
\item
  A student's request to change their method of grading to P/NP in lieu
  of a letter grade counts as one of the four P/NP elected courses
  during the student's undergraduate career, regardless of whether it is
  later nullified.
\item
  Ordinarily a student is limited to one course credit per term on an
  elected P/NP basis. A student wishing to exceed this limitation must
  present a convincing rationale or significant mitigating circumstances
  to the Committee on Petitions.
\end{enumerate}

\subsection{Midterm Grades}\label{midterm-grades}

Midterm grades are not official evaluations and are not part of the
official transcript. The goal of midterm grades is to help students take
responsibility for their academic progress. Faculty submit midterm
grades of A-F for students in all full-term courses. Midterm grades are
not intended to be a guarantee, promise, or contract regarding the final
grade a student will earn in the class. Instead they provide information
for students about their academic performance.

\chapter{CLASS DESIGNATION}\label{sec-class-designation}

Class Designation is determined by the number of course credits earned
following the Fall and Spring Terms.

\begin{itemize}
\tightlist
\item
  First-year student: Fewer than 8.0 course credits earned\\
\item
  Sophomore: 8.0--15.9 course credits earned
\item
  Junior: 16.0--23.9 course credits earned
\item
  Senior: 24.0 or more course credits earned
\end{itemize}

\chapter{REGISTRATION}\label{registration}

Before being allowed to register, students must have been admitted for
study by the Office of Admission, must have settled their account with
the Business Office, and must meet with their Academic Advisor.
Registrations are canceled for failure to pay fees on time.

Students are encouraged to develop a four-year comprehensive educational
plan with the help of their advisor or with other members of the
faculty.

Students who will not graduate during the current academic year register
online during the Spring Term for the Fall Term of the following
academic year during times specified by the Office of the Registrar.
Registration for the Spring and May Terms takes place during the
preceding Fall Term. Once the online registration period is over, all
changes to registration must take place in the Office of the Registrar
using accurate, legible, and completed registration forms. Entering
first-year students and transfer students receive instructions with
their orientation materials, and readmitted students receive
instructions from the Office of the Registrar concerning their
registrations. Registrations are not accepted for a term without
successful petition after the end of the first week of classes of that
term.

\subsection{Changes in Registration}\label{changes-in-registration}

\begin{enumerate}
\def\labelenumi{\arabic{enumi}.}
\tightlist
\item
  Unless a course is registered by a student online during the
  registration period, they must intentionally file a registration form
  with the Office of the Registrar in order to be registered for a
  course. Registration forms for course additions must be approved by
  the student's advisor and instructor of the course.
\item
  Courses may be added or dropped during the first five weeks of the
  Fall, Spring, or Summer Term, including 7-week courses. Students may
  add a May Term course during the first three days of the term. If a
  student needs to add a course after these deadlines due to extenuating
  circumstances, the student must petition the Committee on Petitions by
  completing the appropriate forms.
\end{enumerate}

For courses spanning a fraction of a Term, the last date to add or drop
without a ``W'' will be calculated as a proportionate time frame
comparatively as full-term courses. See Academic Calendar (see p.~12)
for official dates for full-term and half-term courses. 3. A student may
withdraw from one or more courses with the following results: - If a
student withdraws from a course when 2/3 or less of the Term is
completed, a ``W'' (withdrawal) grade will be entered on the student's
permanent record. This grade will not affect the student's GPA. This is
the date listed in the academic calendar as the Last Day to Withdraw
from a course.\\
- For courses spanning a fraction of a Term, the last date to withdraw
with a ``W'' will be calculated using the 2/3 fraction. See Academic
Calendar (p.~15) for official dates for full-term and half-term courses.
- If a student officially withdraws from a course after 2/3 of the Term
is completed, a ``WF'' will be entered on the student's permanent
record. This will affect the student's GPA (see p.~53). - A withdrawn
course, regardless of the date withdrawn, counts as attempted credits on
the transcript. Therefore, it also counts when calculating academic
standing (see p.~61), and satisfactory academic progress (see p.~275)
for financial aid.

\subsection{Course Load}\label{course-load}

A student is full-time for the Fall or Spring Term when enrolled for
three or more course credits. A student enrolled for less than this
course load is a part-time student. For financial aid purposes,
half-time is defined as enrollment in no fewer than two course credits
in each of the Fall and Spring Terms; three-fourths is defined as
enrollment in no fewer than 2.5 course credits in each of the Fall and
Spring Terms. Students who have earned a cumulative GPA of 3.4 or above,
or have earned both 23 credits and a cumulative GPA of 2.5 or above, may
take above 4.99 credits without the need for a petition, but approval
must be secured from the Office of the Registrar. Other students must
submit a petition to the Committee on Petitions and obtain approval in
order to register for five or more course credits. No student shall be
permitted to register for over 5.99 credits in each of the Fall and
Spring Terms. Only petitions from those students who have completed at
least one term as a full-time student will be considered by the
Committee on Petitions. Students may not register for more than one
course credit during May Term.

\chapter{ACADEMIC INTEGRITY POLICY}\label{academic-integrity-policy}

At Coe College, we expect academic integrity of all members of our
community. Academic integrity assumes honesty about the nature of one's
work in all situations. Such honesty is at the heart of the educational
enterprise and is a precondition for intellectual growth. Academic
dishonesty is the willful attempt to misrepresent one's work, cheat,
plagiarize, or impede other students' academic progress. Academic
dishonesty interferes with the mission of the College and will be
treated with the utmost seriousness as a violation of community
standards.

\subsection{Forms of Academic
Dishonesty}\label{forms-of-academic-dishonesty}

\textbf{Cheating} is the attempt to deceive an evaluator by claiming
credit for work one has not done or by knowingly assisting such an
attempt. It includes (but is not limited to) the use of unauthorized
sources of information on in-class or take-home exams, or other
assignments; copying from other students on exams, assignments, or lab
reports; fabrication of data, research, quotations, or other
information; and taking credit for collaborations to which one has not
contributed.

\textbf{Plagiarism} is the use of someone else's words or ideas without
acknowledgement and, when intentional, is a form of academic dishonesty.
The unacknowledged use of words or ideas from any published or
unpublished sources, including Internet resources or other student
papers, constitutes plagiarism. Plagiarism may occur intentionally or
unintentionally through the omission of appropriate citations. Any ideas
or information the student adopts from a source, whether or not directly
quoted, must be acknowledged by specific reference in notes or the text.

\textbf{Any words or phrases} that are taken from a source must be
quoted and cited. Any paraphrase---the restatement of an idea in your
own words---must be cited.

The methods of citation and documentation vary from discipline to
discipline. Students are responsible for determining the appropriate
method for any given assignment or, in the absence of a clearly stated
protocol, using any accepted academic method. Guidelines can be found on
the library website and in the Writing Center.

\textbf{Other forms of academic dishonesty} include (but are not limited
to) deliberately impeding other students' work and misuse of common
academic property, in the libraries, labs, and elsewhere.

\subsection{Sanctions}\label{sanctions}

Instructors have responsibility for determining whether academic
dishonesty has occurred. Instructors shall proceed with sanctions
accordingly. Any act of academic dishonesty that results in one of the
sanctions below shall be detailed in a formal report filed with the
Provost and Dean of the Faculty.

Cases of unintentional plagiarism may be dealt with through educational
procedures such as further assignments requiring the student to practice
documentation and citation methods, or other means determined by the
instructor.

Acts of academic dishonesty will be subject to one or more of the
following sanctions:

\begin{enumerate}
\def\labelenumi{\arabic{enumi}.}
\tightlist
\item
  failure of the assignment, i.e.~exam, paper, lab report, etc.
\item
  failure of the class
\item
  suspension or expulsion
\end{enumerate}

An instructor may impose the first two of these penalties. Suspension or
expulsion may only be carried out by the Provost and Dean of the
Faculty.

Repeated acts of academic dishonesty will result in suspension or
expulsion. When academic dishonesty has been determined to have occurred
a second time, the Provost and Dean of the Faculty shall decide on the
student's status at the College.

\subsection{Procedure}\label{procedure}

When an instance of academic dishonesty is suspected, the instructor
will meet with the student to discuss the incident and will decide
which, if any, of the above sanctions is appropriate.

If warranted, the instructor will send a report to the Provost, with a
copy given to the student, which details the nature of the violation and
the steps taken to address it. The Provost will send a letter to the
student within ten business days of receipt of the faculty member's
report. The letter will reiterate the incident, describe the sanctions,
and inform the student of their right to appeal. The report and letter
from the Provost will remain on file in the Academic Affairs Office
until five years after the student's graduation or severance from Coe.
The Vice President for Student Life will be notified that a report has
been filed. Information in the file will be confidential, to be shared
only at the discretion of the Provost and Dean of the Faculty for a
legitimate educational or legal purpose.

\subsection{Appeals Process}\label{appeals-process}

The student may appeal the charge and/or the sanction within ten
business days of receiving the Provost's letter of notice by emailing
the Provost and Dean of the Faculty requesting an appeals hearing.
Students wishing to appeal are strongly encouraged to consult with the
Director of Academic Achievement, who has been designated by the College
to provide information and advocacy in these matters.

The Provost's Office will convene an Academic Integrity Appeals Board
within ten business days of receipt of the request for appeal to hear
the appeal. The Academic Integrity Appeals Board will be chaired by
either the chair or co-chair of the Academic Policies Committee.
Additional members of the board will include: one additional faculty
member on the Academic Policies Committee, the senior Student Life
officer, one faculty member from the Committee on Admission, Retention,
and Enrollment, and the Associate Dean for Student Academics. ~In any
hearing the Provost may replace the Associate Dean for Student Success,
or other member of the board to avoid conflicts of interest. The student
may choose to have a faculty or staff member present as an observer. The
student and the instructor will each appear as witnesses and each may
request that other pertinent witnesses appear.

A majority vote of the Appeals Board is necessary to uphold or overturn
a sanction. If a sanction is overturned, the Appeals Board may impose a
lesser sanction. The Appeals Board will submit a written finding which
will be sent to the student and the faculty member(s) involved in the
case and which will become part of the student's file. If the appeal's
outcome is to overturn the dishonesty charge, the report in the Academic
Affairs Integrity file will be removed.

\chapter{ACADEMIC STANDING}\label{academic-standing}

All students are expected to meet the College's standards for academic
performance (see chart below). These are in place to keep students on
track towards meeting the graduation requirements of the college (32
credits with a cumulative GPA of 2.0). Students who do not meet Coe's
academic standards may be put on probation or suspended from the
College.

To make sure that students are aware when their academic standing is in
jeopardy, a series of communications are issued throughout an academic
term including D/F notices, academic warnings, and notices of academic
probation. Unless otherwise stated, these communications will be sent to
the student by email and by a letter to the student's Coe mailbox.
Additionally, an email will be sent to the student's advisor.

\subsection{D/F Notices}\label{df-notices}

Academic notices are issued at midterm to students who earn a D, F, or
NP in one or more courses. Students receiving these notices are expected
to meet with their academic advisor to identify appropriate support.

\subsection{Academic Warnings}\label{academic-warnings}

Academic warnings will be issued to students who earn a term grade point
average below 2.00. Students receiving academic warnings are strongly
encouraged to re-evaluate their current approach to their studies and/or
their academic plans. They are expected to meet with their academic
advisor and/or a Learning Commons staff member.

\section{Academic Probation}\label{academic-probation}

The Academic Standing Committee reviews academic records of all students
at the conclusion of both the Fall and Spring Terms. A student is placed
on academic probation if the cumulative GPA falls below the minimum GPA
levels listed below:

\begin{longtable}[]{@{}ll@{}}
\toprule\noalign{}
\textbf{Total Course Credits Attempted} & \textbf{Minimum GPA Levels} \\
\midrule\noalign{}
\endhead
\bottomrule\noalign{}
\endlastfoot
0.0 -- 4.99 & l.50 \\
5.0 -- 7.9 & l.75 \\
8.0 -- 11.9 & 1.80 \\
12.0 -- 15.9 & 1.90 \\
16.0+ & 2.00 \\
\end{longtable}

To return to good standing and be removed from academic probation, a
student shall earn a cumulative GPA greater than the threshold for
academic probation as specified above. If the student's cumulative GPA
decreases at the end of any term on probation, the student is subject to
academic suspension. Students who are placed on academic probation have
a maximum of two (2) consecutive terms, excluding May and Summer, to
return to good standing before they are subject to academic suspension.

Students who are on academic probation: - Are expected to comply with
any requirements outlined by the Academic Standing Committee. - May not
elect P/NP grading in lieu of a letter grade. - May not be excused from
attending class to participate in extra-curricular activities.

The status of academic probation is noted on a student's transcript.

Students on probation who are not meeting the conditions of their
probation and not performing at a passing level in their courses may be
withdrawn from the College during the term on the recommendation of the
Academic Standing Committee. If the student had previously appealed a
suspension, the suspension is reinstated. Otherwise, the student will
exit the College with a leave of absence and is eligible to return on
probation the following term. Students may appeal this decision to the
Provost.

\section{Academic Suspension}\label{academic-suspension}

Students who do not meet the conditions of their academic probation or
who fall below retention thresholds (see below) will be considered for
academic suspension by the Academic Standing Committee, and, if
suspended, will be unable to take courses at Coe College for a period of
at least one academic year. (One academic year is defined as a Fall and
Spring Term, and does not count May or Summer Terms.) A registration
hold preventing registration will be placed on the student's account. A
student already registered for the following term will be dropped from
those courses.

\begin{longtable}[]{@{}ll@{}}
\toprule\noalign{}
\textbf{Total Course Credits Attempted} & \textbf{Minimum GPA Levels} \\
\midrule\noalign{}
\endhead
\bottomrule\noalign{}
\endlastfoot
0.0 -- 4.99 & 0.50 \\
5.0 -- 7.9 & l.00 \\
8.0 -- 11.9 & 1.50 \\
12.0 -- 15.9 & 1.70 \\
16.0 -- 19.9 & 1.80 \\
20.0 -- 23.9 & 1.90 \\
24 + & 2.00 \\
\end{longtable}

Students who are placed on suspension for academic reasons will receive
a letter notifying them of their suspension at their home residence, by
registered mail, as well as a letter in their Coe mailbox and an email
to their Coe email. Suspended students have the right to appeal for
immediate readmission if they can provide evidence of circumstances that
would warrant reconsideration by the deadline indicated in their
suspension letter.

To appeal for immediate readmission:

\begin{itemize}
\tightlist
\item
  A suspended student appealing for readmission must submit an appeals
  letter by mail, email, or delivery in person to the Office of the
  Provost. The letter must include an explanation of any circumstances
  that affected the student's academic performance. The letter should
  provide a detailed plan with specific actions that the student will
  take to improve their academic standing, and explain how the student
  will overcome the obstacles that affected their academic performance.
  Student appeals will not be considered if a written statement is not
  received by the deadline.
\item
  Although not mandatory, the suspended student appealing for
  readmission is urged to schedule to meet with the Academic Standing
  Committee on the date designated in the suspension letter. At this
  meeting the student will have the opportunity to respond to questions
  the Committee may have on the circumstances outlined in the student's
  letter requesting readmission.
\item
  A suspended student appealing for readmission is encouraged to request
  a letter of support for immediate readmission from a faculty advisor
  or other faculty or staff member at Coe, if the letter can shed light
  on the student's ability and motivation to do well in future academic
  endeavors. Similarly, the student is encouraged to consult with one or
  more faculty or staff members to prepare a strong case for
  readmission.
\end{itemize}

Students who have been suspended for academic reasons from the College
once, can return to the College either through a successful appeal for
immediate readmission or through readmission after leaving for at least
one academic year. A student who is suspended for academic reasons more
than once cannot appeal for immediate readmission and must leave the
College for at least one academic year.

\section{Interim Suspension}\label{interim-suspension}

Interim Suspension is a situation where it is determined that a
student's continued presence at the college constitutes an immediate
threat of harm to the student, other individuals, or to the stability
and continuance of normal College functions. The Provost and Dean of
Students or their designee may suspend a student pending disciplinary
proceedings. Such suspension may become effective immediately and
without prior notice.

Interim suspension shall be considered an excused absence until the
conclusion of formal hearings. The student will be offered the
opportunity to make up any academic work missed during the time in which
the interim suspension was imposed. It is the student's responsibility
to make specific arrangements with faculty members to complete academic
work. The Dean of Students will initiate communication with the
appropriate faculty.

\section{Readmission Following
Suspension}\label{readmission-following-suspension}

Coe College's academic suspension policy allows students, who are not
readmitted immediately through appeal, to apply for readmission after at
least one academic year has passed.

To be considered for readmission, the student must submit a letter of
appeal to the Office of the Provost. In the appeals letter, the student
must present evidence that demonstrates how the circumstances that led
to the student's academic suspension have been addressed, and how the
student plans to be successful in his or her academic future.
Readmission is determined by the Academic Standing Committee. The
Academic Standing Committee reserves the right to conduct its own
investigation, review the case, and make a final decision concerning the
student's reinstatement to the College. When appropriate, certain
academic stipulations may be applied. If a suspended student provides
evidence of successful course completion elsewhere and/or written
evidence of motivation and maturity necessary to be academically
successful at Coe, the student may be readmitted on probation at Coe.

\section{Exiting the College}\label{exiting-the-college}

The exit process at Coe College, whether through withdrawing or taking a
leave of absence, is initiated by the student. The exit procedure is
initiated in the Learning Commons (Stewart Memorial Library) and begins
with an interview with the Associate Dean of Student Academics or the
Director of the TRIO-Academic Achievement Program. At the interview, the
student is given an official exit form on which to secure signatures
from the following: the Student Financial Services Office, to verify a
balance due or a credit to be refunded, as well as Student Loan
information, to be aware of financial aid adjustments; the Library, to
ascertain that all materials have been returned; and the Resident
Director of the student's residence hall, to arrange for room checkout.
The completed exit form is returned to the Learning Commons which will
then notify other pertinent areas of the student's withdrawal.

If a student is unable to complete the official withdrawal process, the
intent to withdraw or take a leave of absence can be communicated to one
of the following offices: Registrar, Student Financial Services, Student
Life.

If a student withdraws from all courses during a period of enrollment
for which he or she received financial aid, the Student Financial
Services Office will determine how much, if any, of the student's
financial aid proceeds must be returned to the College, based on a
federally mandated refund formula (see Return of Title IV
Funds/Institutional Refund Policy, p.~268).

Special consideration is given to students who withdrew due to a call to
active duty. Coe's ``Military Call Up/Refund''~and ``Readmission of
Service Member''~policies are published on the Admission/Financial Aid
webpage. Copies of these policies can be requested from the Student
Financial Services Office.

\section{Leave of Absence}\label{leave-of-absence}

A student may find it necessary to interrupt a program of study at the
College. Under this condition, the student may apply for a leave of
absence. A leave of absence may be granted for a period not to exceed 12
months. Students may extend a leave of absence for up to a total of 36
months by contacting the Office of the Registrar via email at
o-registrar@coe.edu. Students that do not renew their Leave of Absence,
or return to Coe, will be officially withdrawn from the College.
Coursework completed while on leave from the College is subject to the
same conditions as work in transfer.

\section{Requesting a leave of
Absence}\label{requesting-a-leave-of-absence}

Students planning to leave of absence from the College should consult
their academic advisor and then they \emph{must declare} their intent to
one of the following designated offices: Student Financial Services,
Registrar, Residence Life, or the Learning Commons. Please refer to the
previous section on exiting the College for the steps that follow.~

\section{On Leave of Absence}\label{on-leave-of-absence}

The Coe community is committed to supporting students while they are
away from the College. Thus, the Associate Dean of Student Academics may
assign a designee who will be the point of contact to each student on a
Leave of Absence. This designee will stay in contact with the student
during their time away from the College, as appropriate, and will assist
with the return process once the student is ready to resume coursework.~

\section{Returning to Coe after a Leave of
Absence}\label{returning-to-coe-after-a-leave-of-absence}

Students planning on returning to the College do not need to apply for
readmission, but must submit a statement of intent to re-enroll to the
Office of the Registrar and the Associate Dean of Student Academics, or
designee, who will assist with the return process. Students are strongly
encouraged to schedule a consultation with the Associate Dean of Student
Academics, or designee, by December 1st for returning the Spring Term,
or August 1st to return the Fall Term. If a student wants to register
during the regular registration time (November for Spring Term or late
March / early April for Fall Term), they will need to submit their
statement of intent and contact the Associate Dean of Student Academics,
or designee, at least two weeks before Registration. The exact date
of~Registration can be found on the Academic Calendar on the Coe website
under the Academics tab.~

\section{Withdrawal from the College and
Readmission}\label{withdrawal-from-the-college-and-readmission}

Admission for work toward a degree terminates and the student is
considered withdrawn from the College if:

\begin{enumerate}
\def\labelenumi{\arabic{enumi}.}
\tightlist
\item
  A full-time student does not enroll at Coe for the next term
  (excluding May Term) and has not completed a Leave of Absence form.
  This does not apply to students in College-approved off-campus study
  programs.
\item
  A part-time student does not enroll for a course at Coe in a 12-month
  period and has not completed a Leave of Absence form.
\end{enumerate}

Students wishing to resume work toward a degree, once admission status
has terminated, must apply for readmission.

Students previously enrolled at Coe and readmitted after an absence of
two years or more may request that all previous work at Coe be
re-evaluated by the Registrar on the same basis as credits offered in
transfer. Re-evaluation means that only courses with grades of C or
better will be counted for credit toward graduation. All courses
affected by the re-evaluation and the grade earned for each course will
remain on the student's permanent record but will not factor into the
cumulative GPA or be counted toward graduation.

\chapter{TRANSCRIPT EVALUATION
POLICIES}\label{transcript-evaluation-policies}

Official transcripts for courses taken at Coe College can only be issued
by the Office of the Registrar and only after the office has received a
written request and payment from the student.

\section{General Policy On Transfer
Credit}\label{general-policy-on-transfer-credit}

To honor its mission and to preserve its academic integrity as a liberal
arts institution, the College accepts a course in transfer for the
equivalent earned credit (4 semester hours = 1 course credit), if that
course meets the spirit of the College's mission and is from a
regionally accredited institution. In cases where it is unclear whether
the course would be acceptable for transfer credit, the Registrar and
the appropriate department chair will consult. Other exceptions are
referred to the Committee on Petitions.

\section{Evaluation Of Credits In
Transfer}\label{evaluation-of-credits-in-transfer}

The Office of the Registrar is responsible for the evaluation of
transfer credit. Credits accepted in transfer do not affect the
cumulative GPA. Grades for the credits accepted are not recorded on
Coe's transcript. Thus, transfer credits increase only the total courses
attempted and the total course credits earned. Credit is not accepted
for coursework earning a grade below ``C'' (2.0 on a 4.0 scale).

\section{Junior or Community College
Credit}\label{junior-or-community-college-credit}

No more than 50\% of the course credits required for a degree at Coe
will be accepted in transfer from 2-year colleges. A maximum of 16
credits will be accepted as transfer credit from 2-year regionally
accredited institutions. Transfer students who complete a regionally
accredited A.A. degree program or a regionally accredited
college-parallel A.S. degree program will be accorded junior status (16
course credits) at Coe. Transfer students who complete an A.A., A.S., or
A.A.S. degree from a regionally accredited institution with which Coe
has a specific articulation agreement will be awarded credit consistent
with that agreement.

\section{European Credit Transfer System
(ECTS)}\label{european-credit-transfer-system-ects}

ECTS credits are a relative rather than an absolute measure of student
workload.~ They specify how much of a year's workload a course unit
represents at the institution or department allocating the credits.~ECTS
is thus based on a full student workload and not limited to contact
hours only. In ECTS, 60 credits represent the workload of a normal
undergraduate academic year of study and normally 30 credits for a
semester and 20 credits for a term. Thus, ECTS credits will~normally~be
transferred to Coe College at a rate of 7.5 ECTS credits: 1 Coe credit.

\section{Occasional Transfer Credit For Degree-Seeking
Students}\label{occasional-transfer-credit-for-degree-seeking-students}

Degree-seeking students sometimes wish to transfer credit from another
institution toward their degree at Coe. Such credit must be approved in
advance of completion of the course by the Registrar. Departments must
approve in advance any courses counting toward a major, a minor, or
teacher certification requirements. Credit from junior or community
colleges is not accepted for students who have junior or higher status
at Coe.

\section{Evaluation of Credits For
Graduation}\label{evaluation-of-credits-for-graduation}

The Office of the Registrar certifies the completion of general degree
(see p.~16) and general education requirements (see p.~23). Credits
toward a major, minor, endorsement, license, authorization, etc. are
approved by the appropriate department chair, administrative
coordinator, or by Academic Policies Committee for interdisciplinary
majors.

\section{Credit By Examination}\label{credit-by-examination}

A maximum of eight course credits in satisfaction of degree requirements
may be applied from credit earned through the Advanced Placement Program
and International Baccalaureate.

\section{Updating Coursework}\label{updating-coursework}

In the natural course of reviewing academic records, a student may be
required to repeat certain courses (or appropriate substitutes) taken
more than four years prior to the review to bring studies in those areas
up to date. Review cases may be brought to the Academic Policies
Committee by any member of the faculty, and this committee makes the
final decision.

\section{Advanced Placement (AP)}\label{advanced-placement-ap}

Coe College's Advanced Placement code is 6101.

Coe College grants college credit for approved Advanced Placement
examination scores of 4 or 5. AP credit cannot be used to fulfill any
general education requirement, including First-Year Seminar, Liberal
Arts Core, Diverse Cultural Perspectives, the writing emphasis
requirement, credit in transfer to reduce the number writing emphasis
courses required, or academic practicum. Each AP exam may earn 1.0
course credit up to a maximum of eight course credits towards
graduation. Credit is granted upon receipt of the results of the
examination, which must be received directly from the Educational
Testing Service. Further information regarding Advanced Placement
examination reporting to the College may be obtained in the Office of
the Registrar and on Coe's website. If an equivalent course is taken at
Coe College, the AP credit is removed from the transcript.

\section{International Baccalaureate
(IB)}\label{international-baccalaureate-ib}

Coe College credit may be awarded for International Baccalaureate work.
Students may earn 1.0 course credit for each higher-level examination
score of 5, 6, or 7 to a maximum of eight course credits. No credit is
granted for standard level examinations. Students may not receive
college credit for both AP and IB in areas of similar content. IB credit
cannot be used to fulfill any general education requirement, including
First-Year Seminar, Liberal Arts Core, Diverse Cultural Perspectives,
the writing emphasis requirement, credit in transfer to reduce the
number of writing emphasis courses required, or academic practicum.
Credit is granted upon receipt of the results of the examination, which
must be received directly from IB. Further information regarding
International Baccalaureate examination reporting to the College may be
obtained in the Office of the Registrar and on Coe's website. If an
equivalent course is taken at Coe College, the IB credit is removed from
the transcript.

\section{Class Attendance}\label{class-attendance}

Regular class attendance is expected, although the instructor of each
course sets the standard expected to be met by the students. The College
expects attendance on all scheduled days, including the first and last
day of a term, as well as the class days immediately preceding and
following College holidays. Students officially representing the College
are excused as necessary prior to the absence. Students on academic
probation are not excused from attending class to participate in
extra-curricular activities.

\section{Final Exams}\label{final-exams}

The final exam schedule is published by the Registrar. It is expected
that final exams will be given during the time scheduled for each
course. There may be extraordinary cases when an individual student has
a compelling reason for taking an exam at a time other than that
scheduled. In such cases the instructor may properly decide to let that
individual take the exam at another time. Having more than two exams
scheduled on one day would justify allowing a student to take a third
exam on another day. The instructor of the course with the highest
course number will move the exam, for that student only, to a mutually
convenient time.

\section{Participation In
Commencement}\label{participation-in-commencement}

Students can participate in Commencement when they have met all the
requirements for graduation for one of Coe's degrees, their financial
obligations to Coe College are met, and they have completed their Intent
to Graduate form and it has been approved. In addition, all students,
unless excused in writing by the Provost, must complete a survey
assessing their educational experience at Coe. All pending graduates
must participate in Commencement exercises unless excused in writing by
the Registrar..

\section{Pending Graduates}\label{pending-graduates}

Students who need to earn no more than 2.0 course credits to complete
all graduation requirements may be permitted to participate in
Commencement exercises as long as the remaining course credits are
scheduled to be completed by August 15 of that year. Pending graduates
will be listed in the Commencement booklet, but since Latin honors (see
Section 13.2.2) are bestowed only after all graduation requirements are
met, they cannot be listed in the Commencement booklet for Latin honors,
although any earned honors will be listed on their diploma and
transcript.

\chapter{STUDENT NOTIFICATION OF RIGHTS UNDER
FERPA}\label{student-notification-of-rights-under-ferpa}

The Family Educational Rights and Privacy Act (FERPA) affords students
certain rights with respect to their education records. They are:

\begin{itemize}
\tightlist
\item
  the right to inspect and review the student's education records within
  45 days of the day the College receives a request for access. Students
  should submit to the Office of the Registrar written requests that
  identify the record(s) they wish to inspect. The Registrar will make
  arrangements for access and notifies the student of the time and place
  where the records may be inspected. If the records are not maintained
  by the Office of the Registrar, the Registrar shall advise the student
  of the correct official to whom the request should be addressed.
\item
  the right to request the amendment of the student's education records
  that the student believes are inaccurate or misleading. Students may
  ask the College to amend records that they believe are inaccurate or
  misleading. They should write the College official responsible for the
  records, clearly identify the part of the record they want changed,
  and specify why it is inaccurate or misleading. FERPA was not intended
  to provide a process to be used to question substantive judgments
  which are correctly recorded. The rights of challenge are not intended
  to allow students to contest, for example, a grade in a course because
  they felt a higher grade should have been assigned. If the College
  decides not to amend the record as requested by the student, the
  College will notify the student of the decision and advise the student
  of his or her right to a hearing regarding the request for amendment.
  Additional information regarding the hearing procedures will be
  provided to the student when notified of the right to a hearing.
\item
  the right to consent to disclosures of personally identifiable
  information contained in the student's education records, except to
  the extent that FERPA authorized disclosure without consent. One
  exception which permits disclosure without consent is disclosure to
  school officials with legitimate educational interests. A school
  official is a person employed by the College in an administrative,
  supervisory, academic or research, or support staff position
  (including law enforcement unit personnel and health staff); a person
  or company with whom the College has contracted (such as an attorney,
  auditor, collection agent, or official of the National Student
  Clearinghouse); or a student serving on an official committee, such as
  an admission, petitions, retention, honors recognition, disciplinary,
  or grievance committee, or assisting another school official in
  performing his or her tasks. A school official has a legitimate
  educational interest if the official needs to review an education
  record in order to fulfill his or her professional responsibility.
\item
  the right to file a complaint with the U.S. Department of Education
  concerning alleged failures by the College to comply with the
  requirements of FERPA.
\end{itemize}

A student is a person who attends or has attended Coe College, as
determined by matriculation and enrollment by the first date of an
academic term. Coe College obtains written permission from the student
before releasing any information from a student's educational record in
most cases. However, as the law allows, on a case-by-case basis,
appropriate parts of a student record may be disclosed, without consent
of the student, to the following parties:

\begin{itemize}
\tightlist
\item
  college employees who have a legitimate need to know.
\item
  persons who need to know in cases of health and safety emergencies.
\item
  accrediting organizations to carry out accrediting functions.
\item
  appropriate parties in connection with financial aid to a student.
\item
  federal, state, and local governmental officials for purposes
  authorized by law.
\item
  individuals who have lawfully obtained court orders or subpoenas.
\item
  organizations conducting educational studies for the College.
\item
  courts during litigation between the College and the student or
  parent.
\item
  victim of crime of violence after final results of a disciplinary
  hearing.
\item
  public after disciplinary proceedings determine student committed
  crime of violence.
\end{itemize}

In many situations, complaints relative to FERPA can be resolved with
the College on an informal basis by contacting the Registrar, in the
lower level of Voorhees Hall.

\textbf{To file a FERPA complaint} with the U.S. Department of
Education, contact the office that administers FERPA at: Family Policy
Compliance Office, U.S. Dept. of Education, 400 Maryland Avenue SW,
Washington, DC 20202-4605.

\part{COURSES OF INSTRUCTION}

\chapter{ACADEMIC PROGRAMS}\label{academic-programs}

\subsection{Asia Term}\label{asia-term-1}

\textbf{ASC-195 Asian Tonal Languages}\\
An introduction to Thai, Vietnamese, and other tonal Asian languages.
Emphasis is on basic communication as well as the distinguishing
features of languages that use tones as part of their linguistic system.

\textbf{ASC-196 Modern South East Asia}\\
An introduction to several Asian cultures such as Thai, Cambodian, and
Vietnamese. This course varies depending on the field of the supervising
Coe faculty member.

\textbf{ASC-444 Independent Study}\\
A student-designed study of some feature of Asian culture, arranged in
consultation with the supervising Coe faculty member. May be taken more
than once for a maximum of 2.0 credits. Prerequisite: consent of
instructor.

\subsection{Community-Based Project}\label{community-based-project}

\textbf{CBP-325 Community-Based Project}\\
Supervised work on service projects proposed by external constituencies
such as non-profit institutions, community agencies, and government
organizations. Students learn about the issues, problems, and techniques
associated with developing, organizing, and participating in projects
that address and solve real-world problems, as well as provide services
and benefits to community and project sponsors. Authorization for the
community-based project is determined by the supervising faculty member.
May be taken more than once. A minimum of 140 hours of work is required.
P/NP basis only.

\subsection{First-Year Seminar}\label{first-year-seminar}

\textbf{FS-110 First-Year Seminar}\\
Required of all first-year students. The First-Year Seminar introduces
students to college-level study with emphasis on critical thinking,
writing, speaking, and research skills. Faculty offer first-year-only
topics courses exploring issues from multiple perspectives within or
across disciplines. Students in all sections prepare portfolios of their
written work and attend a variety of cultural events on and off campus.
Seminar instructors also serve as the primary academic advisors for
first-year students in their respective sections. First-Year Seminars
are writing emphasis and cannot fulfill any distributional, cultural
perspective, or major requirements.

\subsection{World Language}\label{world-language}

\textbf{FSA-100 Foreign Study Abroad}\\
Study abroad during May Term supervised by a Coe College faculty member,
with site visits to places of historical and cultural interest. When
appropriate, may count as a course in the major. May be taken more than
once if offered in different locations. Prerequisite: consent of
instructor. (Offered May Term only)

\textbf{FSA-146 Turkey: History and Culture}\\
Study abroad course on the history and culture of Turkey. An exploration
of a unique secular-Muslim society. Particular historical focus is on
Istanbul and ancient cities on Aegean coast. Prerequisite: consent of
instructor. (Offered May Term only)

\subsection{Internships}\label{internships}

\textbf{INT-494 Interdisciplinary Internship}\\
Supervised work or volunteer experience related to a student's career
interests. A minimum of 140 hours on-site or remote work experience is
required. Authorization and evaluation of the course credit for the
internship is determined by the department through which the student is
completing the internship. P/NP basis only. With departmental approval,
credit may be applied to a major only with consent of department chair.
Prerequisites: consent of the Internship Faculty Advisor.
\textbf{INT-499 Summer Internship -- Non-Credit Bearing}\\
Is a supervised summer work or volunteer experience related to a
student's career interests. The internship is not credit-bearing;
however, it can meet the practicum requirement. A minimum of 140 hours
on-site experience is required. The internship must be approved by the
Faculty Internship Advisor prior to registration. Students are required
to abide by the same guidelines as students completing credit-bearing
internships. Satisfactory completion determined by the supervising
faculty member. P/NP basis only. May be taken more than once.
Prerequisites: consent of the department in which the student is
completing the internship and completion of the Internship request form
housed on the College's online platform for internships.

\subsection{New York Term}\label{new-york-term-1}

See description, Section 14.1.1 New York Term is offered every other
Spring Term in odd-numbered years. There is an extra fee for New York
Term. All Coe financial aid applies, and students are eligible to apply
for additional financial aid based on the additional costs of the term.

\textbf{NYT-250 Fine Arts in New York City}\\
Consists of five 0.4-credit courses: art, film, music, theatre, and
dance. Students attend approximately 35 concerts, plays, and dance
performances and make frequent visits to museums, galleries, and
artists' studios. Each of these events is accompanied by discussion and
seminar sessions with members of the resident New York Term faculty.
(0.4 course credit for each course. Total of 2.0 course credits upon
completion of the five courses.)

\textbf{NYT-394 Internship in New York City}\\
Investigates a student's career interests through work or volunteer
experience. The internship is supervised by a faculty member of the
relevant department, in consultation with Coe's Center for Creativity
and Careers. P/NP basis only. Normally earns 2.0 credits, but may be
taken for 1.0 credits when combined with NYT-444.

\textbf{NYT-444 Independent Study}\\
A plan of study designed by the student in consultation with the
student's faculty advisor, and supervised by the on-campus faculty
advisor. Subject must be particularly appropriate for study in New York
City. Prerequisite: consent of instructor. (Corequisite: NYT-394 for 1.0
credit.)

\section{Occasional Courses}\label{occasional-courses}

Additional courses, not found in the Catalog, may be offered
occasionally and serve one or more of the following purposes:

\begin{enumerate}
\def\labelenumi{\arabic{enumi}.}
\tightlist
\item
  to provide the opportunity for research, creative, or other scholarly
  activity for an instructor jointly with interested students;
\item
  to explore and develop intellectual pursuits which are attractive to
  members of the faculty;
\item
  to respond to student requests for courses which are distinctive,
  unusual, or meet specific needs;
\item
  to bridge between two or more disciplines or curricular categories;
\item
  to serve as one means of developing and testing a possible permanent
  course.
\end{enumerate}

The courses listed in the following section are ones approved by the
College but without plans of being offered in the next few years. The
approved list of courses includes, but is not limited to, the following:

\begin{itemize}
\tightlist
\item
  CHM-104 Introduction to Forensic Science

  \begin{itemize}
  \tightlist
  \item
    An introduction to all aspects of forensic science from obtaining
    specimens to identifying the criminal with accurate forensic tests.
    The course teaches students a basic understanding of the laboratory
    tests and processes of forensic science.
  \end{itemize}
\item
  CHM-105 Food Chemistry

  \begin{itemize}
  \tightlist
  \item
    An introductory course that introduces chemical concepts in the
    context of cooking. Topics include the makeup, shape, and behavior
    of the four major classes of food molecules, effects of chemical
    structure on physical properties, the role of vitamins and
    nutrients, and basic principles of energy. Three class meetings per
    week.
  \end{itemize}
\item
  CHM-442 Materials Chemistry

  \begin{itemize}
  \tightlist
  \item
    Study of the structure and properties of modern materials, including
    glasses, polymers, metals, semiconductors, and superconductors.
    Mechanical, thermal, optical, magnetic, and electrical properties
    are examined and related to structure. Prerequisite:
    Electromagnetism (PHY-265) or Physical Chemistry I (CHM-341) or
    consent of instructor. Corequisite: Materials Physics and Chemistry
    Laboratory (CHM-442L).
  \end{itemize}
\item
  CHM-442L Materials Chemistry Laboratory

  \begin{itemize}
  \tightlist
  \item
    Measurement of structure/property characteristics of materials using
    a variety of instrumentation. Materials studied include glasses,
    polymers, metals, semiconductors, and superconductors. Prerequisite:
    previous or concurrent registration in Materials Physics and
    Chemistry (CHM-442). Corequisite: Materials Physics and Chemistry
    Laboratory (CHM-442L). (0.2 course credit)
  \end{itemize}
\item
  CHM-471 Advanced Chemistry Laboratory II

  \begin{itemize}
  \tightlist
  \item
    Spectroscopic investigations of chemical systems and applications of
    chemical instrumentation for analysis based on current chemical
    literature. One class period and two laboratories per week.
    Prerequisite: Physical Chemistry I (CHM-341).
  \end{itemize}
\item
  EDU-405 Understanding Early Adolescence

  \begin{itemize}
  \tightlist
  \item
    An extension of the study of human development, focusing on the
    growth and development of the middle school age child. Special
    attention is given to the emotional, physical, and cognitive
    characteristics and needs of middle school age children for teachers
    in grades five through eight.
  \end{itemize}
\item
  EDU-415 Middle School Curriculum and Instruction

  \begin{itemize}
  \tightlist
  \item
    Introduction to the organization, structure, and sequence of
    learning experiences for middle grade students. Addresses such
    issues as curriculum integration, teaching teams, pedagogical
    practices for middle school, and developmental appropriateness
    across the range of school subjects. Prerequisite: Practicum in
    Education (WE) (EDU-215).
  \end{itemize}
\item
  EDU-430 Middle School Social Studies

  \begin{itemize}
  \tightlist
  \item
    Overview of the middle school social studies curriculum. A content
    course with a primary focus on geography and its relation to U.S.
    and world history. Prerequisite: Practicum in Education (WE)
    (EDU-215). (0.5 course credit)
  \end{itemize}
\item
  EDU-440 Middle School Mathematics

  \begin{itemize}
  \tightlist
  \item
    Overview of the middle school mathematics curriculum. A content
    course with a primary focus on algebra, problem solving, and number
    theory. Prerequisite: Practicum in Education (WE) (EDU-215). (0.5
    course credit)
  \end{itemize}
\item
  INT-115 May Term in Southern Africa

  \begin{itemize}
  \tightlist
  \item
    Provides opportunities for students to interact with a wide variety
    of communities in southern Africa, all of which are currently
    stressed by impacts of the HIV/AIDS pandemic and climate change.
    Students learn how communities provide health care delivery,
    nutritional support and access to clean water, which are needed to
    sustain the quality of human life in the region. Experiences provide
    hands-on opportunities for students to assist community change in
    these areas. Prerequisite: consent of instructor. (Offered May Term
    only)
  \end{itemize}
\item
  NUR-255 Topics in Health Care

  \begin{itemize}
  \tightlist
  \item
    Offers selected topics on specific health care and/or nursing
    issues, problems, interventions, and theories. Content varies as
    determined by the instructor. May be taken more than once for
    credit, provided the topics are substantially different.
    Prerequisite: sophomore standing. (Offered on an occasional basis)
  \end{itemize}
\item
  PHY-111 Musical Acoustics

  \begin{itemize}
  \tightlist
  \item
    An exploration of the physical principles involved in the
    production, propagation, and perception of musical sounds. Topics
    include simple vibrating systems, properties of waves, and Fourier
    analysis. The primary emphasis is on musical instruments, including
    the voice, but some consideration is also given to room acoustics
    and human perception of sound. Previous musical experience is
    helpful, but not necessary. This course satisfies the non-lab
    science course requirement.\\
  \end{itemize}
\item
  PHY-112/-112L Holography and Optics \& Laboratory

  \begin{itemize}
  \tightlist
  \item
    The making and understanding of holograms are used as the focus for
    a basic physics course in waves and optics. Includes one two-hour,
    (0.0 course credit) lab per week. This course satisfies the general
    education laboratory science requirement.
  \end{itemize}
\item
  PHY-325/-325L Digital Electronics \& Laboratory

  \begin{itemize}
  \tightlist
  \item
    Integrated circuit devices and their applications: the basic logic
    gates, counters, displays, flip-flops, multiplexers, memories. Some
    acquaintance with DC circuit concepts and with the binary number
    system desirable. Includes one two-hour, (0.0 course credit) lab
    weekly. This course satisfies the general education laboratory
    science requirement.
  \end{itemize}
\item
  REL-240 Intertestamental Literature

  \begin{itemize}
  \tightlist
  \item
    A survey of literature composed by Jews during the Hellenistic and
    early Roman periods essential for understanding the emergence and
    development of Rabbinic Judaism, early Christianity and Islam
    (Apocrypha and Pseudepigrapha, Josephus, Philo and the Dead Sea
    Scrolls). Prerequisite: Introduction to Hebrew Bible (WE) (REL-105)
    or consent of instructor.
  \end{itemize}
\item
  REL-278 Mysticism

  \begin{itemize}
  \tightlist
  \item
    A survey of mystical literature in the world's religious traditions.
    This course also addresses the question of the nature of mystical
    experience as well as that of the relation between the mystical
    element of religion and religion as a whole.
  \end{itemize}
\item
  REL-306 Comparative Religion

  \begin{itemize}
  \tightlist
  \item
    A comparative study of the recurring themes and patterns found in
    various religions, past and present. Particular attention is paid to
    the meaning of religious ritual and myth, and the nature of
    religious experience. This course surveys several of the currently
    most influential theories regarding the nature of religion.
    Prerequisite: Eastern Religions (REL-106), Western Religions (WE)
    (REL-108), or consent of instructor.
  \end{itemize}
\item
  SOC-499 Career Related Independent Investigation

  \begin{itemize}
  \tightlist
  \item
    Investigation of a career opportunity through field placement and
    directed reading. This course does not satisfy any of the
    requirements for a major or minor in sociology. Prerequisite:
    declared major in sociology, second term sophomore standing, or
    consent of department chair.
  \end{itemize}
\end{itemize}

\subsection{Skills Development}\label{skills-development}

SKD-115 Summer Bridge Engages students in a one-week course that takes
place prior to Fall Orientation. Students participate in two
mini-classes taught by college faculty. Each mini-class is followed by
small group discussions regarding course content, learning strategies,
college expectations, etc. Additional workshop sessions are held
throughout Summer Bridge on topics such as: financial literacy, campus
resources, goal setting and academic planning. Summer Bridge also offers
opportunities to make social connections through a variety of informal
and planned activities. P/NP basis only. (0.3 course credit) SKD-120
Concepts of Individualized Learning Provides a weekly engagement with an
academic coach to identify interests and explore strengths, applying
this knowledge to the development of a personal academic plan. Topics
include: learning strategies, self-regulation, personal and professional
growth, goal setting, campus engagement and self-reflection. P/NP basis
only. (0.0 course credit) SKD-125 College Foundations Familiarizes
students with the skills and methods of study that lead to competence in
college coursework. Through self-assessment and reflection, students
determine strategies that increase satisfaction and success in the
college environment. P/NP basis only. (0.3 course credit)\\
SKD-130 Personal Finance for College Students A study of managing
finances and making financial decisions that college students encounter.
Areas of study for this project-oriented course include student loans,
credit cards, savings and investments, cars, living on campus versus
apartment living, savings needed for life immediately following
graduation, and travel. (0.5 course credit) INT-100 Professionalism and
Self Presentation Introduces students to the fundamentals of job-seeking
strategies and professional expectations. Students are guided through
development of professional materials, networking techniques, and
interviewing skills. This course does not satisfy Coe's practicum
requirement.

\subsection{Washington Term}\label{washington-term-1}

WSH-284 Topics in Washington, D.C. Experiential learning, study,
writing, and discussion dealing with various subjects related to the
nation's capital. Examples of recent and proposed topics include Art and
Architecture, Campaigns and Elections, Congressional Relations, and
Politics and Communications. WSH-286 Topics in Washington, D.C.:
Non-Western Perspectives Same as Topics in Washington, D.C. (WSH-284)
except the course focuses on topics related to non-Western cultures.
Examples of recent and proposed topics include Globalization and the
U.S.; and People, Politics and Cultures of the Middle East.\\
WSH-464 Washington Term Internship Seminar Group discussion of
internship experiences. Students are exposed to various research
methodologies, readings and guest speakers for understanding Washington
politics. The goal of the course is to expose students to
generalizations about politics and how their internships are either
confirming or challenging those generalizations. WSH-494 The Washington
Experience Internship with an organization related to national or
international politics in Washington, D.C., supervised by one of the
resident staff of the Washington Term. Students establish learning goals
and prepare a portfolio that documents their learning and places it in
the larger context of the literature on American or international
politics. P/NP basis only. One course credit may be counted toward a
major in political science with consent of department, if credit has not
already been received for Internship in Political Science (POL-494).
(2.0 course credits)

\chapter{DEPARTMENTAL PROGRAMS}\label{departmental-programs}

\section{ACCOUNTING, MANAGERIAL}\label{accounting-managerial}

See Section~\ref{sec-accounting}

\section{ACCOUNTING, PUBLIC}\label{accounting-public}

See Section~\ref{sec-accounting}

\section{African American Studies}\label{african-american-studies}

Rodgers (Administrative Coordinator).

The African American Studies major offers students the opportunity to
study African American culture through an interdisciplinary approach.
Students take courses in areas such as literature, history, and
sociology in order to gain a more complete understanding of the major
figures and movements that helped define the culture.

\subsection{African American Studies
Major}\label{african-american-studies-major}

A major in African American studies requires a minimum cumulative 2.0
GPA in all courses counted toward the major.

\begin{enumerate}
\def\labelenumi{\arabic{enumi}.}
\tightlist
\item
  AAM 107 Intro to African American Studies
\item
  +++MISSING INFO: c.aam/eng137.long +++
\item
  HIS 347 African American History
\item
  One of the following:
\end{enumerate}

\begin{itemize}
\tightlist
\item
  HIS 145 History of United States to 1865
\item
  HIS 155 History of United States since 1865
\item
  HIS 227 American Civil War
\end{itemize}

\begin{enumerate}
\def\labelenumi{\arabic{enumi}.}
\setcounter{enumi}{4}
\tightlist
\item
  Five additional courses, at least three of which are numbered 200 or
  above, chosen from the following: -\textbf{(a)} Elective courses:

  \begin{itemize}
  \tightlist
  \item
    AAM 217 Sport and Black Culture
  \item
    AAM 287 Topics in African American Studies
  \item
    AAM 367 Topics in AfricanAmericanLiterature
  \item
    AAM 387 Adv Topics in African American Stud
  \item
    +++MISSING INFO: c.aam447/457.long +++
  \item
    AAM 467 Seminar in African American Lit
  \item
    AAM 494 Internship in African American Stds
  \item
    COM 236 Intercultural Communication
  \item
    COM 357 Sex, Race, \& Gender in Media
  \item
    EDU 187 Human Relations
  \item
    HIS 145 History of United States to 1865
  \item
    HIS 155 History of United States since 1865
  \item
    HIS 227 American Civil War
  \item
    MU 157 Introduction to Jazz History
  \item
    SOC 247 Sociology of Race
  \end{itemize}
\end{enumerate}

-\textbf{(b)} Elective courses (when topic is appropriate): These
courses can count toward the major or minor as determined by the African
American Studies administrative coordinator. - ARH 297 Topics in Art
History: US Pluralism - COM 157 Introduction to Media Analysis - COM 161
Visual Rhetoric - COM 361 Communication \& Social Change - ENG 107
Exploring Literature:US Pluralism - ENG 207 Gender \& Lit:US Pluralism -
ENG 347 Study in Modern or Contemp Amer Lit - ENG 394 Directed Learning
in English - HIS 297 Women in America - HIS 444 Ind Study-History - HIS
472 Seminar American History I - HIS 473 Seminar American History II -
SOC 237 Topics in Sociology:U S Pluralism - SOC 464 Capstone Seminar in
Sociology - THE 488 Special Topics in THE/ FLM

\begin{enumerate}
\def\labelenumi{\arabic{enumi}.}
\setcounter{enumi}{5}
\tightlist
\item
  AAM 444 Ind Study-Afr-Am St (completed during the senior year)
\end{enumerate}

\subsection{African American Studies
Minor}\label{african-american-studies-minor}

\begin{enumerate}
\def\labelenumi{\arabic{enumi}.}
\tightlist
\item
  AAM 107 Intro to African American Studies\\
\item
  +++MISSING INFO: c.aam/eng137.long +++
\item
  HIS 347 African American History\\
\item
  One of the following:
\end{enumerate}

\begin{itemize}
\tightlist
\item
  HIS 145 History of United States to 1865
\item
  HIS 155 History of United States since 1865
\item
  HIS 227 American Civil War
\end{itemize}

\begin{enumerate}
\def\labelenumi{\arabic{enumi}.}
\setcounter{enumi}{4}
\tightlist
\item
  Two additional courses from either list \textbf{5a} or, when
  appropriate, \textbf{5b} as listed above for the major
\end{enumerate}

\subsection{COURSES IN AFRICAN AMERICAN
STUDIES}\label{courses-in-african-american-studies}

\begin{itemize}
\tightlist
\item
  \textbf{AAM 107 Intro to African American Studies} What is African
  American culture, and what accounts for its cultural distinctions?
  This course introduces students to the study of African American
  culture and the field of African American Studies. through an
  interdiciplinary approach (literature, history, music, art and film)
  students will examine central themes and key debates pertinent to
  African American culture and history from its beginning to present.
\item
  \textbf{AAM 137 African American Literature} Reading and discussion of
  the writings of African Americans, with emphasis on the 20th century.
  May include some relevant writings on African Americans by other
  groups. Study of the artistic values and of the social and cultural
  significance of these writings. May be taken more than once, with
  consent of African American Studies administrative coordinator,
  provided the topics are substantially different.
\item
  \textbf{AAM 217 Sport and Black Culture} Examines through cultural
  analysis the complex relationships between sport and Black culture.~
  This course addresses the way sport has evolved from being merely a
  physical activity to a cultural expression in Black communities.~ This
  course emphasizes the historical patterns and current conditions of
  Blacks; participation in sport through various articles, videos, and
  books.~ The course also examines how many Black people have used sport
  as a means of resistance, survival, and social mobility. Students
  learn to analyze cultural expression, to understand race and its
  continuing impact in American life, and to understand how various
  sports pursuits by Black athletes are invested with multiple
  meanings.~Major topics and themes covered include: the concept of
  race, Black culture, the historical presence of Black athletes in
  sport, their current impact, and their dominance in certain sports.
\item
  \textbf{AAM 227 Blackness \& Identity in America} Gathers a wide range
  of scholarship about race and identity to explore what race is, why it
  matters, racial dynamics in organizations, and how best to address
  them. Students will explore questions of identity, privilege,
  ethnicity, gender, and class. In this course, ``race'' is reviewed as
  a shorthand for the interconnected complexity of race, ethnicity,
  culture, and color, and will be carefully analyzed to distinguish
  among such terms and ideas. Students will critically analyze the
  social construct and the popular understandings of race and identity
  as reinforced through cultural institutions. Students will understand
  and evaluate the foundational concepts and theories of race and
  identity and synthesize their knowledge through research,
  presentation, and writing.
\item
  \textbf{AAM 287 Topics in African American Studies} Examines an
  important theme or subject specific to African American experiences
  and culture. Content varies and is determined by the instructor.
  Students learn to understand African American experiences in context,
  to analyze texts and events from multiple disciplinary perspectives,
  and to write in clear, analytical prose. May be taken more than once
  for credit, provided the topics are substantially different.
\item
  \textbf{AAM 367 Topics in AfricanAmericanLiterature} None
\item
  \textbf{AAM 387 Adv Topics in African American Stud} Examines an
  important theme or subject specific to African American experiences
  and culture. Content varies and is determined by the instructor.
  Students learn to understand African American experiences in context,
  to analyze texts and events from multiple disciplinary perspectives,
  and to write in clear, analytical prose. May be taken more than once
  for credit, provided the topics are substantially different.
  Prerequisite: Introduction to African American Studies (AAM-107) or
  consent of instructor.
\item
  \textbf{AAM 444 Ind Study-Afr-Am St} Study of individually chosen
  research topics in African American studies under the direction of a
  faculty member in the area. May be taken for an X status grade with
  consent of instructor prior to registration. Prerequisites:
  Introduction to African American Studies (AAM-107) or African American
  Literature (AAM-137); consent of African American Studies
  administrative coordinator and submission of a written proposal for a
  project.
\item
  \textbf{+++MISSING INFO: c.aam447/457.long +++} +++MISSING INFO:
  c.aam447/457.desc +++
\item
  \textbf{AAM 467 Seminar in African American Lit} Intensive study of
  selected works and subjects in African American literature. May be
  taken more than once, with consent of African American Studies
  administrative coordinator, provided the topics are substantially
  different. Prerequisites: junior standing and The Art of Literary
  Research (ENG-301). May be taken more than once.
\item
  \textbf{AAM 494 Internship in African American Stds} None
\end{itemize}

\section{ANTHROPOLOGY (Minor Only)}\label{anthropology-minor-only}

Fairbanks, Ziskowski

\subsection{Anthropology Minor}\label{anthropology-minor}

\begin{enumerate}
\def\labelenumi{\arabic{enumi}.}
\tightlist
\item
  \textbf{One} of the following:
\end{enumerate}

\begin{itemize}
\tightlist
\item
  ANT 109 Intro to Archaeology:Method \& Thry
\item
  ANT 116 Cultural Anthropology
\end{itemize}

\begin{enumerate}
\def\labelenumi{\arabic{enumi}.}
\setcounter{enumi}{1}
\tightlist
\item
  ANT 215 Ethnographic Methods
\item
  \textbf{One} of the following:
\end{enumerate}

\begin{itemize}
\tightlist
\item
  ANT 284 Topics in Anthropology/Archaeology
\item
  ANT 286 Topics in Anthropology:NWP
\item
  ANT 288 Topics Anthropology/Archaeology:DWP
\end{itemize}

\begin{enumerate}
\def\labelenumi{\arabic{enumi}.}
\setcounter{enumi}{3}
\tightlist
\item
  ANT 450 Anthropological Theory
\item
  \textbf{One} of the following:
\end{enumerate}

\begin{itemize}
\tightlist
\item
  ANT 484 AdvTop: Anthropology or Archaeology
\item
  ANT 486 Advanced Topics in Anthropology:NWP
\item
  ANT 488 Adv Top Anthro/Archaeo:DWP
\end{itemize}

\begin{enumerate}
\def\labelenumi{\arabic{enumi}.}
\setcounter{enumi}{5}
\tightlist
\item
  One additional anthropology course Students must take at least one
  course focused on archaeology and one focused on anthropology. The
  following courses do \textbf{not} satisfy any of the requirements for
  a minor in anthropology:
\end{enumerate}

\begin{itemize}
\tightlist
\item
  ANT 444 Independent Study: Anthropology
\item
  ANT 474 Research Participation:Anthropology
\item
  ANT 494 Internship in Anthropolgy
\end{itemize}

\subsection{Courses in Anthropology}\label{courses-in-anthropology}

\begin{itemize}
\tightlist
\item
  \textbf{ANT 109 Intro to Archaeology:Method \& Thry} An introduction
  to the theoretical approaches and field methodologies of archaeology.
  The goal of this class is to familiarize the student with the history
  and theoretical frameworks of archaeology, in addition to the
  scientific methods with which material culture is collected,
  investigated, and evaluated. Major topics and themes covered in this
  class include: archaeological theory, excavation and survey, artifact
  analysis, death, social systems, economy, religion, and ethical
  practices in cultural heritage management.
\item
  \textbf{ANT 115 Ancient Greece:Hist as Archaeology} Explores the
  various ways in which archaeology can inform our understanding of
  ancient Greek history by visiting the monuments, museums, and
  archaeological sites in modern Greece. The course emphasizes the
  combination of the historical textual evidence and the ancient
  physical evidence to enrich our understanding of Greek social,
  political, and cultural history. THe course also involves an
  examination of the evolution of ancient institutions and practices
  that range from the rise of the Athenian democracy, to the first
  theatrical productions of comedy and tragedy, to the establishment of
  the ancient Olympics. Archaeological theory, archaeiligical practice,
  and historical process constitute the nexus around which the course
  revolves. (Offered May Term only)
\item
  \textbf{ANT 116 Cultural Anthropology} An introduction to cultural
  anthropology, presenting its place within the broader discipline of
  anthropology and outlining its characteristic methodological and
  theoretical approaches to the study of human life. The course
  emphasizes the diversity of approaches to common human experiences in
  a variety of cultural contexts.
\item
  \textbf{ANT 125 Art \& Archaeology Classical World} A survey of the
  art, architecture, and archaeological remains of the Greek and Roman
  civilizations from early Aegean Bronze Age cultures to the fall of the
  Roman empire.
\item
  \textbf{ANT 215 Ethnographic Methods} An introduction to ethnographic
  research methods, presenting a number of tools for collecting and
  analyzing ethnographic data. The course emphasizes the relationship
  between research questions and the methodological tools used by
  anthropologists to study those questions. Prerequisite: Cultural
  Anthropology (ANT-116) or consent of instructor.
\item
  \textbf{ANT 284 Topics in Anthropology/Archaeology} A focused
  examination of an anthropological theme, theory or research method.
  Content varies and is determined by the instructor. May be taken more
  than once for credit, provided the topics are substantially different.
\item
  \textbf{ANT 286 Topics in Anthropology:NWP} A focused examination of
  an anthropological theme, theory or research method. Content varies
  and is determined by the instructor. May be taken more than once for
  credit, provided the topics are substantially different.
\item
  \textbf{ANT 288 Topics Anthropology/Archaeology:DWP} A focused
  examination of an archaeological theme, theory or research method.
  Content varies and is determined by the instructor. May be taken more
  than once for credit, provided the topics are substantially different.
\item
  \textbf{ANT 415 Ancient Greek Pottery Studies} Considers Greek
  ceramics from both scientific and iconological perspectives. Students
  study issues of production, construction and distribution and then
  look at Greek pottery iconographically, focusing on the major
  stylistic periods but emphasizing regional variations through Greece.
  Discussions of connoisseurship and themes prevalent in Greek
  vase-painting such as death, myth, and gender round out the course's
  survey of evidence. Prerequisite: consent of instructor.
\item
  \textbf{ANT 444 Independent Study: Anthropology} Independent study
  under faculty guidance of a research problem chosen by the student.
  May be taken for an X status grade with consent of instructor prior to
  registration. This course does not satisfy any of the requirements for
  a minor in anthropology. Prerequisites: Previous or concurrent
  enrollment in Anthropology Theory (ANT-450) and consent of department
  chair. (Offered by arrangement)
\item
  \textbf{ANT 450 Anthropological Theory} An examination of the
  historical development of anthropological theory, emphasizing
  American, British and French traditions from the 19th century through
  the present. Prerequisite: Cultural Anthropology (ANT-116) or consent
  of instructor.
\item
  \textbf{ANT 474 Research Participation:Anthropology} Individual or
  group investigation with a faculty member on a research topic or
  topics of mutual interest. The student must obtain approval for a
  specific project and make necessary arrangements prior to the term of
  registration for the course. May be taken for an X status grade with
  consent of instructor prior to registration. This course does not
  satisfy any of the requirements for a minor in anthropology.
  Prerequisites: Cultural Anthropology (ANT-116) and consent of the
  instructor. (Offered by arrangement)
\item
  \textbf{ANT 484 AdvTop: Anthropology or Archaeology} Same as ANT-284,
  except at an advanced level. Prerequisite: Introduction to Archaeology
  (ANT-109) or Cultural Anthropology (ANT-116) or consent of instructor.
\item
  \textbf{ANT 486 Advanced Topics in Anthropology:NWP} Same as ANT-286,
  except at an advanced level. Prerequisite: Cultural Anthropology
  (ANT-116) or consent of instructor.
\item
  \textbf{ANT 488 Adv Top Anthro/Archaeo:DWP} Same as ANT-228, except at
  an advanced level. Prerequisite: Introduction to Archaeology (ANT-109)
  or Cultural Anthropology (ANT-116) or consent of instructor.
\item
  \textbf{ANT 494 Internship in Anthropolgy} Placement with a
  career-related organization. A minimum of 140 hours on-site experience
  is required. P/NP basis only. This course does not satisfy any of the
  requirements for a minor in anthropology. Prerequisites: declared
  minor in anthropology, junior standing and consent of department
  chair. (Offered by arrangement)
\end{itemize}

\section{ART AND ART HISTORY}\label{art-and-art-history}

Cohen, Goodson, Knight-Lueth, Rogers (Chair), Thompson.

A major in art or art history is designed to provide a foundation in
theory and practice. Students focus on making and meaning, and gain an
understanding of visual and contextual history and analysis. In studying
the formal, technical and conceptual aspects of art and art history,
students learn to communicate effectively to broad and diverse
audiences.

An art major may also complete an art history minor, but only two
courses may count toward both the major and the minor.

\subsection{Art Major}\label{art-major}

A grade of ``C'' (2.0) or higher must be earned in all courses counted
toward a major in art.

\begin{enumerate}
\def\labelenumi{\arabic{enumi}.}
\tightlist
\item
  ARH-128 Introduction to Art History
\item
  Two 100-level ART- courses
\item
  Four of the following:

  \begin{itemize}
  \tightlist
  \item
    ART 201 Focus Course: Creative Process
  \item
    ART 202 Focus Course: Portfolio Development
  \item
    ART 203 Focus: ContemporaryArtisticPractice
  \item
    ART 211 Focus:ArtistStatements
  \item
    ART 212 Focus Course: Artist Websites
  \item
    ART 213 Focus Course: Critique
  \item
    ART 291 Focus Course: Topics in Studio Art
  \end{itemize}
\end{enumerate}

The Focus Course requirement is considered a coherent set of experiences
that can be tailored to student needs. The completion of four Focus
Courses is considered as a single course credit with respect to
Graduation Requirements (see Chapter~\ref{sec-graduation-requirements}
).

\begin{enumerate}
\def\labelenumi{\arabic{enumi}.}
\setcounter{enumi}{3}
\tightlist
\item
  Two 300-level ART- Courses
\item
  One additional 200-level Art History course
\item
  One of the following:

  \begin{itemize}
  \tightlist
  \item
    ARH 307 Modern and Contemporary Art
  \item
    ARH 310 Contemporary Art, Theory \&Criticism
  \end{itemize}
\item
  Junior Review
\end{enumerate}

A studio portfolio presentation for art majors that is assessed by the
art and art history faculty. Students must complete 20 course credits to
be eligible for the junior review, and it must be completed prior to
enrolling in ART-464 Senior Seminar I.

\begin{enumerate}
\def\labelenumi{\arabic{enumi}.}
\setcounter{enumi}{7}
\tightlist
\item
  ART 464 Senior Seminar I
\item
  ART 474 Senior Seminar II \& Senior Exhibit
\end{enumerate}

In addition to the course requirements, a student must participate in a
Junior Review with the art and art history department's faculty. This
must be completed no later than April of the junior year. The material
submitted at the Review will be related to the senior graduation
requirement, which is a demonstration of proficiency through a public
exhibition of artwork, including portfolio of slides, exhibit
announcement, and résumé.

\textbf{NOTE}: Scheduling may be planned to allow at least one term of
study off-campus at one of the approved foreign or domestic programs.
With departmental approval, courses from these programs may count as 100
or 200 level courses.

\subsection{Art Minor}\label{art-minor}

A grade of ``C'' (2.0) or higher must be earned in all courses counted
toward a minor in art.

\begin{enumerate}
\def\labelenumi{\arabic{enumi}.}
\tightlist
\item
  ARH 128 Introduction to Art History
\item
  Two 100-level ART- courses
\item
  Two 300-level ART- courses
\item
  One 200-level ARH- course
\end{enumerate}

Further information for all these requirements and programs is available
from the Art and Art History Department.

\subsection{Courses in Art}\label{courses-in-art}

\textbf{100-level courses} (\emph{except ART-130, see below}) are open
to any student and may offer the opportunity to explore two or more
areas of art making in a studio environment. ART-130 Art in the
Elementary Classroom does not satisfy any of the requirements for a
major or minor in Art or Art History. It also does not count as a
100-level art course needed for some courses as a prerequisite.

\begin{itemize}
\tightlist
\item
  \textbf{ART101 Art Appreciation} None
\item
  \textbf{ART 102 Sculpture: Material Investigations} Focuses on how
  material choices inform meaning. Students will use a variety of
  traditional and non-traditional materials to make sculptural art
  objects.
\item
  \textbf{ART 103 Ceramic Sculpture} Provides instruction on the
  creation of hand-built ceramic sculptures. Projects will utilize
  pinch, coil building, soft and stiff slab construction, and additive
  and subtractive processes.
\item
  \textbf{ART 105 Narrative Ceramics} Studies the history and evolution
  of ceramics as a storytelling medium. Projects will include
  tile-making, coil-building, surface decoration, and sculptural
  techniques as a means to explore how artists have used clay to create
  compelling narratives.
\item
  \textbf{ART 115 Drawing} Introduces a variety of drawing media,
  including graphite, charcoal, brush and ink, pastels, and collage.
  Students are challenged to observe the visual world around them and to
  respond to that observation with interpretive choices in mark-making
  and composition to produce visual expression. Media may include
  graphite, charcoal, brush and ink, pastels, and collage.
\item
  \textbf{ART 125 Painting} Introduces paint media to represent,
  amplify, and interpret the world. Students are also encouraged to
  discover the possibilities of color, shape, texture and mark-making.
\item
  \textbf{ART 130 Art in the Elementary Classroom} Provides an overview
  of the role of art in the elementary curriculum. Students learn to
  incorporate art activities into other content areas and gain an
  understanding of the objective of elementary classes taught by
  elementary art specialist teachers. This course does not satisfy the
  College's general education fine arts core group requirement nor does
  it satisfy major or minor requirements or prerequisites for Art or Art
  History. Prerequisite: admission to the Education Program or approval
  of the education department. (0.5 course credit)
\item
  \textbf{ART 131 Scultpure: Do, Undo, Redo} Practices the inherent
  creative aspects of making, unmaking and remaking in response to these
  various states of transformation. Artworks are made using clay and
  other sculptural materials while documenting the exploratory processes
  through drawing, photography and more.
\item
  \textbf{ART 135 Ceramics: Form, Function \& Meaning} Focuses on the
  creation of wheel thrown and handbuilt objects within the context of a
  social and cultural lens. Explores the relationship between making and
  meaning in the 21st century. Projects focus on fuctional work and its
  uses.
\item
  \textbf{ART 145 Digital Studio} Explores contemporary digital imaging
  and design. Students produce a series of studio projects ranging from
  digital photography, collage, gif animation, vector graphics, and
  two-dimensional design. Includes introductory-level instruction of
  digital cameras and Adobe Creative Suite: Lightroom, Photoshop,
  Illustrator.
\item
  \textbf{ART 150 Time Based Media} Explores the aesthetic and
  experiential qualities of time. The course includes readings,
  discussions, and screenings of historical and contemporary works plus
  hands-on studio projects using video, sound, performance,
  installation, and more. Includes introductory-level instruction in
  digital cameras and Adobe Lightroom and Premiere Pro. Through media
  production, students cultivate a range of technical skills plus a
  critical understanding of media culture.
\item
  \textbf{ART 151 Layers of Meaning} Pursues meaning through
  multiplicity of imagery. Some work may be done in collage, but other
  materials and techniques, such as drawing, painting, and digital art,
  are used to juxtapose and layer imagery and meaning.
\item
  \textbf{ART 155 Photography: Light Writing} Examines the ways that a
  photographic image can be viewed and interpreted. The course includes
  readings and research on the masters of photography. Technical skills
  include historic black and white analog photography including manual
  camera controls, film and print processing.
\item
  \textbf{ART 175 Printmaking} Focuses on traditional and contemporary
  printmaking techniques, introducing students to the fundamentals of
  materials including paper, ink, presses, and image processing. This
  course explores the development of technical, compositional, and
  conceptual skills through the evolution of printmaking's history of
  multiplicity.
\item
  \textbf{ART 191 Topics in Studio Art} Focuses on a specific theme or
  topic. Topics vary. May be taken more than once for credit, provided
  the topics are substantially different.
\end{itemize}

\textbf{Focus Courses}

Focus Courses are seven-week courses designed to prepare the student for
a serious artistic practice.

\begin{itemize}
\tightlist
\item
  \textbf{ART-201 Focus Course: Creative Process} Applies the tools and
  methods of creative processes. An abbreviated course offered seven
  weeks of a term designed to prepare the student for a serious artistic
  practice. Students propose, iterate, and produce an object or set of
  objects in an exploration of their own creative process. Prerequisite:
  two ART- courses (each 1.0 course credit) or consent of instructor.
  (0.25 course credit)
\item
  \textbf{ART-202 Focus Course: Portfolio Development} Prepares students
  for developing a portfolio that reflects the breadth and skills of
  their own artistic practice and production. Prerequisite: two ART-
  courses (each 1.0 course credit) or consent of instructor. (0.25
  course credit)
\item
  \textbf{ART-203 Focus Course: Contemporary Artistic Practice}
  Highlights the art, writings, routines and habits of contemporary
  artists and practitioners. An abbreviated course offered seven weeks
  of a term designed to prepare the student for a serious artistic
  practice. Prerequisite: two ART- courses (each 1.0 course credit) or
  consent of instructor. (0.25 course credit)
\item
  \textbf{ART-211 Focus Course: Artist Statements/Artist Talks} Guides
  students through the construction of artist statements and artist
  talks to present their work to a wider public. An abbreviated course
  offered seven weeks of a term designed to prepare the student for a
  serious artistic practice. Prerequisite: two ART- courses (each 1.0
  course credit) or consent of instructor. (0.25 course credit)
\item
  \textbf{ART-212 Focus Course: Artist Websites} Guides students through
  best practices in the process of developing and maintaining a
  professional website. An abbreviated course offered seven weeks of a
  term designed to prepare the student for a serious artistic practice.
  Prerequisite: two ART- courses (each 1.0 course credit) or consent of
  instructor. (0.25 course credit)
\item
  \textbf{ART-213 Focus Course: Critique} Offers students the
  opportunity to practice and improve critique skills. An abbreviated
  course offered seven weeks of a term designed to prepare the student
  for a serious artistic practice. Prerequisite: two ART- courses (each
  1.0 course credit) or consent of instructor. (0.25 course credit)
\item
  \textbf{ART-291 Focus Course: Topics in Studio Art} Focuses on a
  specific theme or topic. An abbreviated course offered seven weeks of
  a term designed to prepare the student for a serious artistic
  practice. Topics vary. May be taken more than once for credit,
  provided the topics are substantially different. Prerequisite: two
  ART- courses (each 1.0 course credit) or consent of instructor. (0.25
  course credit)
\end{itemize}

\textbf{300-level courses} offer students the opportunity to expand
their technical and conceptual skills at a more advanced level. Courses
numbered ART 300--349 are offered without prerequisite and open to any
student. Courses numbered ART 350--399 have one or more prerequisites.

\begin{itemize}
\tightlist
\item
  \textbf{ART-301 Socially Engaged Art} Invites collaboration with
  individuals, communities, and institutions in the creation of
  participatory art. The genre explores social forms such as dinner
  parties, conversations, and projects that intervene or intersect with
  real-world systems. Coursework is done collaboratively and
  independently on projects that critically engage with contemporary
  issues and explore a variety of contexts. A variety of media will be
  used as a means to define interests and inform social interventions.
\item
  \textbf{ART-313 Color and Design} Examines the theory and practice of
  color, with emphasis on the use of color as a compositional element.
\item
  \textbf{ART-315 Installation Art} Creates art environments that offer
  a unique experience for the viewer. Projects will include research,
  planning, drafting, and exhibiting artistic creations in spaces across
  campus and throughout the local community.
\item
  \textbf{ART- 325 Contemporary Photographic Genres} Creates,
  researches, and analyzes contemporary digital photography. Readings
  and writings support the weekly production of imagery. Technical
  skills include camera functions, Adobe Photoshop, and Adobe Lightroom.
  Cameras and software provided.
\item
  \textbf{ART-328 Art and Industry} Explores the relationship between
  art, mass production and the steady, habitual effort of making. Topics
  of inquiry include the Arts and Crafts Movement, the birth of industry
  and technology, and how 20th and 21st-century artists have responded
  to the frenzy of mass production and consumerism. Artworks are created
  using varying methods of mass production such as tile making, mold
  making and surface decoration techniques as a way to explore
  repetitive modes of making.
\item
  \textbf{ART-330 Methods of Teaching Art K--12} Includes discussion,
  lecture, and studio work on campus, as well as experience in public
  school classes off campus. This course does not satisfy the College's
  general education fine arts core group requirement nor does it satisfy
  major or minor requirements or prerequisites for Art of Art History.
  Prerequisite: art major or minor and admission to the Education
  Program.
\item
  \textbf{ART-331 Open Studio} Provides an opportunity for independent
  work in the studio with a midterm and final critique. P/NP basis only.
  May be taken more than once. This course does not satisfy any of the
  requirements for a major or minor in art. Prerequisites: Ceramics:
  Form, Function, \& Meaning (ART-135) or consent of instructor.
\item
  \textbf{ART-352 Mark Making} Offers instruction on materials and
  techniques that involve mark making. These may include painting and
  drawing on various surfaces, or digital painting. Through guided
  projects and independent work, students explore the connections
  between form and expression, with the aim of developing a visual
  language that is uniquely their own. Prerequisite: any 100-level ART-
  course.
\item
  \textbf{ART 360 Advanced 3D} Focuses on the conceptual, aesthetic, and
  technical skills necessary to create more advanced and sophisticated
  artworks. Designed to build upon the foundations provided in
  beginning-level Ceramics and Sculpture classes. Prerequisite: any
  100-level ART course.
\item
  \textbf{ART-361 Documentary} Explores theory and practice of
  documentary filmmaking through readings, research projects,
  screenings, and in-class discussions. Students produce studio projects
  ranging from short-form documentaries, podcasts, still photography
  portfolios, and more. Prerequisite: Digital Studio (ART-145),
  Time-Based Media (ART-150), or with permission by instructor.
\item
  \textbf{ART-363 Graphic Design Studio} Focuses on graphic design and
  the communication arts industry. Projects use traditional and digital
  tools, materials and procedures with a focus on finding creative
  visual solutions to communication problems. Prerequisite: Digital
  Studio (ART-145), or Workshop: Digital Toolbox (WKS-204), or Workshop:
  Vector Graphics (WKS-213), or consent of instructor.
\item
  \textbf{ART-364 The Human Form} Focuses on the observation and
  interpretation of the human form. Working from life, students develop
  skills in capturing the gesture and form of the body and use those
  skills to create interpretive and expressive artworks, principally
  using drawing materials, but also exploring with paint and other
  materials. Prerequisite any 100-level ART- course.
\item
  \textbf{ART-370 Video Art and Production} Explores the theoretical and
  technical foundations of video as a visual art medium. Students
  produce a series of short video projects that are presented during
  formal critiques, enabling students to cultivate meaningful dialogue
  about their work. Through production, students develop the technical
  and professional experience needed to enter the media industry.
  Prerequisite: Digital Studio (ART-145), Time-Based Media (ART-150), or
  with permission by instructor.
\item
  \textbf{ART-371 Typography and Design} Investigates the history,
  theory and practice of letterforms and typography in graphics,
  advertising, design and visual communication. Projects address
  principles of typography, letter structure, typeface selection,
  fundamentals of computer typesetting, and typographic layout.
  Prerequisite: Digital Studio (ART-145), or Workshop: Digital Toolbox
  (WKS-204), or Workshop: Vector Graphics (WKS-213), or consent of
  instructor.
\item
  \textbf{ART-373 Screen Printing} Utilizes silkscreen printing as a
  medium that can be integrated with photography, digital imagery, and
  three-dimensional objects. Course offers an exploration of the formal
  elements of design with an emphasis on the use of color. Coursework
  includes studio production, lectures, demonstrations and critiques.
  Prerequisite any 100-level ART- course.
\item
  \textbf{ART-374 Multiples in Printmaking} Utilizes intaglio and
  lithography as a medium to explore artistic vision, personal imagery
  and design. Using the process of intaglio, students investigate the
  states of development of an image by printing multiple variations of
  the plate. Using the process of lithography, students learn how to
  create a limited edition of prints. Prerequisite any 100-level studio
  ART- course.
\item
  \textbf{ART-391 Advanced Topics in Studio Art} See also Art History
  (ARH-391), p.~92 Focuses on a specific advanced studio art theme or
  topic. Topics vary. May be taken more than once for credit, provided
  the topics are substantially different. Prerequisite: any 100- or
  200-level ART- course or consent of instructor.
\item
  \textbf{ART-394 Directed Studies in Art} Investigates a topic in
  studio art selected by the student and instructor to fit the student's
  particular interests and educational needs. May be taken more than
  once for credit. Prerequisites: junior standing and consent of
  department chair.
\end{itemize}

\textbf{400-level courses}

\begin{itemize}
\tightlist
\item
  \textbf{ART-444 Independent Study} Focuses on a topic for independent
  work on a selected project under the direction of a faculty member of
  the department. Prerequisite: consent of instructor.
\item
  \textbf{ART-464 Senior Seminar I} Prepares students for advanced
  research in studio art. Emphasis is on preparation of work toward the
  senior exhibit. Only art majors are admitted to this course. Materials
  fee (where applicable) should be discussed with instructor.
  Prerequisite: declared major in art and successful completion of
  Junior Review.
\item
  \textbf{ART-474 Senior Seminar II \& Senior Exhibition} Prepares
  students for advanced research in studio art. Emphasis is on
  preparation of work toward the senior exhibit. Only art majors are
  admitted to this course. Materials fee (where applicable) should be
  discussed with instructor. Prerequisite: successful completion of
  Senior Seminar I (ART-464).
\item
  \textbf{ART-494 Internship in Art} Investigates an area of interest
  related to the major, through voluntary field placement supervised by
  a faculty member of the art and art history department. A minimum of
  140 hours on-site experience is required. P/NP basis only.
  Prerequisites: declared major in art, junior standing, and consent of
  department chair.
\end{itemize}

\section{Art History}\label{art-history}

\subsection{Art History Major}\label{art-history-major}

A grade of ``C'' (2.0) or higher must be earned in all courses counted
toward a major in art history. An Art History major may also complete an
Art minor, but only two courses may count toward both the major and the
minor.

\begin{enumerate}
\def\labelenumi{\arabic{enumi}.}
\tightlist
\item
  ARH 128 Introduction to Art History\\
\item
  Two 200-level ARH courses\\
\item
  Two additional courses approved by the department, chosen from the
  following:

  \begin{itemize}
  \tightlist
  \item
    Courses in Art History (ARH-\_\_\_)\\
  \item
    ANT 125 Art \& Archaeology Classical World\\
  \item
    HIS 318 Topics in History :Div West Persp\\
  \item
    COM 161 Visual Rhetoric\\
  \item
    FLM 225 Film History\\
  \item
    COM 357 Sex, Race, \& Gender in Media\\
  \end{itemize}
\item
  One of the following

  \begin{itemize}
  \tightlist
  \item
    ARH 307 Modern and Contemporary Art\\
  \item
    ARH 310 Contemporary Art, Theory \&Criticism\\
  \end{itemize}
\item
  ARH 464 Senior Seminar I\\
\item
  ARH 474 Senior Seminar II\\
\item
  One 100- ART course\\
\item
  One 300-level ART course
\end{enumerate}

note: Scheduling may be planned to allow at least one term of study
off-campus at one of the approved foreign or domestic programs. With
departmental approval, up to three courses from these programs may count
toward either the two 200-level courses in item 2 or the two additional
courses in item 3.

Students interested in art history or visual culture at the graduate
level should complete Intermediate French I (FRE-215).

\subsection{Art History Minor}\label{art-history-minor}

A grade of ``C'' (2.0) or higher must be earned in all courses counted
toward a minor in art history.

\begin{enumerate}
\def\labelenumi{\arabic{enumi}.}
\tightlist
\item
  ARH 128 Introduction to Art History\\
\item
  Three additional art history courses, one of which must be ARH-200 or
  above. Can also include one of the following:

  \begin{itemize}
  \tightlist
  \item
    ANT 125 Art \& Archaeology Classical World\\
  \item
    HIS 318 Topics in History :Div West Persp\\
  \item
    COM 161 Visual Rhetoric\\
  \item
    FLM 225 Film History\\
  \item
    COM 357 Sex, Race, \& Gender in Media\\
  \end{itemize}
\item
  One of the following

  \begin{itemize}
  \tightlist
  \item
    ARH 307 Modern and Contemporary Art
  \item
    ARH 310 Contemporary Art, Theory \&Criticism\\
  \end{itemize}
\item
  One 100- or 300-level ART- course
\end{enumerate}

\subsection{COURSES IN ART HISTORY}\label{courses-in-art-history}

\begin{itemize}
\tightlist
\item
  ARH 106 World Art\\
  Traces key themes in art from a global perpective, focusing on the
  ways that cultures and civilizations across time have visually
  expressed social, religious and political values. Cross-cultural
  themes may include: religion and spirituality, word and image,
  violence and death, power and propaganda, gender and society, parks
  and memorials, and ritual and body decoration.\\
\item
  ARH 107 Gender and Art\\
  Explores of the ways in which visual culture reflects and projects
  cultural biases and issues related to gender from prehistory to the
  modern era. Analyzes how gender identities can be shaped by politics,
  religion, and culture, as well as the effect of an artist's sex and/or
  sexual preferences on subject choices, media, and market values.\\
\item
  ARH 128 Introduction to Art History\\
  Examines Western art and architecture from prehistory to the later
  19th century, with emphasis on the ways in which visual culture both
  reflects and shapes societies and civilizations. Explores how works of
  art create and sustain meaning for their original audiences, and how
  some objects or visual solutions transcend their historical moment and
  surface throughout time as familiar cultural icons or references.\\
\item
  ARH 191 Topics in Art History\\
  See also +++MISSING INFO: c.art91.short +++\\
  Focuses on a selected topic or theme in art history or visual culture.
  Topics vary. May be taken more than once for credit, provided the
  topics are substantially different.\\
\item
  ARH 201 Art of the Middle Ages\\
  Explores the art and architecture of the medieval world both
  chronologically and thematically. The course examines issues such as
  patronage of the arts, pilgrimage, the cult of saints, the arts as a
  medium of cultural exchange, and the role of the artist in the Middle
  Ages. Prerequisite: Introduction to Art History (ARH-128) or consent
  of instructor.\\
\item
  ARH 218 The World of Renaissance Art\\
  Explores the visual culture of Europe from the 14th through the 16th
  centuries, focusing on topics such as competition, display, devotion,
  portraits, the printing revolution, death, and gender issues.
  Prerequisite: Introduction to Art History (ART-128) or consent of
  instructor.\\
\item
  ARH 231 Romanticism, Realism, Impressionism\\
  Focuses on the sweeping transformations in the creation, production
  and consumption of visual culture in the 19th century. The rapidly
  changing aesthetics of the dawning modern era generated passionate
  debates about the creation and reception of art during this period.
  This course uses these debates as a series of touchstones for
  understanding the visual and social landscape of the times.
  Prerequisite: Introduction to Art History (ARH-128) or consent of
  instructor.\\
\item
  ARH 248 Baroque, Rococo, and Neoclassicism\\
  Focuses on a thematic exploration of the major art movements in the 17
  th and 18th centuries, emphasizing the pendulum swings of artistic
  practice, exploration, and institutional hierarchies. Inclusion of
  cultural phenomena---the Grand Tour, the Enlightenment, revolutions,
  and the establishment of academies---factor heavily in this course.
  Prerequisite: Introduction to Art History (ARH 128) or consent of
  instructor.\\
\item
  ARH 268 History of Architecture\\
  Investigates major monuments of architectural history from prehistory
  to the present day, with an emphasis on formal and conceptual
  concepts. Key figures, theories, innovations, and functions (both
  original and altered throughout time) are also discussed.
  Prerequisite: Introduction to Art History (ARH 128) or consent of
  instructor.\\
\item
  ARH 296 Topics in Art History:Global Persp\\
  Focuses on a selected topic or theme in art history. Topics vary. May
  be taken more than once for credit, provided the topics are
  substantially different. different. Prerequisite: Introduction to Art
  History (ARH-128) or consent of instructor.\\
\item
  ARH 297 Topics in Art History: US Pluralism\\
  Examines a selected topic or theme in art history. Topics vary. May be
  taken more than once for credit, provided the topics are substantially
  different. Prerequisite: Introduction to Art History (ARH-128) or
  consent of instructor.\\
\item
  ARH 298 Topics in Art History:Div West Pers\\
  Studies a selected topic or theme in art history. Topics vary, and may
  include: Art and Cultural Property; Memory, Environment and Landscape.
  May be taken more than once for credit, provided the topics are
  substantially different. Prerequisite: Introduction to Art History
  (ARH-128) or consent of instructor.\\
\item
  ARH 307 Modern and Contemporary Art\\
  Traces the development of major artistic movements in the 20th
  centrury to the more contemporary trends of the 21st century. Topics
  include: Post-Impressionism, Cubism, Futurism, Surrealism, German
  Expressionism, Dada, Pop Art, Minimalism, Conceptual Art,
  Postmodernism, installation, new media, performance, and digital
  production and distribution. Prerequisite: Introduction to Art History
  (ARH-128) and a 200-level ARH- course or consent of instructor.\\
\item
  ARH 310 Contemporary Art, Theory \&Criticism\\
  Investigates issues in contemporary art. Focuses on art of the late
  20th and early 21st centuries, considering stylistic, historical and
  theoretical developments. Prerequisite: Introduction to Art History
  (ARH-128) and a 200-level ARH- course or consent of instructor.\\
\item
  ARH 391 Advanced Topics in Art History\\
  See also \\
  Focuses on an advanced study of a selected topic or theme in art
  history. Topics vary. May be taken more than once for credit, provided
  the topics are substantially different. Prerequisites: Introduction to
  Art History (ARH-128) and a 200-level ARH- course or consent of
  instructor.\\
\item
  ARH 394 Directed Learning in Art History\\
  Investigates topics in art history selected by the student and
  instructor to fit the student's particular interests and educational
  needs. May be taken more than once for credit. Prerequisites: junior
  standing and consent of department chair.\\
\item
  ARH 444 Independent Study: Art History\\
  Focuses on independent work on a selected project under the direction
  of a faculty member of the Art and Art History department.
  Prerequisite: consent of instructor.\\
\item
  ARH 464 Senior Seminar I\\
  Introduces the research methods, theories, and curatorial practices
  affiliated with the discipline of art history. Additionally, students
  begin crafting their capstone projects in art history, which will be
  completed in Seminar II. Prerequisites: declared major in Art History,
  Introduction to Art History (ARH-128), one 200-level ARH course, and
  junior standing.\\
\item
  ARH 474 Senior Seminar II\\
  Completes the capstone research project (i.e.~research papers,
  exhibitions, or hybrid projects that involve making art and
  contextualizing visual culture). Prerequisite: Successful completion
  of Seminar in Art History I.\\
\item
  ARH 494 Internship in Art History\\
  Investigates an area of interest related to the major, through
  voluntary field placement supervised by a faculty member of the Art
  and Art History department. A minimum of 140 hours on-site experience
  is required. P/NP basis only. Prerequisites: declared major in art
  history, junior standing, and consent of department chair.
\end{itemize}

\section{Asian Studies}\label{sec-asian-studies}

Nordmann (Administrative Coordinator).

The Asian Studies program is designed to acquaint students with the
great cultures of Asia, primarily those of China, India and Japan.

\subsection{Asian Studies Major}\label{asian-studies-major}

A major in Asian studies requires a minimum cumulative 2.0 GPA in all
courses counted toward the major.

The major is interdisciplinary and requires that the students take
courses in the areas of art history, history, literature, and religion,
and suggests that they also take courses in economics, music, and
political science.

1.\textbf{One} of the following:

\begin{verbatim}
-ASC 186 Modern South Asia
-HIS 136 East Asian Civilization
\end{verbatim}

2.Participation in one overseas study program in Asia:

\begin{verbatim}
-an ACM overseas study program (India Studies or Japan Study)
-a Coe foreign exchange program (Chiang Mai University of Thailand, Nagoya Gakuin University of Japan, or Sookmyung University in South Korea)
-Coe’s Asia Term (see p. 30)
\end{verbatim}

3.\textbf{Four} courses, chosen from the following, in consultation with
the Asian Studies administrative coordinator:

\begin{verbatim}
-ASC 106 May Term in Asia
-ASC 176 China and Japan
-ECO 436 Econ Development or ECO 446 International Econ
-HIS 217 American War in Vietnam
-HIS 246 History of Modern China
-HIS 256 History of Modern Japan
-HIS 316 Topics in History:Non-Western Persp (when appropriate)
-HIS 466 Seminar Modern East Asian History
-JPN 115 Elementary Japanese I
-JPN 125 Elementary Japanese II
-JPN 215 Intermediate Japanese I
-JPN 225 Intermediate Japanese II
-+++MISSING INFO: c.phl/rel206.long +++
-POL 286 Asian Politics
-REL 106 Eastern Religions
-A second term of an overseas study program
\end{verbatim}

\subsection{Asian Studies Minor}\label{asian-studies-minor}

1.\textbf{One} of the following:

\begin{verbatim}
-HIS 316 Topics in History:Non-Western Persp
-ASC 186 Modern South Asia
\end{verbatim}

2.\textbf{Five} courses, chosen from the following, in consultation with
the Asian Studies administrative coordinator:

\begin{verbatim}
-ASC 176 China and Japan
-ECO 436 Econ Development or ECO 446 International Econ
-HIS 217 American War in Vietnam
-HIS 246 History of Modern China
-HIS 256 History of Modern Japan
-HIS 316 Topics in History:Non-Western Persp (when appropriate)
-HIS 466 Seminar Modern East Asian History
-JPN 115 Elementary Japanese I
-JPN 125 Elementary Japanese II
-JPN 215 Intermediate Japanese I
-JPN 225 Intermediate Japanese II
-+++MISSING INFO: c.phl/rel206.long +++
-POL 286 Asian Politics
-REL 106 Eastern Religions
\end{verbatim}

\subsection{Courses in Asian Studies}\label{courses-in-asian-studies}

-\textbf{ASC 106 May Term in Asia} Examines experientially the
fast-changing nation of Vietnam through three weeks of study and travel
in the country. Students will investigate and evaluate the contemporary
nation, as well as the history and memory of its wartime experiences,
through readings, written work, class discussions, and visits to
important historical and cultural landmarks. Through these varied
experiences, students will deepen their understandings of the cultural
skills involved in overseas travel, the cultural differences and
commonalities that shape human experiences, and the specific knowledge
of Vietnam's culture and history. -\textbf{ASC 176 China and Japan} An
interdisciplinary examination of themes and ideas central to
understanding the cultures of China and Japan. Materials used include
primary texts in Taoism, Confucianism, and Buddhism, as well as works of
poetry, fiction, and visual art. -\textbf{ASC 186 Modern South Asia}
India, Pakistan, Bangladesh, and Sri Lanka from a modern perspective.
The preoccupation with identity, the reemergence of caste and religious
conflicts, the strengths and weaknesses of political and economic
development, and the implication of trends for the future. Includes
selected topics on culture and history that provide insights to
understanding the present. -\textbf{+++MISSING INFO: c.asc216.long +++}
See \textbf{?@sec-history}

\section{Athletic Training}\label{athletic-training}

Chandler, Libby (Chair).

No students will be admitted to the Athletic Training Education Program
after December 8, 2017

\subsection{Courses in Athletic
Training}\label{courses-in-athletic-training}

-\textbf{AT 100 Basic Athletic Training} An overview of the athletic
training profession designed to lead student's knowledge of the
profession from the general foundations to specific educational
competencies. Included in this course are laboratory sessions used to
instruct and practice athletic training competencies and proficiencies.
Credit is given for Prevention and Care of Athletic Injuries \&
Laboratory (PE 175 OR Basic Athletic Training (AT-100), not both.

Prerequisite: Application to the Athletic Training Education Program
-\textbf{+++MISSING INFO: c.at20\_/30\_/40\_.long +++} +++MISSING INFO:
c.at20\_/30\_/40\_.desc +++ -\textbf{+++MISSING INFO: c.at225/220.long
+++} +++MISSING INFO: c.at225/220.desc +++ -\textbf{+++MISSING INFO:
c.at235/230.long +++} +++MISSING INFO: c.at235/230.desc +++ -\textbf{AT
375 General Medical Conditions} Study of the etiology, clinical
manifestations and management of common diseases/conditions that affect
athletes and others involved in physical activity. Additional topics
include pharmacology, counseling and psychosocial interventions.
Prerequisites: Human Anatomy (BIO-215), Human Physiology (BIO-225), and
admission to Athletic Training Education Program. -\textbf{AT 385
Therapeutic Modalities} Investigation into the concepts,
indications/contraindications, physiological effects, and clinical
implications of physical modalities used in a therapeutic setting.
Prerequisite: Human Physiology (BIO-225). -\textbf{+++MISSING INFO:
c.at395/390.long +++} +++MISSING INFO: c.at395/390.desc +++ -\textbf{AT
425 Administration of Athletic Training} Introduction to the
organization and administration of an athletic training program
including facility planning/design, information/resource management,
legal/ethical considerations, risk management, insurance systems, and
professional development. Prerequisite: senior standing in the Athletic
Training Education Program. -\textbf{AT 494 Internship in Athletic
Training} Comprehensive experience in athletic training through
voluntary field placement supervised by a certified athletic trainer
(ATC®). A minimum of 140 hours on-site experience is required. P/NP
basis only. Prerequisite: senior standing in the Athletic Training
Education Program.

\section{Biochemistry (Collateral
Major)}\label{biochemistry-collateral-major}

Stead (Administrative Coordinator).

\subsection{Collateral Major in
Biochemistry}\label{collateral-major-in-biochemistry}

A major in biochemistry requires a minimum cumulative 2.0 GPA in all
courses counted toward the major. Concurrent completion of a primary
major in chemistry, including Calculus II and two General Physics
courses is required. Students choosing a collateral major in
biochemistry may not select a collateral major in molecular biology

1.+++MISSING INFO: c.bio145/145L.long +++ 2.+++MISSING INFO:
c.chm431/431L.long +++ 3.+++MISSING INFO: c.chm432/432L.long +++
4.\textbf{Two} of the following:

\begin{verbatim}
-+++MISSING INFO: c.bio325/325L.long +++
-+++MISSING INFO: c.bio345/345L.long +++
-BIO 405 Current Topics in Molecular Biology
-+++MISSING INFO: c.bio435/435L.long +++
\end{verbatim}

Genetics (BIO-235) and Laboratory (BIO-235L) and all of the courses
under \#4 are strongly encouraged, especially for those planning
graduate study.

\section{Biology}\label{biology}

R. Christensen, Cozine (Chair), Ellis, Hughes, Keyes, Leonardo,
Leveille, Luttinger, Mommer, Sanchini, Storer (Co-Chair).

In addition to a major in biology, the college also offers collateral
majors in ENVIRONMENTAL SCIENCE (see p.~115), MOLECULAR BIOLOGY (see
p.~144), and NEUROSCIENCE (see p.~156). The biology program provides
students with an early introduction to experimental design and use of
instrumentation to investigate properties of cells and organisms.
Opportunities to work one-on-one with faculty, in honors and independent
study work or in research participation, is also available to students.

\subsection{Biology Major}\label{biology-major}

A major in biology requires a minimum cumulative 2.0 GPA over courses
counted toward the major.

1.+++MISSING INFO: c.bio145/145L.long +++ 2.+++MISSING INFO:
c.bio155/155L.long +++ 3.BIO 202 Topics in Evolution 4.+++MISSING INFO:
c.bio205/205L.long +++ 5.+++MISSING INFO: c.bio235/235L.long +++
6.\textbf{One} of the following:

\begin{verbatim}
-+++MISSING INFO: c.bio355/355L.long +++
-+++MISSING INFO: c.bio375/375L.long +++
\end{verbatim}

7.\textbf{Three} biology electives with the associated laboratories
8.+++MISSING INFO: c.chm121/121L.long +++ 9.+++MISSING INFO:
c.chm122/122L.long +++ 10.\textbf{One} of the following:

\begin{verbatim}
-+++MISSING INFO: c.mth132.long +++ (strongly recommended for students considering attending graduate school)
-STA 100 Statistical Reasoning I-Foundations and STA 110 Stats IIA: Inferential Reasoning
-STA 100 Statistical Reasoning I-Foundations and STA 130 Stats IIB: Experimental Design
\end{verbatim}

11.Attendance at research seminars in junior and senior years

\emph{Strongly recommended:} \textbf{One} of the following sequences:
-+++MISSING INFO: c.phy165/165L.long +++ and +++MISSING INFO:
c.phy175/175L.long +++ -+++MISSING INFO: c.phy185/185L.long +++ and
+++MISSING INFO: c.phy195/195L.long +++

With consent of the department chair, other chemistry or mathematics
courses may be substituted.

\subsection{Laboratory Study in
Biology}\label{laboratory-study-in-biology}

Laboratory and field experiences are essential to scientific
understanding. They offer an introduction to the processes of science,
emphasizing student participation in hypothesis formulation,
experimental design, data collection, data analysis, and the composition
of papers in a format conventional to science. Biology majors must
complete the requisite laboratory courses for the major. Students not
majoring in biology should enroll in a laboratory section of any course
they plan to apply to the laboratory science requirement in general
education.

\subsection{Courses in Biology}\label{courses-in-biology}

-\textbf{BIO 100 Human Biology} A study of Homo sapiens applying the
same tools biologists use to study other organisms. Topics may include
an overview of human evolution, reproduction, genetics, population
genetics, adaptations of selected organs and organ systems, humans'
relationship to and impact on their environment, and other contemporary
issues in human biology. Includes lecture and laboratory. This course
does not satisfy any of the requirements for a major in biology.
-\textbf{+++MISSING INFO: c.bio100L.long +++} +++MISSING INFO:
c.bio100L.desc +++ -\textbf{BIO 103 Anatomy \& Physiology} Introduces
the student to the structure, function, and organization of the human
body. Designed for non-majors as well as students with health-related
career interests. This course does not satisfy any of the requirements
for a major in biology or nursing. -\textbf{+++MISSING INFO:
c.bio103L.long +++} +++MISSING INFO: c.bio103L.desc +++ -\textbf{BIO105
Introductory Biology} None -\textbf{BIO 111 Urban Ecology} Introduces
the ecology of urban systems, investigating how humans and their built
environments interact with and form ecosystems. Two main points of view
will be represented: the city as an ecological system, and the responses
and adaptations of organisms to this ecosystem. We will also investigate
how humans respond to and interact with this biotic and abiotic
environment. Labs will visit relevant areas of Cedar Rapids. One three
hour workshop and one four-hour lab per week. (Offered alternate years.)
-\textbf{+++MISSING INFO: c.bio111L.long +++} +++MISSING INFO:
c.bio111L.desc +++ -\textbf{BIO 115 Marine Biology} Surveys marine
organisms and communities with emphasis on tropical coral reef systems.
The primary field analyses focus on the dynamics of coral reefs, tidal
regions, and turtle grass meadows. An off campus course taught at an
appropriate site and involving extensive field work. Includes lecture
and laboratory. Prerequisite: consent of instructor. (Offered May Term
only) -\textbf{BIO 120 Biology, Health, \& Society} Introduces the basic
principles, terms, and technologies needed to understand the role of
biology in our lives. This course centers around answering questions in
healthcare and in contemporary society, using the application of
biological knowledge. Study skills for mastering college biology are
also taught. Three lectures and one co-curricular skills workshop per
week. Counts as a non-lab science course. This course does not satisfy
any of the requirements for a major in biology. -\textbf{BIO 145
Cellular \& Molecular Biology} A study of biology at the cellular and
molecular level, including cytology, metabolic, and genetic processes.
Three lectures per week. -\textbf{+++MISSING INFO: c.bio145L.long +++}
+++MISSING INFO: c.bio145L.desc +++ -\textbf{BIO 155 Organismal \&
Ecological Biology} The study of organ systems, reproduction, and
embryology, with an emphasis on human biology. Principles of evolution,
ecology, and animal behavior are also studied. Three lectures per week.
-\textbf{+++MISSING INFO: c.bio155L.long +++} +++MISSING INFO:
c.bio155L.desc +++ -\textbf{BIO 165 Ecology and Biology of Birds} An
introduction to the behavior, ecology and biology of birds, emphasizing
identification of local species and such species' adaptations to their
biotic and abiotic environments. Binoculars furnished. This course
satisfies the lab science course requirement. Offered Spring Term on
campus and Summer at the Wilderness Field Station. -\textbf{+++MISSING
INFO: c.bio165L.long +++} +++MISSING INFO: c.bio165L.desc +++
-\textbf{BIO 175 Field Botany} Investigates tree and wildflower species
encountered in Iowa's forests, fields and prairies. Emphasizes hands-on
description of structural characteristics, identification of species
using a dichotomous key and characteristics of flowering plant families.
A basic understanding of the distribution of plants within communities
is also developed. A collection of local plants is required. (Offered in
May Term) -\textbf{BIO 185 Entomology} Introduces the biology of insects
and related arthropods, focusing on how homeostatic problems common to
all organisms have been solved. Topics include physiology, morphology,
systematics, and the relationship among insects, plants, and humans.
-\textbf{+++MISSING INFO: c.bio185L.long +++} +++MISSING INFO:
c.bio185L.desc +++ -\textbf{BIO 195 Introduction to Microbiology} An
elementary-level study of the principles of microbiology. Special
attention is given to the medical aspect of the subject. Three lectures
and one three hour laboratory per week. Designed for students with
allied health career interests. This course does not satisfy any of the
requirements for a major in biology. Prerequisite: Human Biology
(BIO-100) or Cellular and Molecular Biology (BIO-145).
-\textbf{+++MISSING INFO: c.bio195L.long +++} +++MISSING INFO:
c.bio195L.desc +++ -\textbf{BIO 202 Topics in Evolution} Investigates
the roles of selective and non-selective evolutionary mechanisms,
especially as they affect adaptation, speciation, and extinction.
Selected aspects of the fossil record and current theories of hominid
evolution are also part of the course. Counts as a non-lab science
course. Prerequisites: Cellular and Molecular Biology (BIO-145). Note:
Credit is given for either Topics in Evolution (BIO-202) or Organic
Evolution (BIO-125), not both. -\textbf{BIO 205 General Botany} A study
of the basic structure and functions of plants, with emphasis on seed
plants. Functional aspects of structure and development on the cellular,
tissue, and whole plant level receive equal emphasis. Three lectures per
week. Prerequisites: Cellular and Molecular Biology (BIO-145),
Introduction to Biology Laboratory (BIO-140) and Organismal and
Ecological Biology (BIO-155). -\textbf{+++MISSING INFO: c.bio205L.long
+++} +++MISSING INFO: c.bio205L.desc +++ -\textbf{BIO 215 Human Anatomy}
A survey of gross and microscopic human anatomy, providing necessary
background for students in physiology and body mechanics. Three lectures
per week. Prerequisites: Human Biology (BIO-100) or Cellular and
Molecular Biology (BIO-145) or Organismal and Ecological Biology
(BIO-155). This course does not satisfy any of the requirements for a
major in biology. Credit is given for Human Anatomy (BIO-215) or
Integrated Human Anatomy (BIO-245), not both. -\textbf{+++MISSING INFO:
c.bio215L.long +++} +++MISSING INFO: c.bio215L.desc +++ -\textbf{BIO 225
Human Physiology} The functional study of the various systems of the
human body, including the application of physiological principles to the
health professions. Three lectures and one two-hour recitation per week.
This course does not satisfy any of the requirements for a major in
biology. Prerequisites: Human Anatomy (BIO-215) and either Human Biology
(BIO-100), Cellular and Molecular Biology (BIO-145), or Organismal and
Ecological Biology (BIO-155). -\textbf{BIO 235 Genetics} An exploration
of the three main branches of heredity: transmission (classical),
molecular, and population genetics. Transmission genetics examines how
genes and genetic traits are passed from generation to generation.
Molecular genetics probes the structure, function, and regulation of
genes, while population genetics investigates through mathematical
models the distribution and behavior of genes in populations. Three
lectures per week. Prerequisites: Cellular and Molecular Biology
(BIO-145), Introduction to Biology Laboratory (BIO-140), and Organismal
and Ecological Biology (BIO-155). -\textbf{+++MISSING INFO:
c.bio235L.long +++} +++MISSING INFO: c.bio235L.desc +++ -\textbf{BIO 275
Aquatic Ecology} Study of the biota of lakes and streams in northern
Minnesota and adjacent Ontario, and their abiotic environment. Extensive
field sampling enables students to study predation, diel vertical
migration, stream drift, induced morphological defenses, and food-web
interactions. Regional and North American lake and stream biodiversity
and the effects of introduced species and human intervention on aquatic
ecosystems are reviewed. Prerequisites: Cellular and Molecular Biology
(BIO-145) and Introduction to Biology Laboratory (BIO-140), or
Organismal and Ecological Biology (BIO-155) and Laboratory (BIO-150).
(Offered summers at the Wilderness Field Station) -\textbf{+++MISSING
INFO: c.bio275L.long +++} +++MISSING INFO: c.bio275L.desc +++
-\textbf{BIO 285 Animal Behavior} Principles and methods in ethology are
considered with emphasis on temporal and spatial patterns, adaptive
significances, communication, aggression and territoriality, social
organization, and the nature of regulatory mechanisms. Prerequisites:
One of the following: Cellular and Molecular Biology (BIO-145) and
Introduction to Biology Laboratory (BIO-140); Organismal and Ecological
Biology (BIO-155) and Laboratory (BIO-150); or Statistical Methods and
Data Analysis (PSY-301). (Offered summers at the Wilderness Field
Station) -\textbf{+++MISSING INFO: c.bio285L.long +++} +++MISSING INFO:
c.bio285L.desc +++ -\textbf{BIO 295 Spatial Ecology} Study of spatial
patterns of organisms in natural, agricultural, and urban landscapes.
Emphasis is on how the arrangements of organisms both reflect and
feedback on underlying resources and geography. Students learn how to
capture spatial characteristics of individuals, populations, species,
communities, and landscapes sampled in the field, integrate them with
satellite imagery, and analyze them statistically. Use of global
positioning system hardware to provide spatial reference for field
ecology surveys and use of ArcView as a data management, display, and
analysis tool are important features of the course. Three lectures per
week. Prerequisites: Organismal and Ecological Biology (BIO-155) and
Laboratory (BIO-150). -\textbf{+++MISSING INFO: c.bio295L.long +++}
+++MISSING INFO: c.bio295L.desc +++ -\textbf{BIO 315 Integrated Human
Anatomy} Integrates macroscopic human anatomy at the organ and organ
system levels with microscopic study of histology and development from
an evolutionary and sometimes experimental perspective. Three lectures
per week. Prerequisites: Cellular and Molecular Biology (BIO-145),
Introduction to Biology Laboratory (BIO-145L), Organismal and Ecological
Biology (BIO-155) and Laboratory (BIO-155L). Credit is given for
Integrated Human Anatomy (BIO-315) or Human Anatomy (BIO-215), not both.
-\textbf{+++MISSING INFO: c.bio315L.long +++} +++MISSING INFO:
c.bio315L.desc +++ -\textbf{BIO 325 Microbiology} Study of major fields
of microbiology with emphasis on bacteria and viruses. Topics include
bacterial cell structure, genetics, and pathogenesis as well as
immunology and epidemiology. Three lectures per week. Prerequisites:
Cellular and Molecular Biology (BIO-145), Introduction to Biology
Laboratory (BIO-145L), Organismal and Ecological Biology (BIO-155),
Organic Chemistry I (CHM-221) and Organic Laboratory (CHM-322) or
Genetics (BIO-235). -\textbf{+++MISSING INFO: c.bio325L.long +++}
+++MISSING INFO: c.bio325L.desc +++ -\textbf{BIO 335 Conservation and
Ecology of Mammals} The study of mammals of Minnesota and the natural
history, form/function relations, behavior, distributions, and
interaction with their environment and other organisms. Class activities
include direct observations, reading tracks and ``sign.'' Live-trap mark
and recapture studies supplemented by readings, lectures, and
discussions. (Offered summers at the Wilderness Field Station)
-\textbf{+++MISSING INFO: c.bio335L.long +++} +++MISSING INFO:
c.bio335L.desc +++ -\textbf{BIO 345 Techniques in Molecular Biology} An
intensive introduction to biological molecular techniques. Students
learn how to manipulate recombinant DNA and express and detect protein
products in prokaryotic systems. Course methods include cloning,
restriction, mapping, ligation, labeling, purification, elution,
blotting, hybridization, agarose, and acrylamide gels, protein
expression and detection, and immunological techniques. On completing
the course, students should have a working knowledge of most methods
used in molecular biology research. Course consists of two discussion
sessions per week. Prerequisites: Cellular and Molecular Biology
(BIO-145), Introduction to Biology Laboratory (BIO-140), Genetics
(BIO-235) and Laboratory (BIO-230), General Chemistry I (CHM-1XX), and
General Chemistry II (CHM-1XX). -\textbf{+++MISSING INFO: c.bio345L.long
+++} +++MISSING INFO: c.bio345L.desc +++ -\textbf{BIO 355 Animal
Physiology} Study of how organ systems in animal species solve the
fundamental problems of staying alive. This includes how to obtain
energy, how to procure the raw materials necessary for constructing
their bodies, how to gain oxygen from the environment, how to get rid of
unneeded substances and toxic materials. The course addresses how these
ends are achieved through the cooperation of cells, tissues, organs, and
organ systems in a variety of organisms including mammals, birds,
reptiles, amphibians, and fish as well as selected invertebrates. Three
lectures per week. Prerequisites: Cellular and Molecular Biology
(BIO-145), Introduction to Biology Laboratory (BIO-140), Organismal and
Ecological Biology (BIO-155). -\textbf{+++MISSING INFO: c.bio355L.long
+++} +++MISSING INFO: c.bio355L.desc +++ -\textbf{BIO 365 Comparative
Chordate Anat} A survey of the chordates and the development and
evolution of form in vertebrates, using primarily the shark, mudpuppy,
and cat in laboratory. Two lectures, one discussion recitation per week.
Prerequisite: Cellular and Molecular Biology (BIO-145), Organismal and
Ecological Biology (BIO-155) and Topics in Evolution (BIO-202).
-\textbf{+++MISSING INFO: c.bio365L.long +++} +++MISSING INFO:
c.bio365L.desc +++ -\textbf{BIO 375 Integrated Human Physiology} The
study of function in terms of basic physical and chemical processes; an
introduction to the principles of physiligy using the human organism as
a model. Three lectures per week. Prerequisite: Cellular and Molecular
Biology (BIO-145), Organismal and Ecological Biology (BIO-155), and
Introduction to Biology Laboratory (BIO-140) -\textbf{+++MISSING INFO:
c.bio375L.long +++} +++MISSING INFO: c.bio375L.desc +++ -\textbf{BIO 385
Behavior/Ecology of Vertebrates} Identification, evolutionary history,
behavior and ecological relationships of local mammals, birds, reptiles,
amphibians, and fish. Students learn and apply survey techniques
appropriate for particular groups to understand species/habitat
relationships, diversity within groups, and the impact of humans on
diversity. Three lectures per week. Prerequisites: Cellular and
Molecular Biology (BIO-145) and Topics in Evolution (BIO-202).
-\textbf{+++MISSING INFO: c.bio385L.long +++} +++MISSING INFO:
c.bio385L.desc +++ -\textbf{BIO 405 Current Topics in Molecular Biology}
An intensive exploration of Molecular Biology's primary literature,
focusing on recent advances in DNA, RNA and protein manipulation,
characterization, and control. Students read research articles,
critically evaluate published data, and communicate technical
information effectively. Three lectures per week. Prerequisites:
Techniques in Molecular Biology (BIO-345) and Genetics (BIO-235).
-\textbf{BIO 415 Developmental Biology} An exploration of the
developmental processes that occur as an organism grows from fertilized
egg to adult organism. The actions and interactions of cells and tissues
are examined in the living, growing, and rapidly changing environment of
the body. Classical embryology is examined through modern genetic
techniques, with special emphasis on how differential gene expression
makes us who we are. Survey of research methodologies are a prominent
topic. Three lectures per week. Prerequisites: Genetics (BIO-235) and
Laboratory (BIO-230). -\textbf{+++MISSING INFO: c.bio415L.long +++}
+++MISSING INFO: c.bio415L.desc +++ -\textbf{BIO 425 Ecology} Study of
the interrelationship of living organisms with their environment. The
structure and dynamics of populations, communities, and ecosystems are
examined through lectures and laboratory or field observations. Three
lectures per week. Prerequisites: General Botany (BIO-205) or consent of
instructor. -\textbf{+++MISSING INFO: c.bio425L.long +++} +++MISSING
INFO: c.bio425L.desc +++ -\textbf{BIO 435 Cell Physiology} The content,
organization, and function of the cell and its components are studied.
Other topics include the relationship of the cell to its environment,
the cell membrane, enzymes, conversion of matter and energy, cell
growth, and cell division. Three lectures per week. Prerequisites:
Cellular and Molecular Biology (BIO-145), Introduction to Biology
Laboratory (BIO-145L), Organismal and Ecological Biology (BIO-155),
Genetics (BIO-235), and General Chemistry I and Laboratory
(CHM-121/-121L). -\textbf{+++MISSING INFO: c.bio435L.long +++}
+++MISSING INFO: c.bio435L.desc +++ -\textbf{BIO 444 Independent Study}
Independent work on a selected project under direction of a member of
the department. Students are limited to two registrations in Independent
Study. May be taken for an X status grade with consent of instructor
prior to registration. Prerequisites: Courses providing appropriate
background for the project selected and consent of department chair.
-\textbf{BIO 445 Environmental Microbiology} An examination of the role
of microbes in the environment. Topics include microbial diversity,
nutrient cycling including biogeochemistry, bioremediation, and
biodegradation. The use of molecular methods and microsopy to study
microbes in situ are special emphases of the course. Three lectures per
week. Prerequisites: Cellular and Molecular Biology (BIO-145),
Introduction to Biology Laboratory (BIO-145L), Organismal and Ecological
Biology (BIO-155), Organic Chemistry I (CHM-221) with Organic Laboratory
(CHM-322) or Genetics (BIO-235). -\textbf{+++MISSING INFO:
c.bio445L.long +++} +++MISSING INFO: c.bio445L.desc +++ -\textbf{BIO 454
Research Participation} Participation in a research program designed and
directed by a biology department faculty member, including some or all
of the following: library research, research design, and field or
laboratory investigations. Scheduled by arrangement with sponsoring
professor with a time requirement of 10 hours per week minimum per
course credit. P/NP basis only. May be taken for an X status grade with
consent of instructor prior to registration. Research participation
credit does not satisfy any of the requirements for a major in biology.
(0.2 -- 1.0 course credit) -\textbf{BIO 455 Molecular Neurobiology}
Study of the organization and function of the nervous system, focusing
on both cellular physiology and functional neuroanatomy. Topics include
the cellular and molecular properties of brain cells, the neurological
basis of behavior, and disorders of the nervous system. Three lectures
per week. Prerequisites: Cellular and Molecular Biology (BIO-145),
Introduction to Biology Laboratory (BIO-140), Organismal and Ecological
Biology (BIO-155), General Chemistry I (CHM-155), General Chemistry II
(CHM-165) and Genetics (BIO-235) and Lab (BIO-230) -\textbf{+++MISSING
INFO: c.bio455L.long +++} +++MISSING INFO: c.bio455L.desc +++
-\textbf{+++MISSING INFO: c.bio462/463.long +++} +++MISSING INFO:
c.bio462/463.desc +++ -\textbf{BIO 465 Immunology} An integrated series
of lectures designed to familiarize students with cellular, molecular
and biochemical aspects of the development of the immune system and the
immune response, including B-and T-cell function, antibody structure and
function, complement, and major histocompatibility complex formation,
with special emphasis on the integration of these parts during the
response to infectious agents. Special topics include immunomodulation,
vaccination, autoimmunity, and neuro-immune interactions. The laboratory
component utilizes existing instrumentation in the biology department
and cell culture to explore different aspects of the immune system.
Three lectures and one three hour laboratory period per week.
Prerequisites: General Chemistry II (CHM 122/122L), Genetics (BIO-235),
and Integrated Human Physiology (BIO-375) (or concurrent enrollment in
BIO-375). -\textbf{+++MISSING INFO: c.bio465L.long +++} +++MISSING INFO:
c.bio465L.desc +++ -\textbf{BIO 494 Internship in Biology} Investigation
of an area of interest through voluntary field placement supervised by a
faculty member of the department. A minimum of 140 hours on-site
experience is required. P/NP basis only. This course does not satisfy
any of the requirements for a major in biology. Prerequisites: declared
major in biology, junior standing, and consent of department chair.

-\textbf{OCC 365 Oak Ridge Science Semester} See
\textbf{?@sec-ACM-and-international-programs}

-\textbf{HSS 494 Health Professions Externship} This capstone experience
for qualified junior and senior students, could finalize a career
decision or create an opportunity for a different career path. Using the
academic knowledge already obtained through course work, the student
will now apply it to real world situations as experiential learning with
professional mentors in our community and will gain additional knowledge
of what the daily activities and difficulties can, or may be present in
their particular internship profession. They will increase their own
communication skills, learn interpersonal skills with mentors, staff,
clients/patients, gain professional ethics skills, and gain additional
topic knowledge while working with their mentors. Some will gain ``hands
on'' training while others will gain their knowledge thru observation of
techniques, personal interactions, and through researching pertinent
subject matters. Some may gain public speaking experiences, while others
may increase their professional writing abilities. Each intern
experience will have its own unique set of learning outcomes. Junior or
Senior standing with appropriate GPA.

\section{Business Administration and
Economics}\label{business-administration-and-economics}

Baranowski, Carstens, Chen, J. Christensen, Crain, Eichhorn (Co-Chair),
Hall, Hayes, D., Knudson, Kuennen, Melcher, Nguyen (Co-Chair), Perry,
Trout, Westberg.

The Stead Department of Business Administration and Economics offers
majors in \textbf{ACCOUNTING}, \textbf{BUSINESS ADMINISTRATION},
\textbf{ECONOMICS}, \textbf{INTERNATIONAL BUSINESS}, and INTERNATIONAL
ECONOMICS, and minors in economics and international economics. Students
pursuing one or more of these are strongly encouraged to consult with
departmental faculty to develop specific plans of study.

\subsection{BUSINESS ADMINISTRATION}\label{business-administration}

The study of business administration is designed to develop in its
students a high degree of competence in the application of business
theories and sound judgment to the solution of business operating
problems. Graduates are able to contribute effectively to the activities
of their organizations by applying business skills required to plan,
analyze, and control company activities. Beyond enhancing the value of a
firm, the graduate should have a firm knowledge from which to achieve
professional growth in the workplace and/or graduate school. Students
majoring in business administration have two options: the business
administration major and the international business major (see p.~132).
The business major is designed to prepare students in the core areas of
business administration, with specific focus on management, finance, and
marketing. Students are also provided the opportunity to specialize in a
particular area through selection of elective courses. The international
business major is designed to prepare students for specific challenges
related to operating an organization in an international context.\\
Because of significant overlap in the business administration and the
international business requirements, only one of the two majors may be
elected by a student. See p.~132 for the International Business major
requirements.

\subsection{Business Administration
Major}\label{business-administration-major}

A major in business administration requires a minimum cumulative 2.0 GPA
in all courses counted toward the major.

\begin{enumerate}
\def\labelenumi{\arabic{enumi}.}
\tightlist
\item
  ACC 171 Principles of Accounting I
\item
  ECO 175 Principles of Macroeconomics
\item
  BUS 190 Statistical Analysis
\item
  ACC 215 Principles of Accounting II
\item
  ECO 215 Prin Microeconomics
\item
  BUS 315 Business Law I
\item
  One of the following Management courses: BUS 250 Principles of
  Management BUS 300 Human Resource Management BUS 387 Adv Top:Human Res
  Mgt BUS 470 Strategic Management
\item
  One of the following Marketing courses: BUS 330 Principles of
  Marketing BUS 460 Advertising BUS 465 Advanced Topics in Marketing BUS
  467 Consumer Behavior
\item
  One of the following Finance courses ECO 225 Money \& Banking BUS 350
  Principles of Finance BUS 352 Corporate Finance BUS 495 Investment
  Analysis ECO 495 International Finance
\item
  Two additional business courses numbered 200 or above
\end{enumerate}

Students planning to attend graduate school in business are strongly
encouraged to take calculus, additional statistics courses, intermediate
level macro/microeconomics, quantitative management courses, and
strategic management. Students should also discuss their graduate school
plans with a member of the department to determine other courses that
may be expected for admission to specific graduate programs

\subsection{Concentrations in
Business}\label{concentrations-in-business}

Students completing a major in Business Administration may elect to also
complete a concentration of courses within a specific discipline in the
major. Courses taken to fulfill requirements in the business major may
be used to also fulfill requirements in an area of concentration.
Students may declare only one concentration.

Concentrations are offered in management, marketing, and sports
management. To complete a concentration, a student must meet the
following requirements: 1. Completion of a major in Business
Administration 2. Completion of five additional courses in which a
cumulative 2.0 GPA must be earned. At least three of the five courses
must be completed at Coe.

\subsection{Management Concentration in
Business}\label{management-concentration-in-business}

\begin{enumerate}
\def\labelenumi{\arabic{enumi}.}
\tightlist
\item
  BUS 250 Principles of Management
\item
  Four of the following:\\
  BUS 300 Human Resource Management\\
  BUS 375 Business Ethics\\
  BUS 387 Adv Top:Human Res Mgt\\
  BUS 395 Organizational Behavior\\
  BUS 446 International Business Management\\
  BUS 464 Seminar in Management\\
  BUS 454 Research in Business\\
  BUS 470 Strategic Management
\end{enumerate}

\subsection{Marketing Concentration in
Business}\label{marketing-concentration-in-business}

\begin{enumerate}
\def\labelenumi{\arabic{enumi}.}
\tightlist
\item
  BUS 355 Principles of Marketing\\
\item
  Four of the following: BUS 280 Business Websites \textbf{and} BUS 281
  Social Media in Business\\
  BUS 420 Sports Marketing and Communications\\
  BUS 460 Advertising\\
  BUS 461 Marketing Decision Making/Analysis\\
  BUS 465 Advanced Topics in Marketing\\
  BUS 470 Strategic Management\\
  BUS 454 Research in Business\\
  PR 205 Public Relations
\end{enumerate}

\subsection{Sports Management Concentration in
Business}\label{sports-management-concentration-in-business}

\begin{enumerate}
\def\labelenumi{\arabic{enumi}.}
\tightlist
\item
  SMT 100 Introduction to Sports Management\\
\item
  Four of the following:\\
  BUS 400 Legal Issues in Sports\\
  BUS 420 Sports Marketing and Communications\\
  BUS 454 Research in Business\\
  ECO 125 Sports Analytics\\
  SMT 300 Sport/Recreation Event Management\\
  SMT 454 Seminar in Sports Management\\
  SMT 494 Internship in Sports Management
\end{enumerate}

\section{Accounting}\label{sec-accounting}

Students majoring in accounting have two options: the managerial
accounting major and the public accounting major. The managerial
accounting major is designed for students interested in pursuing a
career in private industry or the not-for-profit sector. The public
accounting major is designed for students interested in practicing
public accounting and those desiring the Certified Public Accounting
(CPA) designation. Because of the identical core course requirements in
the managerial accounting and public accounting majors, only one of the
two majors can be elected by a student. A student receiving a managerial
accounting major may pursue a post-baccalaureate major in public
accounting, but when completed, the public accounting major would
replace the managerial accounting major in the student's records.

Requirements for becoming a CPA differ from state to state. All states
require that students pass the Uniform CPA examination and many require
the completion of 150 semester hours of education prior to
certification. The 150 hours can be completed in undergraduate or
graduate programs. To become a CPA in Iowa, the Iowa State Board of
Accountancy requires that students pass the CPA examination, complete
150 semester hours of education, and complete specified course work in
accounting and business subjects. Students majoring in accounting,
business, or economics are required to have a minimum of 24 of their 150
semester hours in financial accounting, auditing, taxation, and
management accounting courses above the elementary level. If students
majoring in any other field wish to become a CPA in Iowa, they must, in
addition, complete at least 24 semester hours in business-related
courses (not including internships). (One course credit completed at Coe
is counted as four semester hours.) Accounting students at Coe should
work closely with faculty in the department to determine their plan of
study. Students wishing to become a CPA should also check the state
board's education requirements for the state in which they wish to
practice.

\subsection{Managerial Accounting
Major}\label{managerial-accounting-major}

A major in managerial accounting requires a minimum cumulative 2.0 GPA
in all courses counted toward the major.

1.ACC 171 Principles of Accounting I 2.ECO 175 Principles of
Macroeconomics 3.BUS 190 Statistical Analysis 4.ACC 215 Principles of
Accounting II 5.ECO 215 Prin Microeconomics 6.ACC 301 Intermediate
Accounting I 7.ACC 315 Managerial Accounting 8.BUS 315 Business Law I
9.ACC 334 Computer Applications - Accounting 10.ACC 351 Intermediate
Accounting II 11.ACC 413 Auditing 12.ACC 415 Adv Managerial Accounting

\subsection{Public Accounting Major}\label{public-accounting-major}

A major in public accounting requires a minimum cumulative 2.0 GPA in
all courses counted toward the major. Students wishing to prepare for
the examination to become a Certified Public Accountant should select
the Public Accounting major, which meets the guidelines for the breadth
of business and accounting courses recommended by The American Institute
of Certified Public Accountants. Students at Coe can meet the
150-semester hour requirement by completing 37.5 credits.

If degree requirements are not completed by the end of four years,
students pursuing a major in public accounting, who are in good academic
standing, and who are full-time students, receive a 60\% discount on
tuition in the consecutive fifth year. They continue to be eligible to
apply for student loans and any available federal and state government
grants

1.ACC 171 Principles of Accounting I 2.ECO 175 Principles of
Macroeconomics 3.BUS 190 Statistical Analysis 4.ACC 215 Principles of
Accounting II 5.ECO 215 Prin Microeconomics 6.ACC 301 Intermediate
Accounting I 7.ACC 315 Managerial Accounting 8.BUS 315 Business Law I
9.ACC 334 Computer Applications - Accounting 10.ACC 351 Intermediate
Accounting II 11.ACC 413 Auditing 12.ACC 422 Tax Accounting 13.ACC 463
Advanced Auditing 14.\textbf{One} of the following:

\begin{verbatim}
-BUS 410 Business Law II
-BUS 412 Entrepreneurship & New Business Frm
-BUS 464 Seminar in Management
-BUS 470 Strategic Management
\end{verbatim}

15.\textbf{Three} of the following:

\begin{verbatim}
-ACC 313 Fraud Examination
-ACC 401 Fund Acctg & Adv Top
-ACC 451 Advanced Accounting
-ACC 472 Advanced Tax Accounting
\end{verbatim}

\subsection{Courses in Accounting}\label{courses-in-accounting}

-\textbf{ACC 171 Principles of Accounting I} +++MISSING INFO:
c.aac171.desc +++ -\textbf{ACC 215 Principles of Accounting II}
Continuation of the study of financial accounting, including the
statement of cash flow. A study of managerial accounting and the
analysis of financial information for management decision making.
Prerequisite: Principles of Accounting I (ACC-171). -\textbf{ACC 301
Intermediate Accounting I} An in-depth study of financial accounting
topics. Emphasis is on income determination, evaluation, and
presentation of assets, liabilities, and owners' equity. Prerequisite:
Principles of Accounting II (ACC-215). -\textbf{ACC 313 Fraud
Examination} An examination of real world cases that represent the many
types of accounting fraud in order to gain an understanding of the
behavioral and social factors that motivate offenders. Emphasis is
placed on understanding the exposures to loss and the appropriate
prevention, detection, and investigation approaches, including legal
requirements for testifying as an expert witness. Prerequisites:
Principles of Accounting II (ACC-215) or consent of instructor. (Offered
alternate years) -\textbf{ACC 315 Managerial Accounting}

\backmatter
\printindex

\end{document}
